\documentclass[11pt]{article}

% -----------------------------------------------------------------------------
% Preamble supplied by the authors.  The paper body is intentionally contained
% in this file; in particular, there are no \input or \subfile commands.
% -----------------------------------------------------------------------------
\usepackage{fullpage}
\usepackage[T1]{fontenc}
\usepackage[utf8]{inputenc}
\usepackage{lmodern} % Scalable fonts for the compiled conference manuscript.
\usepackage{xcolor}
\definecolor{red}{HTML}{F88379}
% Unique PDF destinations also cover shared theorem counters and restatements.
\usepackage[colorlinks=true,hypertexnames=false]{hyperref}
\hypersetup{
  linkcolor=blue,
  citecolor=orange,
  filecolor=purple,
  urlcolor=orange,
  menucolor=purple,
  runcolor=purple,
  pdftitle={Consensus Clustering below Two: Online Recourse and Fully Dynamic Algorithms},
  pdfsubject={Online consensus clustering with recourse; collaborative draft with reserved abstract, introduction, and dynamic sections}
}
\usepackage{amsmath,amssymb,amsthm,amsfonts,graphicx,color,subcaption,
  enumerate,bm,array,mathtools}
\usepackage{multicol,multirow}
\usepackage{thmtools,thm-restate}
\usepackage{ragged2e}
\usepackage{lineno}
\sloppy
\usepackage{framed}
\usepackage[framemethod=tikz]{mdframed}
% Comment this out for no line numbers.
% \linenumbers

% The supplied preamble uses algorithm2e.  All algorithms below are typeset as
% framed figures so that this source also compiles in minimal TeX installations.
\IfFileExists{algorithm2e.sty}{
  \usepackage[linesnumbered,ruled,vlined]{algorithm2e}
}{}

\makeatletter
\renewcommand{\nllabel}[1]
 {{\let\@currentlabel\algocf@currentlabel
  \let\@currentcounter\algocf@currentcounter
  \label{#1}}}%

\renewcommand{\algocf@nl@sethref}[1]{%
  \renewcommand{\theHAlgoLine}{\thealgocfproc.#1}%
  \hyper@refstepcounter{AlgoLine}%
  \gdef\algocf@currentlabel{#1}%
  \gdef\algocf@currentcounter{AlgoLine}%
 }%
\makeatother

\usepackage[capitalise,nameinlink]{cleveref}

\graphicspath{{graphics/}}
\crefname{claim}{Claim}{Claims}
\crefname{property}{Property}{Properties}
\crefname{algocf}{Algorithm}{Algorithms}
\Crefname{algocf}{Algorithm}{Algorithms}

% Overleaf has the bibliography setup used in the supplied preamble.  The
% fallback keeps the same source compilable in a minimal local TeX installation.
\newif\ifusebiblatex
\IfFileExists{biblatex.sty}{
  \usebiblatextrue
  \IfFileExists{trad-alpha.bbx}{
    \usepackage[style=trad-alpha,natbib=true,maxcitenames=4]{biblatex}
  }{
    \usepackage[style=alphabetic,natbib=true,maxcitenames=4]{biblatex}
  }
  \addbibresource{ref.bib}
}{
  \usebiblatexfalse
}
\usepackage{dirtytalk}

% A dedicated algorithm float keeps the pseudocode presentation independent
% of the particular algorithm2e version installed by the compiler.
\DeclareCaptionType{paperalgorithm}[Algorithm][List of Algorithms]
\crefname{paperalgorithm}{Algorithm}{Algorithms}
\Crefname{paperalgorithm}{Algorithm}{Algorithms}

\makeatletter
\g@addto@macro\bfseries{\boldmath}
\makeatother

\newtheorem{theorem}{Theorem}[section]
\newtheorem{lemma}[theorem]{Lemma}
\newtheorem{claim}[theorem]{Claim}
\newtheorem{fact}[theorem]{Fact}
\newtheorem{corollary}[theorem]{Corollary}
\newtheorem{observation}[theorem]{Observation}
\newtheorem{proposition}[theorem]{Proposition}
\newtheorem{property}[theorem]{Property}
\theoremstyle{definition}
\newtheorem{definition}[theorem]{Definition}
\newtheorem{axiom}[theorem]{Axiom}
\theoremstyle{remark}
\newtheorem{remark}[theorem]{Remark}
\newtheorem*{remark*}{Remark}

% -----------------------------------------------------------------------------
% Notation.  These commands replace the external macros file used by the old
% project, keeping the manuscript self-contained apart from ref.bib.
% -----------------------------------------------------------------------------
\DeclareMathOperator*{\argmin}{arg\,min}
\DeclareMathOperator*{\argmax}{arg\,max}
\newcommand{\A}{\mathcal{A}}
\newcommand{\C}{\mathcal{C}}
\newcommand{\F}{\mathcal{F}}
\newcommand{\Hh}{\mathcal{H}}
\newcommand{\I}{\mathcal{I}}
\newcommand{\Pp}{\mathcal{P}}
\newcommand{\Qq}{\mathcal{Q}}
\newcommand{\Ss}{\mathcal{S}}
\newcommand{\Y}{\mathcal{Y}}
\newcommand{\OPT}{\operatorname{OPT}}
\newcommand{\cost}{\operatorname{cost}}
\newcommand{\err}{\operatorname{err}}
\newcommand{\dist}{\operatorname{dist}}
\newcommand{\Agree}{\ensuremath{\mathsf{Agree}}}
\newcommand{\Pivot}{\ensuremath{\mathsf{Pivot}}}
\newcommand{\eps}{\varepsilon}
\newcommand{\1}{\mathbf{1}}
\newcommand{\E}{\mathbb{E}}
\newcommand{\Prb}{\mathbb{P}}
\newcommand{\Vtwo}{\binom{V}{2}}
\newcommand{\polylog}{\operatorname{polylog}}
\newcommand{\wtO}{\widetilde{O}}
\newcommand{\restr}[2]{#1|_{#2}}
\newcommand{\same}{\mathsf{same}}
\newcommand{\sep}{\mathsf{sep}}
\newcommand{\card}[1]{\left\lvert #1 \right \rvert}
\newcommand{\Otilde}[1]{\widetilde{O}\left(#1\right)}
\newcommand{\kushagra}[1]{\textcolor{red}{\newline Kushagra: #1}}


% ---- command for algorithm2e ----------------------------------------------

% ---- custom colors ---------------------------------------
\definecolor{oldcol}{RGB}{110,110,110}
\definecolor{newcol}{RGB}{0,90,200}
% \Old{text}  -> comment on its own line
% \Old*{text} -> end-of-line comment after a statement
% end custom colors ---------------------------------------

\NewDocumentCommand{\Old}{s m}{%
  \IfBooleanTF{#1}
    {\tcp*{\textcolor{oldcol}{#2}}}
    {\tcp{\textcolor{oldcol}{#2}}}}

% \New{code} -> a normal code line (ends with \; automatically), colored
\NewDocumentCommand{\New}{m}{\textcolor{newcol}{#1}}


\SetCommentSty{textit}
\SetKwProg{On}{on}{:}{end}

% end command for algorithm2e ----------------------------------------------




% note commands
\newcommand{\colnote}[3]{\textcolor{#1}{$\ll$\textsf{#2}$\gg$\marginpar{\tiny\bf #3}}}
\newcommand{\anote}[1]{\colnote{blue}{#1--Amir}{AC}}
\definecolor{greenlantern}{rgb}{0.0, 0.5, 0.0}
\newcommand{\tienlong}[1]{\colnote{greenlantern}{#1--Long}{TL}}
\newcommand{\bug}[1]{\colnote{red}{#1--Bug}{Bug}}
\allowdisplaybreaks

\title{Consensus Clustering below Two:\\
Online Recourse and Fully Dynamic Algorithms}
\author{Anonymous Authors}
\date{}

\begin{document}
\maketitle

\begin{abstract}
placeholder, To be completed
\end{abstract}

\section{Introduction}\label{sec:introduction}

placeholder, To be completed

\section{Preliminaries}\label{sec:preliminaries}

\tienlong{Need restructured to add the following:}
- (Implicit) Dynamic model: only output when query, works for fully dynamic (insertion and deletion)

- Explicit dynamic model: output updated vertics in every timestep, works for insertion only 

State the main theorems:

\begin{restatable}{theorem}{OnlineMain}\label{thm:online-main}
There is a randomized online algorithm that, given a set of $n$ vertices and every oblivious vertex-arrival instance with $m$ input clusterings,
has the following guarantees. For some parameter $\zeta \in (0, 1)$, with probability at least $1-\zeta$,
its outputs satisfy
\begin{equation}\label{eq:online-ratio}
  \cost_t(C_t)
  \le
  \left(2-\frac{2}{3\eta+5}+\eps\right)\OPT_t
  \qquad\text{simultaneously for all }t\in[n].
\end{equation}
where $\eta$ is the state-of-the-art competitive ratio of the online correlation clustering problem.
For every realization of its random choices and every vertex $v$, it recourse is given by,
\begin{equation}\label{eq:online-recourse-main}
  \operatorname{rec}(v)
  =
  O_{\eta}\!\left(
    \eps^{-2}
    \log\frac{2n}{\zeta}\,
    \log\frac{2mn}{\zeta}
  \right).
\end{equation}

\end{restatable}

\begin{restatable}{theorem}{ImplicitDynamic}\label{thm:implicit-dynamic}
There is a randomized fully dynamic algorithm that, given $m$ input clusterings and a known update bound $U$, has the following guarantees for every oblivious sequence of
at most $U$ vertex insertions and deletions.
For every $\varepsilon\in(0,1]$ and $\zeta\in(0,1)$,
with probability at least $1-\zeta$, its maintained
clusterings satisfy
$\operatorname{cost}_t(C_t) \le (1.6944+\varepsilon)\operatorname{OPT}_t$
simultaneously after every update in time
$O\left(\varepsilon^{-3}\log^2\frac{U}{\varepsilon\zeta}\log^2(N)\right),$
where $N$ is the maximum number of active vertices. Membership and same-cluster queries take $O(\log(n))$ worst-case time, and listing all cluster identifiers takes $O(n)$ time.
\end{restatable}

\begin{restatable}{theorem}{ExplicitDynamic}\label{thm:explicit.dynamic}
There is an absolute constant $ \eta\ge1 $ with the following property.
For every known horizon $ n $, every oblivious vertex-arrival instance with $ m $ input partitions, every $ 0<\eps<2/(3\eta+5) $, and every $ \zeta\in(0,1) $, there is a randomized explicit dynamic consensus-clustering algorithm such that, with probability at least $ 1-\zeta $,
\begin{align}
    \cost_t(C_t)\le\left(2-\frac{2}{3\eta+5}+\eps\right)\OPT_t
    \qquad\text{simultaneously for every }t\in[n].
    \label{eq:explicit.dynamic.ratio}
\end{align}
The algorithm stores an output cluster identifier for every vertex and reports the arriving vertex's identifier and every change to an old vertex's identifier.
Under the constant-time input-label access of \cref{subsec:dynamic-model}, its expected amortized update time, including these reports, is
\begin{align}
    O\!\left(
        \eps^{-3}\log^2\frac{2n}{\zeta}
        \log\frac{2n}{\eps\zeta}\cdot\polylog(n+2)
    \right).
    \nonumber
\end{align}
In particular, for fixed $ \eps $ and inverse-polynomial $ \zeta $, the expected amortized update time is $ \Otilde{1} $.
\end{restatable}

\kushagra{Written only for the Online part.}

We first define the consensus-clustering objective, the vertex-arrival
model, and the measure of changes to previously assigned vertices.
We then state the auxiliary correlation-clustering guarantee used in the
online algorithm. All subsequent definitions and results in this section
concern the online part of the paper.

\paragraph{Conventions.}
For a positive integer $r$, write $[r]=\{1,\ldots,r\}$.
For a finite set $W$ and a nonnegative integer $k$, let $\binom{W}{k}$
denote the family of its $k$-element subsets; in particular,
$\binom{W}{2}$ is the set of unordered pairs of distinct elements of $W$.
For sets $X,Y$, their symmetric difference is
$X\triangle Y=(X\setminus Y)\cup(Y\setminus X)$.
The notation $\1[\mathsf{E}]$ is $1$ when the statement $\mathsf{E}$ holds
and $0$ otherwise. We use $\Prb$ and $\E$ for probability and expectation.
Unless a base is displayed, logarithms are natural.
A subscript on asymptotic $O(\cdot)$ notation lists the parameters on
which the hidden constant may depend.

\subsection{Partitions and the consensus objective}\label{subsec:partitions}

A \emph{partition} $C$ of a finite set $W$ is a collection of nonempty,
pairwise disjoint subsets whose union is $W$. Its parts are called
\emph{clusters}. The empty set has the unique empty partition.
For $u,v\in W$, write $u\sim_C v$ if they belong to the same cluster.
The set of pairs joined by $C$ is
\[
  E_W^+(C)=\bigl\{\{u,v\}\in\tbinom{W}{2}:u\sim_C v\bigr\}.
\]
We omit the subscript $W$ when the ground set is clear.

The distance between two partitions counts the pairs on which their
decisions to join or separate disagree.

\begin{definition}[Partition distance]\label{def:partition-distance}
For partitions $C,D$ of the same set $W$, define
\begin{equation}\label{eq:partition-distance}
  d_W(C,D)=|E_W^+(C)\triangle E_W^+(D)|.
\end{equation}
We write $d(C,D)$ when both partitions have an unambiguous ground set.
\end{definition}

For $U\subseteq W$, the \emph{restriction} $C|_U$ is the partition obtained
by intersecting every cluster of $C$ with $U$ and removing empty
intersections. The following elementary properties will be used throughout:
partition distance satisfies the triangle inequality, and restriction can
only decrease it.

\begin{observation}\label{obs:partition-metric}
For partitions $C,D,F$ of $W$ and a subset $U\subseteq W$,
\begin{align}
  d_W(C,D)&\le d_W(C,F)+d_W(F,D), \label{eq:partition-triangle}\\
  d_U(C|_U,D|_U)&\le d_W(C,D). \label{eq:restriction-distance}
\end{align}
\end{observation}
\begin{proof}
Associate a partition with its binary incidence vector on $\binom{W}{2}$:
a coordinate is $1$ precisely when its pair is joined. Partition distance
is the number of coordinates on which two such vectors differ. This gives
the triangle inequality coordinate by coordinate. Restriction deletes
coordinates outside $\binom{U}{2}$ and therefore cannot increase the
number of disagreements.
\end{proof}

Let $m$ be a positive integer, and let
$\Pp=(P_1,\ldots,P_m)$ be an indexed collection of input partitions of $W$.
Repeated partitions are allowed and retain their separate indices.
The \emph{consensus cost} of a partition $C$ and the optimum cost are
\begin{equation}\label{eq:consensus-objective}
  \cost_{\Pp}(C)=\frac1m\sum_{i=1}^m d_W(C,P_i),
  \qquad
  \OPT(\Pp)=\min_{\text{partitions }C\text{ of }W}\cost_{\Pp}(C).
\end{equation}
We use the average, rather than the sum, of the input distances. This
normalization does not affect approximation ratios.
Applying \eqref{eq:partition-triangle} to each input gives the useful
comparison inequality
\begin{equation}\label{eq:cost-comparison}
  \cost_{\Pp}(C)\le \cost_{\Pp}(F)+d_W(C,F)
\end{equation}
for any two partitions $C,F$ of $W$.

\subsection{The online model and recourse}\label{subsec:online-model}

Fix a vertex set $V=\{v_1,\ldots,v_n\}$ of size $n\ge1$ and $m$
input partitions $P_1,\ldots,P_m$ of $V$.
The vertices arrive in the order $v_1,\ldots,v_n$.
We assume that $n$ and $m$ are known to the algorithm.
At time $t\in[n]$, the active vertex set and the restricted input collection
are
\[
  V_t=\{v_1,\ldots,v_t\},\qquad
  P_i^t=P_i|_{V_t}\quad(i\in[m]),\qquad
  \Pp^t=(P_1^t,\ldots,P_m^t).
\]
For convenience, $V_0=\varnothing$.
When $v_t$ arrives, the algorithm receives its cluster label in each
input partition. More precisely, input $i$ provides a persistent label
$c^i(v)$ for each vertex $v$, with
$c^i(u)=c^i(v)$ if and only if $u\sim_{P_i}v$.
Previously revealed labels do not change.
Thus arrival reveals how the new vertex relates to every earlier vertex
in every input, while all relations between earlier vertices remain fixed.

Write
\begin{equation}\label{eq:prefix-objective}
  \cost_t(C)=\cost_{\Pp^t}(C),
  \qquad
  \OPT_t=\OPT(\Pp^t).
\end{equation}
For the analysis, fix an optimal partition $F_t^*$ of $V_t$ for each $t$,
using a deterministic tie-breaking rule independent of the algorithm's
random choices. These reference partitions are not available to the
algorithm. Restricting an optimal partition at time $t+1$ to $V_t$
and using \eqref{eq:restriction-distance} shows that
$\OPT_t\le\OPT_{t+1}$.

After processing $v_t$, an online algorithm outputs a partition $C_t$ of
$V_t$. For a number $R\ge1$, it is \emph{$R$-competitive} on an execution
if $\cost_t(C_t)\le R\OPT_t$ at every time $t$.
An \emph{oblivious instance} is one in which the full input partitions
and the arrival order are fixed independently of the algorithm's random
choices. All probability guarantees below concern this model.

We measure changes to the output through persistent cluster identifiers.
Specifically, the algorithm assigns each active vertex $v$ an identifier
$\lambda_t(v)$, and two active vertices belong to the same output cluster
if and only if their identifiers agree.
For a vertex $v_s$ arriving at time $s$, its \emph{recourse} is
\begin{equation}\label{eq:recourse-definition}
  \operatorname{rec}(v_s)
  =\sum_{t=s+1}^n
    \1[\lambda_t(v_s)\ne\lambda_{t-1}(v_s)].
\end{equation}
The initial assignment is free. Every subsequent identifier change
counts, including a change that only renames a cluster.
A per-vertex recourse bound controls \eqref{eq:recourse-definition}
separately for every vertex, rather than only its sum over all vertices.

\subsection{An auxiliary online correlation-clustering algorithm}
\label{subsec:agreement}

A \emph{complete signed graph} $H$ on $W$ assigns either a positive or
a negative sign to every pair in $\binom{W}{2}$.
Let $E^+(H)$ denote its set of positive pairs; all other pairs are negative.
The correlation-clustering error of a partition $C$ of $W$ is
\begin{equation}\label{eq:cc-objective}
  \err_H(C)=|E^+(H)\triangle E_W^+(C)|,
  \qquad
  \OPT_{\rm CC}(H)
  =\min_{\text{partitions }C\text{ of }W}\err_H(C).
\end{equation}
Thus an error is either a positive pair separated by $C$ or a negative pair
joined by $C$. The positive relation of $H$ need not itself be a partition.
Nevertheless, the same coordinatewise argument as in
\cref{obs:partition-metric} gives
\begin{equation}\label{eq:signed-comparison}
  d_W(C,F)\le \err_H(C)+\err_H(F)
\end{equation}
for any partitions $C,F$ of $W$.

For $U\subseteq W$, write $H|_U$ for the signed graph induced on $U$.
A vertex-arrival signed-graph sequence consists of graphs $H_t$ on $V_t$
such that $H_{t+1}|_{V_t}=H_t$: signs between old vertices never change.
We use the deterministic Online Agreement algorithm of
Cohen-Addad et al.~\cite{cohen2022online}, denoted by $\Agree$, on such
sequences. A \emph{persistent execution} means that one copy of this
algorithm processes the whole sequence without being restarted.
Its output at time $t$ is denoted by $D_t$.

Informally, $\Agree$ maintains a constant-factor approximation at every
prefix while changing each vertex's cluster identifier only logarithmically
many times. Let $\eta\ge1$ denote its approximation constant and
$q_{\rm A}>0$ the constant in its per-vertex recourse bound.
We use the following interface.

\begin{theorem}[Online Agreement {\cite[Theorems E.2 and F.1]{cohen2022online}}]
\label{thm:agree-blackbox}
There are absolute constants $\eta\ge1$ and $q_{\rm A}>0$ and a
deterministic algorithm $\Agree$ such that, on every vertex-arrival
sequence of complete signed graphs, a persistent execution satisfies
\[
  \err_{H_t}(D_t)\le\eta\,\OPT_{\rm CC}(H_t)
  \qquad\text{for every }t\in[n].
\]
Its output admits persistent cluster identifiers for which every vertex
has at most $q_{\rm A}\lceil\log_2 n\rceil$ changes.
\end{theorem}

Fix an implementation and its internal parameters and tie-breaking rules
once and for all. We retain the symbolic constant $\eta$ in our competitive
ratio, rather than assigning it a numerical value not specified by the
interface. The auxiliary objective \eqref{eq:cc-objective} is unweighted.
It will be used only to interpret the majority of three inputs; the
algorithm's final output is evaluated by the full consensus objective
\eqref{eq:prefix-objective}.

\subsection{The Fully Dynamic Model}
\label{subsec:dynamic-model}

There are $m$ indexed input partitions $P_1,\ldots,P_m$.
Starting from $V_0=\emptyset$, each update inserts or deletes
one vertex. At time $t$, the active set is $V_t$, and the
current input instance is
\[
    \mathcal P_t=(P_1|_{V_t},\ldots,P_m|_{V_t}).
\]
Write
\[
    \operatorname{cost}_t(C)
        =\operatorname{cost}_{\mathcal P_t}(C),
    \qquad
    \operatorname{OPT}_t
        =\min_C \operatorname{cost}_t(C).
\]

Each input partition is represented by persistent cluster
labels occupying one machine word. An insertion provides
constant-time random access to any selected input label of
the new vertex. Labels of surviving vertices remain unchanged.
The update sequence is oblivious: vertex identities, input
labels, and insertion and deletion choices are fixed
independently of the algorithm's randomness. The maintained clustering is accessed through queries.
A membership query returns the current cluster identifier
of an active vertex, and a same-cluster query tests whether
two active vertices belong to the same cluster.

\paragraph{Explicit Dynamic Model}\label{par:explicit_dynamic_model} % (fold)
We call a dynamic model \emph{explicit} if for each update, the algorithm reports explicitly which part of the solution has changed.
In particular, for consensus clustering and correlation clustering problems, we represent the the solution by a cluster identifier for each vertex. 
Then, for each update, the algorithm reports the new vertex and its cluster identifier, and for each vertex whose cluster identifier has changed, the algorithm reports the vertex and its new cluster identifier.
We refer to the algorithm in this model as an explicit dynamic algorithm.
% pararaph Explicit Dynamic Model (end)

\section{Online Consensus Clustering}\label{sec:online}
In this section, we prove our guarantee for the online consensus clustering. Our main online result gives a competitive ratio strictly below $2$ with polylogarithmic recourse for every vertex. Let us restate the main result.


\OnlineMain*

For fixed $\eta, \eps$, the bound in
\eqref{eq:online-recourse-main} is polylogarithmic in the instance size
and the inverse failure probability. We compute candidate costs exactly
with respect to all $m$ inputs. The theorem concerns approximation and
recourse; it does not assert a near-linear running time.

\paragraph{Degenerate cases.}
If $n=1$, the unique partition is optimal and incurs no recourse.
If $m=1$, following the single input is optimal.
If $m=2$, following $P_1^t$ is optimal at every time, since for every
partition $F$ of $V_t$,
\[
  \cost_t(F)
  =\frac{d(F,P_1^t)+d(F,P_2^t)}2
  \ge\frac{d(P_1^t,P_2^t)}2
  =\cost_t(P_1^t).
\]
Following an input uses its persistent labels and incurs no recourse.
We therefore assume $n\ge2$ and $m\ge3$ for the rest of the section.


\subsection{Algorithm}
\label{sec:oncon-algorithm}

In this section, we describe our online consensus clustering algorithm,
which we call \texttt{oncon}. The input consists of $m$ partitions
$P_1,\ldots,P_m$ of a vertex set $V=\{v_1,\ldots,v_n\}$, revealed through
vertex arrivals. The values of $n$ and $m$ are known in advance, but the
partitions are not. When vertex $v_t$ arrives, the algorithm receives its
cluster label $c_i(v_t)$ in each input partition $P_i$. These labels are
persistent: $c_i(u)=c_i(v)$ if and only if $u$ and $v$ belong to the same
cluster of $P_i$.

Let $V_t=\{v_1,\ldots,v_t\}$ and $P_i^t=P_i|_{V_t}$. After each arrival,
the algorithm outputs a partition $C_t$ of $V_t$, represented by cluster
identifiers $\lambda_t(v)$ for $v\in V_t$. Two vertices belong to the same
output cluster precisely when their identifiers agree. The consensus
cost of a partition $C$ of $V_t$ is
\[
    \operatorname{cost}_t(C)
    =
    \frac{1}{m}\sum_{i=1}^{m} d(C,P_i^t),
\]
where $d(C,P_i^t)$ counts the unordered vertex pairs on which the two
partitions disagree. Our goal is to maintain a low-cost partition at
every prefix while changing each vertex's identifier only a small
number of times after its arrival. The approximation and recourse
guarantees are established in the analysis.

For a positive integer $r$, write $[r]=\{1,\ldots,r\}$. If $n=1$ or
$m\leq 2$, the algorithm simply follows $P_1$, using its persistent
labels as output identifiers. Henceforth, assume that $n\geq 2$ and
$m\geq 3$. To describe the algorithm, we introduce the following
definitions.

\paragraph{Majority graph.}
Let $Q=\{i,j,k\}\subseteq[m]$ be a set of three distinct input indices.
At time $t$, the majority graph $H_Q(t)$ is the complete signed graph on
$V_t$ in which a pair $\{u,v\}$ is positive if at least two of
$P_i^t,P_j^t,P_k^t$ place $u$ and $v$ in the same cluster. Otherwise,
the pair is negative. The indices in $Q$ are distinct, although the
corresponding input partitions may coincide.

When a vertex arrives, its incident signs are determined from its
labels in the three sampled inputs. Signs between previously arrived
vertices remain unchanged. Thus,
$H_Q(1),\ldots,H_Q(n)$ forms a valid vertex-arrival sequence of signed
graphs.

\paragraph{Auxiliary clustering algorithm.}
The positive pairs of a majority graph need not define a partition.
We therefore use the deterministic online correlation-clustering
algorithm \texttt{Agree} from the preliminaries to obtain a partition
from this graph. A disagreement with a signed graph is a positive pair separated by a
partition or a negative pair joined by it. Let $\eta\geq 1$ denote the
approximation factor of \texttt{Agree}: at every prefix, its output has
at most $\eta$ times the minimum possible number of disagreements.
A persistent execution of \texttt{Agree} also changes each vertex's
cluster identifier only $O(\log n)$ times.

\paragraph{Parameters.}
The algorithm takes an accuracy parameter
\[
    0<\varepsilon<\frac{2}{3\eta+5}
\]
and a failure parameter $\zeta\in(0,1)$. Define
\[
    K=
    \left\lceil
        \frac{4}{\varepsilon}\log\frac{2n}{\zeta}
    \right\rceil,
    \qquad
    w=\frac{3(\eta+1)}{3\eta+5},
    \qquad
    \theta=\frac{\varepsilon}{4}.
\]
Here $K$ is the number of candidate processes, $w$ is the probability
of sampling an input candidate, and $\theta$ determines the switching
threshold. Logarithms are natural.

\paragraph{Candidate processes.}
A candidate process maintains a partition at every arrival, together
with persistent cluster identifiers. We independently sample $K$
candidate processes, indexed by $a\in[K]$. Each candidate is generated
as follows.
\begin{enumerate}
    \item With probability $w$, choose an index $i_a$ uniformly from
    $[m]$. This is an \emph{input candidate}. At time $t$, it outputs
    \[
        A_a(t)=P_{i_a}^t
    \]
    and uses the input labels $c_{i_a}(v)$ as its cluster identifiers.

    \item With probability $1-w$, choose a three-element set
    $Q_a\subseteq[m]$ uniformly at random. This is a
    \emph{triple candidate}. Initialize one persistent copy of
    \texttt{Agree} on the sequence
    \[
        H_{Q_a}(1),\ldots,H_{Q_a}(n).
    \]
    At time $t$, let $A_a(t)$ be the partition produced by this copy.
\end{enumerate}
Write $\lambda_t^a(v)$ for the local identifier of vertex $v$ in
candidate $a$ at time $t$.

All samples are drawn once, before processing the arrivals, and are
retained throughout the execution. Different candidates may sample
overlapping inputs or the same triple. Every candidate processes every
arrival, including while it is not selected for the output. In
particular, selecting or leaving a candidate never restarts its copy
of \texttt{Agree}.

\paragraph{Current costs and history scores.}
For each candidate $a$ and time $t$, define
\[
    f_a(t)=\operatorname{cost}_t(A_a(t)),
    \qquad
    G_a(t)=\max_{1\leq u\leq t} f_a(u),
    \qquad
    M_t=\min_{a\in[K]}G_a(t).
\]
Thus, $f_a(t)$ is the candidate's current cost, $G_a(t)$ is its largest
cost so far, and $M_t$ is the smallest available history score.
The algorithm computes current costs exactly against all $m$ input
partitions. Starting with $G_a(0)=0$, it updates history scores by
\[
    G_a(t)=\max\{G_a(t-1),f_a(t)\}.
\]

\paragraph{Selection phases.}
While $M_t=0$, the algorithm follows a candidate with zero history
score. Once $M_t$ becomes positive, selection proceeds in phases.
A positive phase begins with baseline $B=M_t$ and threshold
\[
    T=(1+\theta)B.
\]
Both $B$ and $T$ remain fixed throughout that phase. A candidate $a$
is \emph{eligible} at time $t$ if $G_a(t)<T$.

The algorithm retains its selected candidate while that candidate
remains eligible. If the selected candidate becomes ineligible while
$M_t<T$, the algorithm switches to an eligible candidate. If
$M_t\geq T$, it starts a new phase and selects a minimum-history
candidate. A phase change is processed before checking the eligibility
of the previously selected candidate.

\paragraph{Algorithm \texttt{oncon}.}
We now give the complete procedure. Whenever a choice is required,
break ties by choosing the smallest candidate index.
\begin{enumerate}
    \item \textbf{Initialization.}
    Sample and initialize the $K$ candidate processes according to
    the distribution above. Set $G_a(0)=0$ for every $a\in[K]$.
    Leave the selected index $a_0$ and the phase baseline $B$
    undefined.

    \item \textbf{Processing an arrival.}
    When vertex $v_t$ arrives, perform the following steps.
    \begin{enumerate}
        \item \textbf{Advance all candidates.}
        Supply the new vertex and its relevant input labels to every
        candidate process. For each triple candidate, reveal the new
        vertex and its incident signs to its persistent copy of
        \texttt{Agree}. Obtain $A_a(t)$ and its local identifiers
        $\lambda_t^a(v)$ for every $a\in[K]$.

        \item \textbf{Update all scores.}
        Compute
        \[
            f_a(t)=\frac{1}{m}\sum_{i=1}^{m}d(A_a(t),P_i^t)
        \]
        and set
        \[
            G_a(t)=\max\{G_a(t-1),f_a(t)\}
        \]
        for every $a\in[K]$. Then compute
        \[
            M_t=\min_{a\in[K]}G_a(t).
        \]

        \item \textbf{Select a candidate.}
        Apply the first applicable case below.
        \begin{enumerate}
            \item If $M_t=0$, retain the previous candidate if
            $a_{t-1}$ is defined and $G_{a_{t-1}}(t)=0$.
            Otherwise, select a candidate $a_t$ with $G_{a_t}(t)=0$.

            \item Otherwise, if $B$ is undefined or
            $M_t\geq(1+\theta)B$, begin a new positive phase:
            set
            \[
                B\leftarrow M_t,
                \qquad
                T\leftarrow(1+\theta)B,
            \]
            and choose
            \[
                a_t\in\operatorname*{arg\,min}_{a\in[K]}G_a(t).
            \]

            \item Otherwise, if $G_{a_{t-1}}(t)<T$, retain the
            previous candidate by setting $a_t=a_{t-1}$.

            \item Otherwise, switch to a candidate $a_t$ satisfying
            \[
                G_{a_t}(t)<T.
            \]
            Such a candidate exists because this case is reached
            only when $M_t<T$.
        \end{enumerate}

        \item \textbf{Produce the output.}
        Output the partition
        \[
            C_t=A_{a_t}(t),
        \]
        using the global cluster identifiers
        \[
            \lambda_t(v)=\bigl(a_t,\lambda_t^{a_t}(v)\bigr)
            \qquad\text{for every }v\in V_t.
        \]
    \end{enumerate}
\end{enumerate}

Including the candidate index in each global identifier distinguishes
identifiers belonging to different candidate processes. Changes to an
old vertex's output identifier can therefore arise either from a
switch of the selected candidate or from an internal identifier change
within the selected candidate. All candidate processes continue
running across switches and phase boundaries.

\subsection{Analysis}
\label{subsec:oncon-analysis}

In this subsection, we analyse the algorithm \texttt{oncon} and prove
\cref{thm:online-main}. To prove the theorem, we establish the following
lemmas, proving each immediately after its statement.

Throughout the analysis, fix an oblivious input instance with
$n\geq 2$ and $m\geq 3$. We use the notation and parameters from
\cref{sec:oncon-algorithm}. In particular,
\[
    0<\varepsilon<\frac{2}{3\eta+5},
    \qquad
    \theta=\frac{\varepsilon}{4}.
\]
Define
\[
    c=\eta+1,
    \qquad
    R_0=2-\frac{2}{3\eta+5},
    \qquad
    R=R_0+\frac{\varepsilon}{2}.
\]
These choices give $1<R_0<R<2$.

Recall that each history score $G_a(t)$ is nondecreasing in $t$.
Consequently, their minimum $M_t$ is also nondecreasing. These
properties hold for every realization of the sampled candidates.

Our first lemma bounds the entire history of a triple candidate
using its majority errors against a single reference partition.
For a time $t$, a partition $F$ of $V_t$, and a three-element set
$Q\subseteq[m]$, define
\[
    B_Q^t(F)
    =
    E^+(H_Q(t))\triangle E^+(F).
\]
Thus, $B_Q^t(F)$ consists of the pairs on which the majority graph
$H_Q(t)$ disagrees with $F$.

\begin{lemma}[History bound for a triple candidate]
\label{lem:oncon-history}
For every realization of the sampled candidates, every triple
candidate $a$, every time $t\in[n]$, and every partition $F$ of $V_t$,
\[
    G_a(t)
    \leq
    \cost_t(F)+c\,|B_{Q_a}^t(F)|.
\]
In particular,
\[
    G_a(t)
    \leq
    \OPT_t+c\,|B_{Q_a}^t(F_t^*)|.
\]
\end{lemma}

\begin{proof}
Fix a triple candidate $a$, a time $t$, and a partition $F$ of $V_t$.
Write $Q=Q_a$. For every earlier time $u\leq t$, define
\[
    F_u=F|_{V_u},
    \qquad
    B_u=B_Q^t(F)\cap\binom{V_u}{2}.
\]
Since signs between previously arrived vertices never change,
\[
    \err_{H_Q(u)}(F_u)=|B_u|.
\]
The partition $F_u$ is a feasible solution for correlation clustering
on $H_Q(u)$. Therefore, \cref{thm:agree-blackbox} gives
\[
    \err_{H_Q(u)}(A_a(u))
    \leq
    \eta\,\OPT_{\rm CC}(H_Q(u))
    \leq
    \eta\,|B_u|.
\]
Applying \eqref{eq:signed-comparison}, we obtain
\begin{equation}\label{eq:oncon-history-distance}
    d(A_a(u),F_u)
    \leq
    \err_{H_Q(u)}(A_a(u))
    +
    \err_{H_Q(u)}(F_u)
    \leq
    c\,|B_u|.
\end{equation}

Moreover, restriction cannot increase partition distance, so
\[
    \cost_u(F_u)
    =
    \frac1m\sum_{i=1}^m d(F|_{V_u},P_i^u)
    \leq
    \frac1m\sum_{i=1}^m d(F,P_i^t)
    =
    \cost_t(F).
\]
Using \eqref{eq:cost-comparison}, we conclude that
\begin{align*}
    f_a(u)
    &=
    \cost_u(A_a(u))\\
    &\leq
    \cost_u(F_u)+d(A_a(u),F_u)\\
    &\leq
    \cost_t(F)+c\,|B_u|\\
    &\leq
    \cost_t(F)+c\,|B_Q^t(F)|.
\end{align*}
The right-hand side is independent of $u$. Taking the maximum over
all $u\leq t$ proves the first assertion. Choosing $F=F_t^*$ proves
the second assertion.
\end{proof}

The next lemma combines the complementary expected history bounds
of input and triple candidates. The sampling probability $w$ is
chosen so that the additional terms in these two bounds cancel.

\begin{lemma}[Expected history score]
\label{lem:oncon-expectation}
For every fixed time $t\in[n]$ and candidate index $a\in[K]$,
\[
    \E[G_a(t)]\leq R_0\OPT_t.
\]
\end{lemma}

\begin{proof}
Fix a time $t$ and a candidate index $a$. For each pair
$e\in\binom{V_t}{2}$, let $k_e$ be the number of input partitions
whose decision on $e$ disagrees with that of $F_t^*$, and define
\[
    p_e=\frac{k_e}{m},
    \qquad
    X_t=\sum_{e\in\binom{V_t}{2}}p_e^2.
\]
Counting disagreements pair by pair gives
\[
    \sum_{e\in\binom{V_t}{2}}p_e
    =
    \frac1m\sum_{i=1}^m d(F_t^*,P_i^t)
    =
    \OPT_t.
\]

First, condition on $a$ being an input candidate. For any fixed
input index $i$, the quantity $\cost_u(P_i^u)$ is nondecreasing in
$u$: decisions on old pairs remain unchanged, and newly revealed
pairs contribute nonnegative additional cost. Hence an input
candidate satisfies $G_a(t)=f_a(t)$.

For two independent uniform input indices $i,j\in[m]$, the
partitions $P_i^t$ and $P_j^t$ disagree on a pair $e$ precisely when
one agrees with $F_t^*$ on $e$ and the other disagrees. This event
has probability $2p_e(1-p_e)$. Consequently,
\begin{align}
    \E[G_a(t)\mid a\text{ is an input candidate}]
    &=
    \frac{1}{m^2}
    \sum_{i=1}^m\sum_{j=1}^m d(P_i^t,P_j^t)\notag\\
    &=
    \sum_{e\in\binom{V_t}{2}}2p_e(1-p_e)\notag\\
    &=
    2\OPT_t-2X_t.
    \label{eq:oncon-input-history}
\end{align}

Next, condition on $a$ being a triple candidate. View its uniformly
sampled triple as three distinct input indices in a uniformly
random order. For a fixed pair $e$, the probability that the inputs
in two specified positions both disagree with $F_t^*$ on $e$ is
\[
    \frac{k_e(k_e-1)}{m(m-1)}
    \leq
    \left(\frac{k_e}{m}\right)^2
    =
    p_e^2.
\]
A majority error on $e$ requires at least two of the three sampled
inputs to disagree with $F_t^*$ on $e$. Taking a union bound over
the three pairs of sampled positions yields
\[
    \Prb\!\left[
        e\in B_{Q_a}^t(F_t^*)
        \,\middle|\,
        a\text{ is a triple candidate}
    \right]
    \leq 3p_e^2.
\]
By linearity of expectation,
\begin{equation}\label{eq:oncon-majority-error}
    \E\!\left[
        |B_{Q_a}^t(F_t^*)|
        \,\middle|\,
        a\text{ is a triple candidate}
    \right]
    \leq 3X_t.
\end{equation}
Applying \cref{lem:oncon-history}, we obtain
\begin{equation}
    \E[G_a(t)\mid a\text{ is a triple candidate}]
    \leq
    \OPT_t+3cX_t.
    \label{eq:oncon-triple-history}
\end{equation}

Combining \eqref{eq:oncon-input-history} and
\eqref{eq:oncon-triple-history} according to the sampling
probabilities gives
\begin{align*}
    \E[G_a(t)]
    &\leq
    w(2\OPT_t-2X_t)
    +(1-w)(\OPT_t+3cX_t)\\
    &=
    (1+w)\OPT_t
    +\bigl(-2w+3c(1-w)\bigr)X_t.
\end{align*}
Since
\[
    w=\frac{3c}{3c+2},
\]
we have
\[
    -2w+3c(1-w)=0,
    \qquad
    1+w
    =
    2-\frac{2}{3c+2}
    =
    R_0.
\]
Substitution proves the lemma.
\end{proof}

We next use independent sampling to ensure that a candidate with
a sufficiently small history score exists at every prefix.

\begin{lemma}[A good history at every prefix]
\label{lem:oncon-good-history}
With probability at least $1-\zeta$,
\[
    M_t\leq R\OPT_t
    \qquad\text{simultaneously for every }t\in[n].
\]
\end{lemma}

\begin{proof}
Fix a time $t\in[n]$. First suppose that $\OPT_t=0$. Every input partition $P_i^t$ then
coincides with $F_t^*$. Input candidates have history score zero,
and every triple has $B_Q^t(F_t^*)=\varnothing$. By
\cref{lem:oncon-history}, triple candidates also have history score
zero. Hence $M_t=0$ for every realization of the samples.

Now suppose that $\OPT_t>0$. By Markov's inequality and
\cref{lem:oncon-expectation}, each candidate $a$ satisfies
\[
    \Prb[G_a(t)>R\OPT_t]
    \leq
    \frac{\E[G_a(t)]}{R\OPT_t}
    \leq
    \frac{R_0}{R}.
\]
Therefore,
\[
    \Prb[G_a(t)\leq R\OPT_t]
    \geq
    1-\frac{R_0}{R}
    =
    \frac{\varepsilon}{2R}
    \geq
    \frac{\varepsilon}{4},
\]
where we used $R<2$.

At this fixed time $t$, each history score is a deterministic
function of its candidate's initial sample and the fixed input
instance. Every candidate runs independently of the selection rule,
and the copies of \texttt{Agree} are deterministic. Since the
initial candidate samples are independent, their success events
at time $t$ are independent. It follows that
\begin{align*}
    \Prb[M_t>R\OPT_t]
    &\leq
    \left(1-\frac{\varepsilon}{4}\right)^K\\
    &\leq
    \exp\left(-\frac{\varepsilon K}{4}\right)\\
    &\leq
    \frac{\zeta}{2n}.
\end{align*}
The same bound holds when $\OPT_t=0$, since the failure probability
is then zero.

Taking a union bound over all $n$ prefixes gives
\[
    \Prb\!\left[
        \exists t\in[n]:M_t>R\OPT_t
    \right]
    \leq
    n\cdot\frac{\zeta}{2n}
    =
    \frac{\zeta}{2}
    \leq
    \zeta.
\]
No independence between different prefixes is needed.
\end{proof}

The next lemma shows that the selection rule converts a small
minimum history score into a low-cost output. Its guarantee is
deterministic once the candidate processes are fixed.

\begin{lemma}[Selection guarantee]
\label{lem:oncon-selection}
For every realization of the sampled candidates, the selection
rule of \texttt{oncon} always has a valid candidate to choose and
satisfies
\[
    \cost_t(C_t)\leq(1+\theta)M_t
    \qquad\text{for every }t\in[n].
\]
\end{lemma}

\begin{proof}
Fix any realization of the sampled candidates. If $M_t=0$, a candidate with history score zero exists by the
definition of $M_t$. The selection rule chooses such a candidate,
whose current cost is also zero. Consider a positive phase with baseline $B$ and threshold
$T=(1+\theta)B$. At the start of the phase, the algorithm sets
$B=M_t$ and selects a candidate attaining this minimum. Its history
score is $B<T$, so it is eligible.

At a later time when the same phase continues, we have $M_t<T$;
otherwise, the algorithm would first start a new phase.
Consequently, a candidate attaining $M_t$ is eligible. Thus,
whenever the selected candidate becomes ineligible, an eligible
replacement exists.

Throughout the phase, the selected candidate has history score
strictly below $T$. Moreover, monotonicity of $M_t$ gives $B\leq M_t$.
Therefore,
\[
    \cost_t(C_t)
    =
    f_{a_t}(t)
    \leq
    G_{a_t}(t)
    <
    T
    =
    (1+\theta)B
    \leq
    (1+\theta)M_t.
\]
At a phase boundary, the baseline is reset before the output is
chosen, so the same argument applies to the new phase. Together
with the zero-score case, this proves the lemma.
\end{proof}

Finally, we bound the recourse by counting switches between
candidates and identifier changes within the candidate processes.

\begin{lemma}[Per-vertex recourse]
\label{lem:oncon-recourse}
For every realization of the sampled candidates and every vertex $v$,
\[
    \operatorname{rec}(v)
    =
    O_{\eta}\!\left(
        \varepsilon^{-2}
        \log\frac{2n}{\zeta}
        \log\frac{2mn}{\zeta}
    \right).
\]
\end{lemma}

\begin{proof}
Fix any realization of the sampled candidates. An \emph{outer switch}
is a time $t\geq 2$ at which $a_t\neq a_{t-1}$. Let $S$ be the total
number of outer switches, and let $h$ be the number of positive
phases. Define
\[
    H=
    1+
    \left\lceil
        \log_{1+\theta}\left(m\binom{n}{2}\right)
    \right\rceil.
\]
We first bound $h$ and $S$.

Each partition distance is an integer. Since consensus costs are
averages of $m$ such distances, every current cost, history score,
and minimum history score belongs to the grid
\[
    \left\{
        \frac{j}{m}:
        j=0,1,\ldots,m\binom{n}{2}
    \right\}.
\]
In particular, every positive phase baseline lies between $1/m$
and $\binom{n}{2}$.

If $h>0$, let $B_1,\ldots,B_h$ be the successive phase baselines.
The rule for starting a new phase gives
\[
    B_{j+1}\geq(1+\theta)B_j
    \qquad\text{for }j<h.
\]
Consequently,
\[
    \frac{(1+\theta)^{h-1}}{m}
    \leq
    B_h
    \leq
    \binom{n}{2},
\]
which implies $h\leq H$. This bound is immediate when $h=0$.
Since $\theta=\varepsilon/4$ and
$\log(1+\theta)\geq\theta/2$ for $0<\theta<1$, we have
\[
    H=O\!\left(\varepsilon^{-1}\log(2mn)\right).
\]

During the zero-score stage, a selected candidate is replaced only
when its history score becomes positive. That candidate can never
again have history score zero. Therefore, there are at most $K-1$
switches during this stage.

Within a positive phase, a selected candidate is replaced only
when its history score reaches the fixed threshold of that phase.
Since history scores do not decrease, it can never become eligible
again during the same phase. Hence there are at most $K-1$ switches
within each positive phase. Allowing one additional switch upon
entry into each positive phase gives
\[
    S
    \leq
    (K-1)+hK
    \leq
    K(1+H).
\]

We next account for changes within the candidate processes. For a
vertex $v=v_r$ and a candidate $a$, define its internal recourse by
\[
    \operatorname{rec}_a(v)
    =
    \sum_{t=r+1}^{n}
    \1[\lambda_t^a(v)\neq\lambda_{t-1}^a(v)].
\]
Input candidates have zero internal recourse. By
\cref{thm:agree-blackbox}, every triple candidate satisfies
\[
    \operatorname{rec}_a(v)
    \leq
    q_{\rm A}\lceil\log_2 n\rceil.
\]

Recall that the output identifier is
\[
    \lambda_t(v)
    =
    \bigl(a_t,\lambda_t^{a_t}(v)\bigr).
\]
If this identifier changes when the selected candidate changes,
charge the event to that outer switch. Otherwise, the selected
candidate remains the same, so its local identifier for $v$ must
have changed. Charge the event to that internal identifier change.
No internal event is charged more than once. Therefore,
\begin{align*}
    \operatorname{rec}(v)
    &\leq
    S+\sum_{a=1}^{K}\operatorname{rec}_a(v)\\
    &\leq
    K(1+H)+Kq_{\rm A}\lceil\log_2 n\rceil.
\end{align*}
This accounting uses the fact that all candidate processes persist
throughout the execution. Returning to a previously selected
candidate does not restart its internal recourse budget.

Finally, substitute
\[
    K=O\!\left(
        \varepsilon^{-1}\log\frac{2n}{\zeta}
    \right),
    \qquad
    H=O\!\left(
        \varepsilon^{-1}\log(2mn)
    \right).
\]
Since $\varepsilon<1$ and
$\log n,\log(2mn)\leq\log(2mn/\zeta)$, we obtain
\[
    \operatorname{rec}(v)
    =
    O_{\eta}\!\left(
        \varepsilon^{-2}
        \log\frac{2n}{\zeta}
        \log\frac{2mn}{\zeta}
    \right).
\]
Here the absolute constant $q_{\rm A}$ of the fixed implementation
of \texttt{Agree} is absorbed into the implicit constant.
The argument holds for every realization of the samples.
\end{proof}

We are now ready to prove \cref{thm:online-main}.

\begin{proof}[Proof of \cref{thm:online-main}]
The cases $n=1$ and $m\leq 2$ were handled before the description
of the algorithm: following the indicated input is optimal and
incurs no recourse. Assume that $n\geq 2$ and $m\geq 3$.
By \cref{lem:oncon-good-history}, with probability at least
$1-\zeta$,
\[
    M_t\leq R\OPT_t
    \qquad\text{for every }t\in[n].
\]
On this event, \cref{lem:oncon-selection} gives
\[
    \cost_t(C_t)
    \leq
    (1+\theta)M_t
    \leq
    (1+\theta)R\OPT_t
    \qquad\text{for every }t\in[n].
\]
Since $\theta=\varepsilon/4$, $R=R_0+\varepsilon/2$, and $R<2$,
\[
    (1+\theta)R
    =
    R_0+\frac{\varepsilon}{2}
    +\frac{\varepsilon R}{4}
    <
    R_0+\varepsilon.
\]
It follows that, simultaneously for all $t\in[n]$,
\[
    \cost_t(C_t)
    \leq
    \left(
        2-\frac{2}{3\eta+5}+\varepsilon
    \right)\OPT_t.
\]
This includes prefixes with $\OPT_t=0$.

Independently of this event, \cref{lem:oncon-recourse} bounds the
recourse of every vertex for every realization of the random samples.
These are the two guarantees asserted in \cref{thm:online-main}.
\end{proof}


\section{Fully Dynamic Consensus Clustering}
\label{sec:dynamic}
In this section, we prove our guarantee for fully dynamic
consensus clustering. Our algorithm maintains a
$(1.6944+\varepsilon)$-approximation under vertex insertions
and deletions, with polylogarithmic expected amortized
update time and logarithmic time for membership and
same-cluster queries. We restate the main result. 

\ImplicitDynamic*

\kushagra{I think in the theorem statements it is better to write the simplified forms of the update time, so that it is easier to read. Maybe we can write $O(\log^2 U \, \, \log^2 N)$.}


% \kushagra{Currently, $C_t$ is only defined for the online model. It needs to be defined for the dynamic model too. Somewhere we need to write "Let, $C_t$ denote the partition of $V_t$ maintained by the algorithm after the $t$th update. - defined in the preliminaries." }

\subsection{The Offline Algorithm}
\label{subsec:offline}
We first study the problem on a fixed vertex set.
Our offline algorithm generates candidate clusterings from
sampled inputs and selects one using an independent
evaluation sample. The dynamic algorithm will maintain
these candidates and their estimated costs under vertex
updates. We prove the following guarantee.

\begin{theorem}[Offline consensus clustering]
\label{thm:offline}
There is a randomized offline algorithm that, given
$m$ input clusterings of $n$ vertices, has the following
guarantees. For every $\varepsilon\in(0,1]$ and
$\zeta\in(0,1)$, with probability at least $1-\zeta$,
it returns a clustering $C$ satisfying
\[
    \operatorname{cost}(C)
        \le (1.6944+\varepsilon)\operatorname{OPT}
\]
in time $O\left(
        n\varepsilon^{-3}
        \log^2\frac{n}{\varepsilon\zeta}
    \right)$.
\end{theorem}

\kushagra{Same. Simplify the running time to $O(n \log^2 n).$}

\subsubsection{Algorithm}
\label{subsec:offline-overview}
\label{subsec:offline-buckets}

Let $\mathcal P=(P_1,\ldots,P_m)$ be the input clusterings
of a vertex set $V$ of size $n$. Each cluster $P \in \mathcal{P}$ is represented by a cluster label in $[n]$, with constant-time access to each label. If a vertex $v \in P_i$ then we define the cluster label of that vertex as $c_i(v) \in [n]$. Set
\[
    L=\log\frac{n+2}{\varepsilon\zeta},
    \qquad
    T=\Theta(\varepsilon^{-1}L),
    \qquad
    s=\Theta(\varepsilon^{-2}L),
\]
with sufficiently large constants. Each trial produces two candidate clusterings. We sample
seven input indices independently and uniformly from $[m]$,
with replacement. For the sampled tuple $Q=(i_1,\ldots,i_7)$,
define the majority graph $G_Q$ by
\[
    \{u,v\}\in E(G_Q)
    \quad\Longleftrightarrow\quad
    \sum_{j=1}^{7}
        \mathbf{1}[c_{i_j}(u)=c_{i_j}(v)]\ge 4.
\]
We obtain the first candidate by running \textsc{Pivot}~\cite{ACN2008} on $G_Q$. At each step, \textsc{Pivot} chooses a uniformly random vertex from the remaining vertices. It forms one cluster from that vertex and all its remaining neighbors, then removes the cluster and repeats until no vertices remain. The second is an independently and uniformly sampled input partition. We perform $T$ independent trials and collect their $2T$ candidates in a multiset $\mathcal Y$.

To run \textsc{Pivot} efficiently, we represent neighborhoods
using buckets. Fix the sampled tuple $Q=(i_1,\ldots,i_7)$, and let
\[
    \mathcal{R}=\binom{[7]}{4}
\]
be the family of all four-element subsets of the sampled
positions. For each $R\in\mathcal{R}$ and vertex $u\in V$,
define
\[
    B_R(u)
    =
    \left\{
        v\in V :
        c_{i_j}(v)=c_{i_j}(u)
        \text{ for every }j\in R
    \right\},
\]
where $c_{i_j}(x)$ denotes the label of the cluster containing
$x$ in the input partition $P_{i_j}$. Thus, $B_R(u)$ consists
of the vertices that are clustered together with $u$ in
all four sampled positions belonging to $R$.

Two distinct vertices are adjacent in $G_Q$ if and only if
they share a bucket for some $R\in\mathcal{R}$: agreeing in
at least four of the seven sampled positions is equivalent
to agreeing in every position of some $R\in\mathcal{R}$.
Consequently,
\[
    N_{G_Q}[u]
    =
    \bigcup_{R\in\mathcal{R}} B_R(u).
\]
Since $|\mathcal{R}|=\binom{7}{4}=35$, each vertex has
35 bucket memberships, one for each choice of $R$.
A pivot's remaining closed neighborhood can therefore
be enumerated from its bucket lists.
For selection, draw an independent evaluation sample
$S=(P_{j_1},\ldots,P_{j_s})$, with all indices sampled
independently and uniformly from $[m]$, with replacement.
The estimated cost of a candidate $C$ is
\[
    \widehat{\operatorname{cost}}_S(C)
        =\frac{1}{s}\sum_{\ell=1}^{s}d(C,P_{j_\ell}).
\]
We return a candidate of minimum estimated cost. To evaluate these distances, we use the Mirkin-distance
identity. For two clusterings $C,D$, let $Z$ be the partition
formed by the nonempty intersections of their clusters.
Then
\[
    d(C,D)
        =\sum_{A\in C}\binom{|A|}{2}
         +\sum_{B\in D}\binom{|B|}{2}
         -2\sum_{H\in Z}\binom{|H|}{2}.
\]
Thus, distances can be computed from cluster sizes and
intersection sizes. ~\cref{alg:offline-consensus}
gives the complete procedure.
\kushagra{Need to write an English description of ~\cref{alg:offline-consensus}.}
\RestyleAlgo{boxed}
\begin{algorithm}[t]
\caption{Near-linear-time offline consensus clustering}
\label{alg:offline-consensus}
\DontPrintSemicolon
\KwIn{Input clusterings $\mathcal P=(P_1,\ldots,P_m)$ on $V$,
      accuracy $\varepsilon\in(0,1]$, and failure probability
      $\zeta\in(0,1)$.}
\KwOut{A clustering $\widehat C$ of $V$.}

Set $T,s$ as above and initialize $\mathcal Y\gets\emptyset$\;

\For{$a=1,\ldots,T$}{
    Sample $Q_a$ as seven independent uniform indices
    from $[m]$, with replacement\;
    Construct the bucket lists for $G_{Q_a}$\;
    Initialize the remaining-vertex array $A\gets V$
    and the clustering $C_a\gets\emptyset$\;

    \While{$A\neq\emptyset$}{
        Choose a uniformly random vertex $p$ from $A$\;
        Enumerate
        $\displaystyle
            D=\bigcup_{R\in\binom{[7]}{4}}
                \bigl(B_R(p)\cap A\bigr)$,
        discarding duplicates\;
        Add $D$ as a cluster of $C_a$\;
        Remove every vertex of $D$ from its bucket lists
        and from $A$\;
    }

    Independently sample a uniform index $J_a\in[m]$\;
    Add $C_a$ and $P_{J_a}$ to $\mathcal Y$\;
}

Draw $S$ as $s$ independent uniform input partitions,
with replacement, independently of all trials\;
Compute $\widehat{\operatorname{cost}}_S(C)$
for every $C\in\mathcal Y$ using the Mirkin-distance identity\;

\KwRet{$\displaystyle
    \widehat C\in
    \operatorname*{arg\,min}_{C\in\mathcal Y}
        \widehat{\operatorname{cost}}_S(C)$}\;
\end{algorithm}

\begin{lemma}
\label{lem:offline-running-time}
\cref{alg:offline-consensus} runs in
$O(nTs)=O(n\varepsilon^{-3}L^2)$ time.
\end{lemma}

\begin{proof}
For each $R\in\binom{[7]}{4}$, radix-sorting the vertices
by their four-coordinate bucket keys constructs the bucket
lists in $O(n)$ time. Since there are $35$ choices of $R$,
constructing all buckets takes $O(n)$ time per trial.

Store each bucket as a doubly linked list, with deletion
handles for its members. A swap array supports uniform
sampling and deletion of remaining vertices in constant
time, and timestamps remove duplicates during enumeration.
The neighborhood identity ensures that each iteration
performs exactly one step of \textsc{Pivot}. Every enumerated
vertex belongs to the cluster being removed and appears
in at most $35$ inspected lists. Finding and removing a
cluster $D$ therefore takes $O(|D|)$ time. Summing over
the clusters gives $O(n)$ time per trial, and hence
$O(nT)$ time for candidate generation.

Assign each candidate's clusters labels in $[n]$.
Counting these labels and radix-sorting the label pairs
$(c_C(v),c_D(v))$ gives all cluster and intersection sizes
in $O(n)$ time. The Mirkin-distance identity therefore
computes each distance in $O(n)$ time. Evaluating all
$2T$ candidates against the $s$ sampled inputs takes
$O(nTs)$ time. The total running time is consequently
\[
    O(nT+nTs)
        =O(nTs)
        =O(n\varepsilon^{-3}L^2).
\]
\end{proof}

\subsubsection{Analysis}
\label{subsec:offline-analysis}

We bound the expected cost of the better of the two candidates
produced by one trial, averaging over the sampled inputs and the
random choices of \textsc{Pivot}. We then use independent repetition
and evaluation to prove Theorem~\ref{thm:offline}.

For the analysis, write $\operatorname{cost}(C)$ for
$\operatorname{cost}_{\mathcal{P}}(C)$ and $\mathrm{OPT}$ for
$\mathrm{OPT}(\mathcal{P})$. Fix an optimal clustering $C^\star$. For each pair $e=\{u,v\}$, define
\[
    w_e^+
    = \frac{1}{m}\sum_{i=1}^{m}
      \mathbf{1}\bigl[c_i(u)=c_i(v)\bigr],
    \qquad
    w_e^- = 1-w_e^+.
\]
Thus, placing the endpoints of $e$ together contributes $w_e^-$
to the consensus cost, while separating them contributes $w_e^+$. Let $G^\star$ be the majority graph of all input partitions,
with ties broken arbitrarily, and set
\[
    m_e=\min\{w_e^+,w_e^-\},
    \qquad
    u_e=1-2m_e.
\]
Following the sign of $e$ in $G^\star$ costs $m_e$, while
opposing it costs $1-m_e$. Hence, $u_e$ is the additional cost
of opposing the majority decision. Let $c_e^\star$ denote the contribution of $e$ to the cost of
$C^\star$. Since every clustering either joins or separates
the endpoints of each pair,
\[
    c_e^\star\in\{m_e,1-m_e\},
    \qquad
    \mathrm{OPT}
    =\sum_{e\in\binom{V}{2}}c_e^\star
    \geq \sum_{e\in\binom{V}{2}}m_e.
\]

A triplet is \emph{bad} in a graph if exactly two of its three
pairs are edges. A partition cannot realize this pattern because
joining the endpoints of both edges forces all three vertices
into the same cluster. Write $\mathcal B(G)$ for the set of
bad triplets in $G$.

Fix the seven sampled input partitions indexed by $Q$, and let
$G_Q$ be their majority graph. We analyze \textsc{Pivot} on this
fixed graph, taking probability only over its random pivot choices.
Recall that each step chooses a uniformly random remaining vertex
and removes that vertex and its remaining neighbors as one cluster.
For each such triplet $t$, let $p_t$ be the probability that
a vertex of $t$ is chosen as a pivot while all three vertices remain.

Let $C$ be the output of \textsc{Pivot} on the fixed sampled graph
$G_Q$. For each pair $e$, let $s_e$ be its contribution to the
consensus cost when the output follows its sign in $G_Q$.
Thus, $s_e=w_e^-$ if $e$ is an edge, and $s_e=w_e^+$ otherwise.
Opposing its sign costs $1-s_e$.

\begin{lemma}[Exact pair identity]
\label{lem:pair-disagreement}
\label{lem:pair-cost}
For each pair $e$, the probability that \textsc{Pivot} decides
$e$ against its sign in $G_Q$ is
\[
    D_e
    = \frac{1}{3}
      \sum_{\substack{t\in\mathcal{B}(G_Q)\\
                      e\in\binom{t}{2}}} p_t.
\]
Consequently, taking expectation over the random pivots,
\[
    \mathbb{E}_{\mathrm{piv}}[\operatorname{cost}(C)]
    = \sum_e s_e+\sum_e D_e(1-2s_e).
\]
\end{lemma}

\kushagra{Is this a new lemma? Isn't this lemma proved in ACN paper?}

\begin{proof}
Consider $e=\{u,v\}$. The first pivot whose cluster contains
either endpoint decides the pair.
If that pivot is $u$ or $v$, the decision follows the sign of $e$.
Suppose a third vertex $w$ pivots while both endpoints remain.
If $uv$ is a nonedge, the decision contradicts its sign when
$w$ is adjacent to both endpoints, placing them in the same cluster. If $uv$ is an edge, the decision contradicts its sign when
$w$ is adjacent to exactly one endpoint, separating them.
These are precisely the cases in which $\{u,v,w\}$ is a bad
triplet. While all three vertices remain, each is equally likely
to be chosen as the next pivot. Thus, $w$ causes this disagreement
with probability $p_t/3$.
For a fixed pair, these events are disjoint because the first
such event permanently decides the pair. Summing gives the formula
for $D_e$.

The pair $e$ follows its sign
with probability $1-D_e$ and opposes it with probability $D_e$.
Its expected contribution is therefore
\[
    (1-D_e)s_e+D_e(1-s_e)
    =s_e+D_e(1-2s_e).
\]
Summing over all pairs proves the identity.
\end{proof}

Let $F_e$ indicate that $G_Q$ and $G^\star$ disagree on the
sign of $e$. When $F_e=0$, following the sampled sign agrees
with the full majority decision and costs $m_e$. When $F_e=1$,
it opposes the full majority decision and costs $1-m_e=m_e+u_e$.
Therefore,
\[
    s_e=m_e+F_eu_e,
    \qquad
    1-2s_e=u_e(1-2F_e).
\]
Write
\[
    T_Q=\mathcal B(G_Q),
    \qquad
    B=\mathcal B(G_Q)\cap\mathcal B(G^\star).
\]
Substituting the expressions for $s_e$ and $D_e$ into
Lemma~\ref{lem:pair-cost} and grouping by triplet gives
\[
    \mathbb E_{\mathrm{piv}}[\operatorname{cost}(C)]
        =\sum_e m_e+\sum_e F_eu_e
         +\sum_{t\in T_Q}\frac{p_t}{3}
             \sum_{e\in\binom{t}{2}}u_e(1-2F_e).
\]
We use the full majority signs on the shared bad triplets as
a reference. For each $t\in T_Q$, define
\[
    X_t=
        \sum_{e\in\binom{t}{2}}u_e(1-2F_e)
        -\mathbf 1[t\in B]
            \sum_{e\in\binom{t}{2}}u_e.
\]
For $t\in B$, $X_t$ measures the change in the triplet's
contribution caused by wrong sampled signs. For
$t\in T_Q\setminus B$, it equals the entire triplet contribution.
Substituting this expression gives
\begin{align}
    \mathbb E_{\mathrm{piv}}[\operatorname{cost}(C)]
        &=\sum_e m_e
          +\sum_{t\in B}\frac{p_t}{3}
              \sum_{e\in\binom{t}{2}}u_e \notag\\
        &\quad+\sum_e F_eu_e
          +\sum_{t\in T_Q}\frac{p_t}{3}X_t.
    \label{eq:offline-decomposition}
\end{align}
The first line uses the full majority signs and the bad
triplets shared by both graphs. The second line accounts for
the direct cost of wrong sampled signs and the change in
triplet contributions. If all three signs of $t\in T_Q$ are
sampled correctly, then $t\in B$ and $X_t=0$.

A uniformly sampled input disagrees with the sign of $e$
in $G^\star$ with probability $m_e$. Since the seven inputs
are sampled independently, their majority disagrees with
$G^\star$ exactly when at least four make the minority
decision. Define
\[
    \phi_7(x)
    =\Pr[\operatorname{Bin}(7,x)\geq4]
    =\sum_{j=4}^{7}\binom{7}{j}x^j(1-x)^{7-j}.
\]
Thus, $\Pr[F_e=1]=\phi_7(m_e)$, and linearity of expectation
gives
\[
    \mathbb E_Q\!\left[\sum_e F_eu_e\right]
        =\sum_e\phi_7(m_e)u_e.
\]

We next show that each bad triplet contributes at most one
positive unit to the sampling correction.

\begin{lemma}
\label{lem:triplet-correction}
For every $t\in T_Q$, we have $X_t\leq 1$.
Moreover, if $t\in B$, then
\[
    X_t=-2\sum_{\substack{e\in\binom{t}{2}\\F_e=1}}u_e
    \leq 0.
\]
\end{lemma}

\begin{proof}
If $t$ is bad in both graphs, each pair contributes
$u_e((1-2F_e)-1)$ to $X_t$. A correctly sampled pair has
$F_e=0$ and contributes zero. A wrongly sampled pair has
$F_e=1$ and contributes $-2u_e$. Since $u_e\geq0$, we have
$X_t\leq0$.


Now suppose $t$ is bad in $G_Q$ but not in $G^\star$.
Write $t=\{u,v,w\}$, with edges $uw,vw$ and nonedge $uv$
in $G_Q$.  Since $uv$ is a nonedge in $G_Q$, following its sign separates the endpoints and costs $s_{uv}=w_{uv}^+=1-w_{uv}^-$.
The pairs $uw,vw$ are edges, so following their signs joins
the endpoints and costs $s_{uw}=w_{uw}^-$ and $s_{vw}=w_{vw}^-$.
Thus,
\[
    s_{uv}+s_{uw}+s_{vw}
    =1+\underbrace{
        w_{uw}^-+w_{vw}^- -w_{uv}^-
      }_{\geq0}
    \geq1.
\]
Since $t$ is not bad in $G^\star$, there is no reference
contribution to subtract. Thus,
\[
    X_t
    =3-2(s_{uv}+s_{uw}+s_{vw})
    \leq1.
\]
\end{proof}
We now assign the positive triplet corrections to mis-sampled pairs. Every triplet with $X_t>0$ contains a mis-sampled pair, since
$X_t=0$ when all three signs are sampled correctly.
Assign each such triplet to one of its mis-sampled pairs, and define
\[
    \Gamma_e
    =\sum_{t\text{ assigned to }e}\frac{p_t}{3}X_t.
\]
Discarding the nonpositive corrections gives
\[
    \sum_{t\in T_Q}\frac{p_t}{3}X_t\leq\sum_e\Gamma_e.
\]
By Lemma~\ref{lem:triplet-correction}, each assigned triplet has
$0<X_t\leq1$. Since every triplet assigned to $e$ contains $e$,
\[
    0\leq\Gamma_e
    \leq\sum_{\substack{t\in T_Q\\e\in\binom{t}{2}}}\frac{p_t}{3}
    =D_e\leq1.
\]
Moreover, $\Gamma_e=0$ whenever $F_e=0$, because only mis-sampled
pairs receive a charge. Next, unanimous sampled decisions allow us to sharpen the charge bound.

\begin{lemma}
\label{lem:sampled-unanimity}
If all seven sampled inputs make the same decision on a pair $e$,
then no bad triplet of $G_Q$ contains $e$. Consequently, $\Gamma_e=0$.
\end{lemma}

\begin{proof}
Write $e=\{u,v\}$.
If all sampled inputs separate $u,v$, each input can place a third
vertex $w$ with at most one endpoint.
Thus, the agreement counts for $uw$ and $vw$ sum to at most seven.
They cannot both reach four, so $w$ cannot be adjacent to both
endpoints in $G_Q$. This excludes a bad triplet containing
the nonedge $uv$.

If all sampled inputs join $u,v$, then in every input, $w$ is
joined to $u$ if and only if it is joined to $v$.
Hence $uw$ and $vw$ have identical agreement counts, so $w$ is
adjacent to both endpoints or neither. This excludes a bad triplet
containing the edge $uv$.
\end{proof}

The event that all seven sampled inputs disagree with $G^\star$
on $e$ has probability $m_e^7$. It is a sampling error, but the
lemma shows that its triplet charge is zero.
Since $0\leq\Gamma_e\leq1$ and only sampling errors receive charges,
\[
    \mathbb{E}_Q[\Gamma_e]
    \leq
    \Pr[F_e=1\text{ and the sampled decision is not unanimous}]
    =\phi_7(m_e)-m_e^7.
\]

For $0\leq x\leq1/2$, define
\[
    \psi_7(x)=
    \begin{cases}
        \displaystyle
        \frac{\phi_7(x)(2-2x)-x^7}{x},
        & 0<x\leq1/2,\\[6pt]
        0, & x=0.
    \end{cases}
\]
Combining the direct sign-error cost and the triplet charges
gives
\begin{align}
    &\mathbb E_Q\!\left[
        \sum_e F_eu_e
        +\sum_{t\in T_Q}\frac{p_t}{3}X_t
    \right] \notag\\
    &\qquad\leq \sum_e
        \bigl[\phi_7(m_e)u_e+\phi_7(m_e)-m_e^7\bigr] \notag\\
    &\qquad=\sum_e
        \bigl[\phi_7(m_e)(2-2m_e)-m_e^7\bigr] \notag\\
    &\qquad=\sum_e\psi_7(m_e)m_e.
    \label{eq:sampling-price}
\end{align}

We now bound the better of the two candidates using a fixed
weighted average of their costs. Let $C'=P_J$ be the independently sampled input partition.
For a pair $e$, it joins the endpoints with probability $w_e^+$
and then pays $w_e^-$; otherwise, it separates them and pays $w_e^+$.
Its expected contribution is therefore
\[
    z_e=w_e^+w_e^-+w_e^-w_e^+
       =2m_e(1-m_e).
\]

For any fixed $0<\beta<1$, the smaller of two costs is at most
their weighted average. Hence,
\[
    \mathbb{E}\bigl[\min\{\operatorname{cost}(C),
                          \operatorname{cost}(C')\}\bigr]
    \leq
    \beta\,\mathbb{E}[\operatorname{cost}(C)]
    +(1-\beta)\sum_e z_e.
\]
It therefore suffices to choose $\beta$ and bound the right-hand
side by $\gamma\,\mathrm{OPT}$.
We do this by combining the first line of
\eqref{eq:offline-decomposition} with the sampling bound
\eqref{eq:sampling-price}. Let $D'_e$ collect the contributions from common bad triplets:
\[
    D'_e
    =\sum_{\substack{t\in B\\e\in\binom{t}{2}}}\frac{p_t}{3}
    \leq D_e\leq1.
\]
Using $D'_e$, the first line of \eqref{eq:offline-decomposition}
can be written as
\begin{align*}
    &\sum_e m_e+\sum_{t\in B}\frac{p_t}{3}
        \sum_{e\in\binom{t}{2}}u_e\\
    &\qquad=\sum_e(1-D'_e)m_e
      +\sum_{t\in B}\frac{p_t}{3}
       \sum_{e\in\binom{t}{2}}(1-m_e).
\end{align*}
This follows because each pair contributes
$(1-D'_e)m_e+D'_e(1-m_e)=m_e+D'_eu_e$. The approximation bound now reduces to two inequalities: one for individual pairs and one for bad triplets.

\begin{lemma}[Charging inequalities]
\label{lem:charging-inequalities}
Let
\begin{equation}
    \beta=\frac{6144}{20101},
    \qquad
    \gamma=\frac{34058}{20101}=2-\beta<1.6944
    \label{eq:approx-constants}
\end{equation}
For every $x\in[0,1/2]$,
\begin{equation}
    \beta\bigl[x+\psi_7(x)x\bigr]
    +(1-\beta)2x(1-x)
    \leq\gamma x
    \label{eq:joint-pair}
\end{equation}
For every bad triplet $t$ in $G^\star$,
\begin{equation}
    \sum_{e\in\binom{t}{2}}
    \Bigl[
        \beta\bigl(1-m_e+\psi_7(m_e)m_e\bigr)
        +(1-\beta)z_e
    \Bigr]
    \leq\gamma\sum_{e\in\binom{t}{2}}c_e^\star
    \label{eq:joint-triplet}
\end{equation}
\end{lemma}

The proof of Lemma~\ref{lem:charging-inequalities} is given in
Appendix~\ref{app:offline-charging}.

\begin{lemma}[Joint charging]
\label{lem:joint-charging}
For $\beta$ and $\gamma$ as in \eqref{eq:approx-constants},
\[
    \beta\,\mathbb{E}[\operatorname{cost}(C)]
    +(1-\beta)\,\mathbb{E}[\operatorname{cost}(C')]
    \leq\gamma\,\mathrm{OPT}.
\]
\end{lemma}

\begin{proof}
Fix $Q$. Lemma~\ref{lem:charging-inequalities} gives
\eqref{eq:joint-pair} and \eqref{eq:joint-triplet} for the
chosen constants. For each pair $e$, substitute $x=m_e$ into
\eqref{eq:joint-pair} and multiply both sides by $1-D'_e$.
For each common bad triplet $t\in B$, multiply both sides
of \eqref{eq:joint-triplet} by $p_t/3$.
Add all these inequalities. We obtain
\begin{align*}
    &\sum_e(1-D'_e)
      \Bigl[\beta m_e+\beta\psi_7(m_e)m_e+(1-\beta)z_e\Bigr]\\
    &\quad+
      \sum_{t\in B}\frac{p_t}{3}
      \sum_{e\in\binom{t}{2}}
      \Bigl[\beta(1-m_e)+\beta\psi_7(m_e)m_e
            +(1-\beta)z_e\Bigr]\\
    &\leq
      \gamma\sum_e(1-D'_e)m_e
      +\gamma\sum_{t\in B}\frac{p_t}{3}
       \sum_{e\in\binom{t}{2}}c_e^\star.
\end{align*}
For each pair $e$, the pair inequality has weight $1-D'_e$
and the triplet inequalities have total weight $D'_e$. The
$\beta m_e$ and $\beta(1-m_e)$ terms therefore contribute
$\beta[(1-D'_e)m_e+D'_e(1-m_e)]=\beta(m_e+D'_eu_e)$.
The term $\beta\psi_7(m_e)m_e+(1-\beta)z_e$ receives total
weight $(1-D'_e)+D'_e=1$. On the right, $m_e\leq c_e^\star$,
so the same weights bound the sum by $\gamma\sum_e c_e^\star
=\gamma\,\mathrm{OPT}$.
Consequently,
\[
    \beta\sum_e(m_e+D'_eu_e)
    +\beta\sum_e\psi_7(m_e)m_e
    +(1-\beta)\sum_e z_e
    \leq\gamma\,\mathrm{OPT}.
\]

Finally, average over $Q$. By
\eqref{eq:offline-decomposition} and \eqref{eq:sampling-price},
\[
    \mathbb E[\operatorname{cost}(C)]
    \leq \mathbb E_Q\!\left[\sum_e(m_e+D'_eu_e)\right]
        +\sum_e\psi_7(m_e)m_e.
\]
Since $\sum_e z_e=\mathbb E[\operatorname{cost}(C')]$,
the weighted expected cost is at most $\gamma\,\mathrm{OPT}$.
\end{proof}

The smaller cost is at most the weighted average, so
Lemma~\ref{lem:joint-charging} gives
\begin{equation}
    \mathbb E\!\left[
        \min\{\operatorname{cost}(C),\operatorname{cost}(C')\}
    \right]\leq\gamma\,\mathrm{OPT}.
    \label{eq:offline-single-trial}
\end{equation}

Theorem~\ref{thm:offline} follows from \eqref{eq:offline-single-trial}
by independent repetition and selection using the evaluation sample.

\begin{proof}[Proof of Theorem~\ref{thm:offline}]
The running time follows from Lemma~\ref{lem:offline-running-time}.
If $\mathrm{OPT}=0$, the input partitions coincide, so estimated
costs are exact and a sampled-input candidate is optimal.
Otherwise, set $\delta=\varepsilon/4$.

By \eqref{eq:offline-single-trial} and Markov's inequality,
a trial produces a candidate of cost at most
$(\gamma+\delta)\,\mathrm{OPT}$ with probability $\Omega(\delta)$.
The $T$ independent trials therefore contain such a candidate
with probability at least $1-\zeta/2$.

Condition on $\mathcal Y$ and let $C_0$ be a candidate of minimum
true cost. For any candidate $C$ with
$\operatorname{cost}(C)>(1+\delta)\operatorname{cost}(C_0)$,
write $X(P)=d(C,P)-d(C_0,P)$ and $R=d(C,C_0)$.
The triangle inequality gives $|X(P)|\leq R$, and averaging it
over the inputs gives $R\leq\operatorname{cost}(C)+
\operatorname{cost}(C_0)$. Thus, for a uniform input $P$,
\[
    \mathbb E[X(P)]
        =\operatorname{cost}(C)-\operatorname{cost}(C_0)
        >\frac{\delta R}{2+\delta}.
\]
Since the evaluation sample is independent of $\mathcal Y$,
Hoeffding's inequality bounds the probability that $C$ has
estimated cost at most that of $C_0$ by $\exp(-s\delta^2/18)$.
A union bound over the $2T$ candidates gives
$\operatorname{cost}(\widehat C)\leq
(1+\delta)\operatorname{cost}(C_0)$ except with probability
$2T\exp(-s\delta^2/18)\leq\zeta/2$.

Combining the two events, with probability at least $1-\zeta$,
\[
    \operatorname{cost}(\widehat C)
    \leq (1+\delta)(\gamma+\delta)\,\mathrm{OPT}
    \leq (1.6944+\varepsilon)\,\mathrm{OPT},
\]
where the last inequality uses $\delta=\varepsilon/4$ and
$\varepsilon\leq1$.
\end{proof}

\subsection{The Fully Dynamic Algorithm}
\label{subsec:dynamic-algorithm}

\kushagra{This section is again not well written. Just like the previous subsection for Offline Algorithm here also make two sub-subsections Algorithm and Analysis.  At the start of the subsection mention the theorem statement you want to prove just like you did for the offline algorithm. Just give prompt to GPT it will organize the writing like this.}

We work in the model of Section~\ref{subsec:dynamic-model}.
Given the update bound $U$, accuracy $\varepsilon$, and failure
probability $\zeta$ in Theorem~\ref{thm:implicit-dynamic}, set
\[
    \Lambda=\log\frac{U+2}{\varepsilon\zeta},
    \qquad
    T=\Theta(\varepsilon^{-1}\Lambda),
    \qquad
    s=\Theta(\varepsilon^{-2}\Lambda),
\]
with sufficiently large constants. Before the updates begin, draw the
$T$ independent trial samples $(Q_a,J_a)$ and an independent evaluation
sample $S=(P_{j_1},\ldots,P_{j_s})$ as in the offline algorithm.
We retain these sampled indices throughout the sequence.

At time $t$, trial $a$ maintains a \textsc{Pivot} candidate $C_{a,t}$
on $G_{Q_a,t}$, the majority graph of its seven sampled partitions
restricted to $V_t$, and the input candidate $C'_{a,t}=P_{J_a}|_{V_t}$.
Let $\mathcal Y_t$ be the family of these $2T$ candidates. The trials
maintain their partitions independently of $S$. After each update,
we store the index of a candidate $C_t$ that minimizes
\[
    \widehat{\operatorname{cost}}_{S,t}(C)
    =\frac{1}{s}\sum_{\ell=1}^{s}d(C,P_{j_\ell}|_{V_t})
    \qquad\text{over }C\in\mathcal Y_t.
\]

\paragraph{Maintaining a \textsc{Pivot} candidate.}
In trial $a$, a vertex $v$ receives an independent uniform priority
$\rho_a(v)\in(0,1)$ upon insertion and keeps it until deletion.
A reinsertion receives a fresh priority. We use the random-ranking
model, in which priority comparisons take constant time.
The candidate $C_{a,t}$ is defined by choosing the lowest-priority
remaining vertex at each step of \textsc{Pivot}. Since the update
sequence is oblivious, at each fixed time the priorities on $V_t$
are independent and uniform, even conditional on $Q_a$. After a
pivot and its neighbors are removed, the remaining order is still
uniform. Thus this construction has the offline \textsc{Pivot}
distribution on $G_{Q_a,t}$.

For the maintenance analysis, fix one trial and write $Q$ and $\rho$
for its sample and priorities. We use the offline bucket representation.
Each vertex belongs to $K=\binom{7}{4}=35$ buckets, and two vertices are
adjacent exactly when they share a bucket. Inserting or deleting a vertex
changes only its own bucket memberships. The repair works for any such
representation with at most $K$ clique buckets per vertex and
constant-time comparisons of bucket keys.

\paragraph{Pivots, owners, and blocking certificates.}
Let $M$ be the greedy maximal independent set obtained by scanning vertices in increasing priority and selecting a vertex whenever it has no previously selected neighbor. These selected vertices are exactly the pivots of \textsc{Pivot} in that order.
Since $M$ is independent and each bucket is a clique, each bucket
contains at most one pivot and each vertex has at most $K$ neighboring
pivots. The cluster owner of $v$ is the earliest pivot in $N[v]$.
\begin{equation}
    e(v)=\operatorname*{arg\,min}_{p\in M\cap N[v]}\rho(p).
    \label{eq:dynamic-owner}
\end{equation}
Two vertices belong to the same cluster exactly when they have the same
owner. Each bucket stores its unique pivot, or a null value if it
contains none. Inspecting these values and including $v$ itself when
$v\in M$ determines $e(v)$ in $O(K)$ time. We maintain a blocking
certificate $\ell(v)$ satisfying
\begin{equation}
\begin{aligned}
    \ell(v)&=v &&\text{if }v\in M,\\
    \ell(v)&\in M\cap N(v),\qquad
    \rho(\ell(v))<\rho(v) &&\text{if }v\notin M.
\end{aligned}
\label{eq:dynamic-certificate}
\end{equation}
For a nonpivot, the certificate may be any earlier neighboring pivot, whereas the owner is always the earliest one. Thus, a certificate can remain valid even when the owner changes. For each pivot $p$, maintain its follower group
\[
    F(p)=\{v:\ell(v)=p\},
\]
including $p$ itself. When $p$ disappears from the pivot set, we refresh the certificates of its surviving followers. Each vertex stores references to its bucket records and its position in its follower group, so moving it between groups takes constant time.

\paragraph{Updating the candidates.}
The queue $\mathcal H$ stores vertices awaiting repair, ordered by
priority. The procedure \textsc{Refresh}$(v)$ clears the old certificate
and unlinks $v$ from its follower group, if it has one. It inspects the
pivots of $v$'s buckets. If an earlier neighboring pivot exists, it sets
$\ell(v)$ to the earliest such pivot and attaches $v$ to its follower
group. Otherwise, it inserts $v$ into $\mathcal H$. A vertex is rechecked
when extracted because an earlier pivot may have appeared while it
was waiting.

Algorithm~\ref{alg:dynamic-consensus} gives the updates for all trials.
Within each repair, we update owners after each pivot removal or
installation and update the score after each insertion, deletion, or
owner change, using the rules below.

\RestyleAlgo{boxed}
\begin{algorithm}[htbp]
\caption{Fully dynamic consensus clustering}
\label{alg:dynamic-consensus}
\DontPrintSemicolon
\KwIn{Input-label access, $\varepsilon,\zeta$, and update bound $U$.}
\KwOut{The index of the maintained candidate after each update.}

Choose $T,s$ and sample $(Q_a,J_a)$ for all trials and $S$ for evaluation\;
Initialize all candidate and score records as empty\;
\ForEach{update at vertex $v$}{
    \For{$a=1,\ldots,T$}{
        Initialize $\mathcal H\gets\emptyset$ for trial $a$\;
        \eIf{the update inserts $v$}{
            Draw an independent uniform priority $\rho_a(v)$\;
            Add $v$ to its buckets and initialize its owner and score records\;
            Apply \textsc{Refresh}$(v)$\;
        }{
            Remove $v$ from its owner and score records\;
            \eIf{$v\in M$}{
                Detach $F(v)$ and let $L\gets F(v)\setminus\{v\}$\;
                Remove $v$ from $M$ and clear its bucket-pivot entries\;
                Remove $v$ and its bucket memberships\;
                Apply \textsc{Refresh} to every vertex in $L$\;
            }{
                Remove $v$ from its follower group and its buckets\;
            }
        }
        \While{$\mathcal H\neq\emptyset$}{
            Extract the minimum-priority vertex $x$\;
            \eIf{$x$ has an earlier neighboring pivot}{
                Set $\ell(x)$ to the earliest neighboring pivot
                and attach $x$ to its follower group\;
            }{
                Let $R\gets M\cap N(x)$ and $L\gets\bigcup_{p\in R}F(p)$\;
                Remove every $p\in R$ from $M$, clear its bucket-pivot
                entries, and detach $F(p)$\;
                Install $x$ as a pivot with $\ell(x)=x$ and $F(x)=\{x\}$,
                and record it in its buckets\;
                Apply \textsc{Refresh} to every vertex in $L$\;
            }
        }
        Update the input candidate $P_{J_a}|_{V_t}$ and its score at $v$\;
    }
    Store the index of a candidate of minimum score\;
}
\end{algorithm}

\begin{lemma}[Exact maintenance]
\label{lem:dynamic-pivot-correctness}
After each update, the repair in Algorithm~\ref{alg:dynamic-consensus} restores \eqref{eq:dynamic-certificate} and maintains the exact greedy pivot set and its \textsc{Pivot} partition for the retained priorities. Each vertex enters $\mathcal H$ at most once during an update. Newly installed pivots remain pivots for the rest of that update, and every removed pivot was present at its start.
\end{lemma}

\begin{proof}
The set $M$ remains independent because every conflicting pivot is removed before a new pivot is installed. Refreshing the affected followers ensures that, after each repair step, every nonpivot has either a valid earlier certificate or a queue entry.
When $x$ becomes a pivot, each pivot $p$ it displaces and every follower $u$ of $p$ satisfy $\rho(u)\geq\rho(p)>\rho(x)$. Thus, all new queue entries come after $x$, and minimum-priority extraction processes vertices in strictly increasing order. Newly installed pivots therefore remain pivots for the rest of the update. Each displaced pivot stays blocked by its remover and cannot return, so all removed pivots were present at the start.

A queued vertex belongs to no follower group, so it cannot be refreshed again while waiting. Any earlier pivot found when it is processed must have been installed since it was queued and is therefore permanent. Otherwise, the vertex itself becomes a permanent pivot. Hence each vertex enters the queue at most once, and the procedure terminates.
At termination, every nonpivot has an earlier neighboring pivot. Inducting over the priority order, the greedy scan selects every vertex in $M$ and excludes every vertex outside $M$. Thus, $M$ is the exact greedy pivot set, and the owner rule \eqref{eq:dynamic-owner} gives its \textsc{Pivot} partition.
\end{proof}

\paragraph{Update-time analysis.}
We count the work needed for the pivot set and certificates. For update $t$, let $n_t=\max\{|V_{t-1}|,|V_t|\}$. Let $A_t$ count \textsc{Refresh} calls, including the initial call on
insertion. Let $F_t$ count pivot insertions and removals, and let $H_t$
count queue entries.
Vertex records and nonempty bucket records are indexed by balanced dictionaries. Existing vertices retain their bucket and follower references, so refreshing them requires no dictionary lookup. Inserting or deleting a vertex incurs $O(K\log(Kn_t+2))$ dictionary work.
Each refresh costs $O(K)$ outside the queue, and each pivot change updates $O(K)$ bucket-pivot entries. Extracting and rechecking a queued vertex costs $O(K+\log(n_t+2))$. Thus, the work for update $t$ is
\begin{equation}
    O\!\left(
        KA_t+KF_t+KH_t
        +H_t\log(n_t+2)
        +K\log(Kn_t+2)
    \right).
    \label{eq:pivot-update-work}
\end{equation}
We also need to count changes in the actual cluster owners. Let $R_t$ be the number of surviving vertices whose owner changes between the completed states before and after update $t$.

\begin{lemma}[Queue and owner sensitivity]
\label{lem:dynamic-sensitivity}
For each fixed update, conditional on the sampled indices $Q$,
\[
    \mathbb E[H_t]\leq1,
    \qquad
    \mathbb E[F_t]\leq1,
    \qquad
    \mathbb E[R_t]=O(\log(n_t+2)).
\]
\end{lemma}

\begin{proof}
Conditioning on $Q$ fixes the graphs before and after the update, while their priorities induce a uniform random order.
Let $S_t$ be the influence set of Censor-Hillel, Haramaty, and Karnin~\cite[Theorem~1]{censor-hillel-2016}, which satisfies $\mathbb E[|S_t|]\leq1$. Its closure starts at the updated vertex when its pivot status requires a change. It includes an old pivot with an earlier influenced neighbor, and an old nonpivot once all its earlier old-pivot neighbors have entered the closure.

Every displaced pivot in our repair has an earlier newly installed pivot as a neighbor. Every surviving old nonpivot enters the queue only after all its earlier old-pivot neighbors have been removed. Induction over the repair events therefore places all queued vertices and all vertices whose pivot membership changes in $S_t$. By Lemma~\ref{lem:dynamic-pivot-correctness}, each is counted at most once, so
\[
    H_t\leq|S_t|,
    \qquad F_t\leq|S_t|.
\]

The owner $e(v)$ is the greedy eliminator studied by Behnezhad et al.~\cite[Theorem~3]{behnezhad-2019}. Their theorem bounds the expected number of eliminator changes after deleting a fixed vertex by $O(\log(n_t+2))$. Insertion compares the same two graphs in reverse, with a uniform ranking on their union, and obeys the same bound.
\end{proof}
We account for certificate refreshes using owner changes and a potential that counts vertices whose certificate differs from their owner:
\[
    \Phi_t=\bigl|\{v\in V_t:\ell_t(v)\neq e_t(v)\}\bigr|.
\]

\begin{lemma}[Certificate amortization]
\label{lem:certificate-amortization}
For every update $t$ and every fixed choice of distinct priorities,
\begin{equation}
    A_t\leq
    \mathbf{1}[\text{update }t\text{ is an insertion}]
    +(K+1)R_t+\Phi_{t-1}-\Phi_t.
    \label{eq:certificate-amortization}
\end{equation}
\end{lemma}

\begin{proof}
Consider a surviving vertex whose owner remains $p$. The pivot $p$ remains present throughout the repair. No newly installed neighboring pivot can be earlier than $p$, since it would remain present and change the final owner. Thus, $p$ stays the earliest neighboring pivot. If the certificate is already $p$, no refresh occurs. Otherwise, the first refresh chooses $p$, reducing the vertex's contribution to $\Phi$ from one to zero and preventing further refreshes.

Now consider a surviving vertex whose owner changes. Each refresh follows the removal of an old blocking pivot, since newly installed pivots are permanent during the update. There are at most $K$ old pivots in its closed neighborhood, so it is refreshed at most $K$ times. Its contribution to $\Phi$ can increase by at most one. Thus, its refresh count plus its potential increase is at most $K+1$.
An inserted vertex needs one initial refresh and ends with its certificate equal to its owner: it is either blocked immediately or becomes a permanent pivot. A deleted vertex is not refreshed and removes a nonnegative contribution from $\Phi$. Summing these per-vertex bounds proves the inequality.
\end{proof}

\begin{lemma}[Maintaining one \textsc{Pivot} candidate]
\label{thm:single-pivot-maintenance}
Starting empty, the pivot and certificate records of one trial can be maintained over any oblivious sequence of $1\leq u\leq U$ vertex updates in expected total time
\[
    O\!\left(K^2\sum_{t=1}^{u}\log(n_t+2)\right).
\]
\end{lemma}

\begin{proof}
Since the process starts empty, $\Phi_0=0$. Summing \eqref{eq:certificate-amortization} and using $\Phi_u\geq0$ gives
\[
    \sum_{t=1}^{u}A_t
    \leq u+(K+1)\sum_{t=1}^{u}R_t.
\]
Taking expectations and applying Lemma~\ref{lem:dynamic-sensitivity},
\[
    \mathbb E\!\left[\sum_{t=1}^{u}A_t\right]
    =O\!\left(K\sum_{t=1}^{u}\log(n_t+2)\right).
\]
Substituting this bound, together with $\mathbb E[F_t]\leq1$ and $\mathbb E[H_t]\leq1$, into \eqref{eq:pivot-update-work} gives the claimed total time. The dictionary work also fits this bound because
$\log(Kn_t+2)=O(K\log(n_t+2))$ for $K\geq1$.
The bucket memberships, follower entries, and queue records use
$O(K(n+1))$ space with $n$ active vertices.
\end{proof}
For the sampled majority graph, $K=35$. If $N$ bounds the active set size, this gives $O(\log(N+2))$ expected amortized update time for the pivot and certificate records of one trial.

\paragraph{Maintaining candidate scores.}

A vertex can change owner while its blocking certificate remains valid. To maintain the estimated costs, we therefore store the actual owners and report their changes separately.

For each pivot $p$, maintain its owner list
\[
    O(p)=\{v:e(v)=p\}.
\]
Each bucket also maintains a balanced tree of its active vertices, ordered first by their current owner's priority and then by vertex identifier. If $M\cap N[v]=\emptyset$, temporarily assign $v$ a unique singleton label and owner priority $+\infty$. These temporary labels are distinct from pivot labels, so the intermediate assignments still define a partition.

\emph{Removing a pivot.} After removing $p$, visit $O(p)$ and recompute each surviving member's owner from its bucket-pivot values. Update its owner-list position and its entries in all its bucket trees.

\emph{Installing a pivot.} First assign $p$ to itself. Then search its bucket trees for vertices whose owner priority exceeds $\rho(p)$. Each such vertex now has owner $p$. Reassign it and update all its tree entries before continuing, which prevents duplicate reports from overlapping buckets.

An arriving vertex is initially assigned its earliest current neighboring pivot, or a temporary singleton if none exists. A departing vertex is removed from these indexes before surviving owners are repaired. Perform the owner updates after each individual pivot removal and installation in Algorithm~\ref{alg:dynamic-consensus}; certificate refreshes still follow the order specified there.
Each reported change costs $O(K\log(n_t+2))$ time, and the initial searches for each installed pivot have the same bound.
We now relate the intermediate owner changes to $R_t$, which counts only changes between the completed states before and after the update. Let $B_t$ count all reported owner changes during the repair, including creation or removal of the updated vertex's owner record.

\begin{lemma}
\label{lem:owner-reporting}
For every update $t$,
\[
    B_t=O(K(R_t+1)).
\]
Consequently, for $K=35$,
\[
    \mathbb E[B_t]=O(\log(n_t+2)).
\]
\end{lemma}

\begin{proof}
Suppose a surviving vertex has the same owner $p$ before and after the update. The pivot $p$ remains present throughout the repair. No earlier neighboring pivot can be newly installed, since it would remain present and change the final owner. Thus, the owner stays $p$ throughout, and the vertex is never reported.

For a vertex whose owner changes, each report is caused by removing its current owner or installing a new earlier neighboring pivot. Removed pivots belong to the old pivot set, and newly installed pivots remain until the repair ends. Both sets are independent, so each contributes at most $K$ pivots to the vertex's closed neighborhood. Hence, the vertex is reported at most $2K$ times. Temporary singleton assignments are already counted among these changes.
The inserted or deleted vertex contributes $O(K)$ additional record changes. Summing gives $B_t=O(K(R_t+1))$. Taking expectations and applying Lemma~\ref{lem:dynamic-sensitivity} proves the second bound.
\end{proof}

\paragraph{Updating the scores.}
We use the Mirkin-distance identity to express each estimated cost through block sizes and intersections. For a candidate $C$ with block labels $c_C(v)$, maintain
\[
\begin{aligned}
    a_p&=|\{v\in V_t:c_C(v)=p\}|,\\
    b_{\ell,p,b}
    &=|\{v\in V_t:c_C(v)=p,\ c_{j_\ell}(v)=b\}|.
\end{aligned}
\]
Define the score
\begin{equation}
    \sigma_t(C)
    =s\sum_p\binom{a_p}{2}
     -2\sum_{\ell=1}^{s}\sum_{p,b}\binom{b_{\ell,p,b}}{2}.
    \label{eq:dynamic-score}
\end{equation}
By the Mirkin-distance identity, the difference
$\widehat{\operatorname{cost}}_{S,t}(C)-\sigma_t(C)/s$ depends only on
$S$ and $V_t$. Thus a minimum-score candidate also minimizes the
estimated cost.
These counters can be updated locally whenever a vertex enters or leaves a block. Since $\binom{q+1}{2}-\binom q2=q$, inserting $v$ into block $p$ changes its score by
\[
    \Delta_+\sigma
    =sa_p-2\sum_{\ell=1}^{s}b_{\ell,p,c_{j_\ell}(v)},
\]
using the counts before insertion. Similarly, deleting $v$ from block $p$ changes the score by
\[
    \Delta_-\sigma
    =-s(a_p-1)
     +2\sum_{\ell=1}^{s}\bigl(b_{\ell,p,c_{j_\ell}(v)}-1\bigr),
\]
using the counts before deletion. Update the counters after applying the score change. Moving a vertex between blocks consists of a deletion from its old block followed by an insertion into its new block.

Store the block sizes in one balanced dictionary and the nonzero intersection counts in a separate dictionary for each evaluation input. Each dictionary has at most $n$ entries, where $n$ is the current number of vertices. Cache the evaluation labels of each vertex upon arrival and discard empty count records. Each insertion, deletion, or move therefore takes $O(s\log(n_t+2))$ time, with $O((n+1)s)$ space.
For a Pivot candidate, apply these updates whenever its owner changes, including assignments to temporary singletons. For a sampled-input candidate, only the inserted or deleted vertex changes the counters, since surviving input labels remain fixed.

\begin{lemma}[Maintaining a candidate and its score]
\label{lem:candidate-score-maintenance}
For $K=35$ and a fixed evaluation sample of size $s\geq1$, one Pivot candidate and its score can be maintained over $u$ updates in expected total time
\[
    O\!\left(s\sum_{t=1}^{u}\log^2(n_t+2)\right).
\]
With $n$ active vertices, the space usage is $O((n+1)s)$.

One sampled-input candidate and its score require $O(s\log(n_t+2))$ worst-case time per update, with the same space bound.
\end{lemma}

\begin{proof}
Updating the owner lists and bucket trees takes $O((B_t+F_t+1)\log(n_t+2))$ time: each owner report updates its index entries, and each pivot change initiates bucket searches. Updating the score counters takes $O(s(B_t+1)\log(n_t+2))$ time. Thus, the additional work for update $t$ is
\[
    O\!\left(
        (B_t+F_t+1)\log(n_t+2)
        +s(B_t+1)\log(n_t+2)
    \right).
\]
By Lemmas~\ref{lem:owner-reporting} and \ref{lem:dynamic-sensitivity}, $\mathbb E[B_t]=O(\log(n_t+2))$ and $\mathbb E[F_t]\leq1$. The expected additional work is therefore $O(s\log^2(n_t+2))$. Summing and adding the pivot-and-certificate work from Lemma~\ref{thm:single-pivot-maintenance} proves the total-time bound.

The owner lists and bucket trees use $O(n)$ space. The intersection counters and cached evaluation labels use $O((n+1)s)$ space.
For a sampled-input candidate, surviving labels remain fixed, so each update changes only the inserted or deleted vertex's counters.
\end{proof}

Theorem~\ref{thm:implicit-dynamic} follows from the offline approximation
argument and the bounds on maintaining the candidates and their scores.

\begin{proof}[Proof of Theorem~\ref{thm:implicit-dynamic}]
Fix an update time $t$. By exact maintenance and the distribution of
the active priorities, the candidates in $\mathcal Y_t$ are distributed
as independent offline trials on $\mathcal P_t$ and are independent
of $S$. Apply the repetition and evaluation argument in the proof of
Theorem~\ref{thm:offline} with accuracy $\varepsilon$ and failure
probability $\zeta/U$. The choices of $T$ and $s$ suffice because
$\log(4TU/\zeta)=O(\Lambda)$. A union bound over at most $U$ updates
gives
\[
    \operatorname{cost}_t(C_t)
        \leq (1.6944+\varepsilon)\,\mathrm{OPT}_t
    \qquad\text{simultaneously after all updates}
\]
with probability at least $1-\zeta$.

For a nonempty sequence of $u\leq U$ updates,
Lemma~\ref{lem:candidate-score-maintenance} bounds the expected total
maintenance time of all candidates and their scores by
\[
    O\!\left(Ts\sum_{t=1}^{u}\log^2(n_t+2)\right)
    =O\!\left(u\varepsilon^{-3}\Lambda^2\log^2(N+2)\right).
\]
This also covers initialization in $O(Ts)$ time, reading and caching
the sampled labels, and choosing a minimum score in $O(T)$ time per
update. Dividing by $u$ gives the expected amortized update time.

For a membership query, the stored candidate index identifies the
records to query. We locate the vertex in its balanced dictionary and
read its stored owner or input label in $O(\log(n+2))$ worst-case time.
A same-cluster query compares two such labels. Traversing the selected
candidate's records lists all memberships in $O(n+1)$ time.
\end{proof}
\ifusebiblatex
  \printbibliography
\else
  \bibliographystyle{alpha}
  \bibliography{ref}
\fi

\newpage
% add appendix
\appendix
\section{Verification of the Offline Charging Inequalities}
\label{app:offline-charging}

\begin{proof}[Proof of Lemma~\ref{lem:charging-inequalities}]
Throughout the proof, take $\beta,\gamma$ as in
\eqref{eq:approx-constants}. Expanding the definition of
$\psi_7$ gives
\[
    \psi_7(x)=70x^3-238x^4+308x^5-181x^6+40x^7.
\]
We verify the two charging inequalities using the Bernstein
basis on $[0,1]$:
\[
    B_{j,d}(y)=\binom{d}{j}y^j(1-y)^{d-j},
    \qquad j=0,\ldots,d.
\]
These polynomials are nonnegative and sum to one. A polynomial
expressed in this basis therefore lies between its smallest
and largest coefficients.

\paragraph{The pair inequality.}
Write $\psi_7(x)=xH(x)$, where
\[
    H(x)=70x^2-238x^3+308x^4-181x^5+40x^6.
\]
The degree-six Bernstein coefficients of $H(y/2)$ are
\[
    \left(0,0,\frac76,\frac{161}{80},
        \frac73,\frac{145}{64},\frac{63}{32}\right).
\]
Their maximum is $7/3$, so $H(x)\leq7/3$ for
$0\leq x\leq1/2$. Since $\gamma=2-\beta$,
\begin{align*}
    &\gamma x-\beta[x+\psi_7(x)x]-(1-\beta)2x(1-x)\\
    &\qquad=x^2[2(1-\beta)-\beta H(x)]\\
    &\qquad\geq x^2\left[2(1-\beta)-\frac73\beta\right]
        =\frac{13578}{20101}x^2\geq0.
\end{align*}
This proves \eqref{eq:joint-pair}.

\paragraph{The bad-triplet inequality.}
Let $t=\{u,v,w\}$ be bad in $G^\star$, with edges $uw,vw$
and nonedge $uv$. Set
\[
    m_1=m_{uw},\qquad m_2=m_{vw},\qquad m_3=m_{uv}.
\]
In each input partition, separating $u$ from $v$ requires
separating at least one from $w$. Averaging gives
$w_{uv}^-\leq w_{uw}^-+w_{vw}^-$, so
\[
    m_1+m_2+m_3
        =w_{uw}^-+w_{vw}^-+1-w_{uv}^-\geq1.
\]
Thus,
\begin{equation}
    0\leq m_i\leq\frac12,
    \qquad m_1+m_2+m_3\geq1.
    \label{eq:triplet-domain}
\end{equation}
The restriction of $C^\star$ to $t$ has one of five forms:
\[
\begin{array}{c|l}
    \text{Clusters} & \text{Cost}\\
    \hline
    \{u,v,w\} & C_1=m_1+m_2+(1-m_3)\\
    \{u,w\},\{v\} & C_2=m_1+(1-m_2)+m_3\\
    \{v,w\},\{u\} & C_3=(1-m_1)+m_2+m_3\\
    \{u,v\},\{w\} & C_4=(1-m_1)+(1-m_2)+(1-m_3)\\
    \{u\},\{v\},\{w\} & C_5=(1-m_1)+(1-m_2)+m_3.
\end{array}
\]
Define
\[
    f(x)=\beta[1-x+\psi_7(x)x]+(1-\beta)2x(1-x).
\]
It suffices to prove
\[
    f(m_1)+f(m_2)+f(m_3)\leq\gamma C_j,
    \qquad j=1,\ldots,5,
\]
throughout \eqref{eq:triplet-domain}.

Set $x_i=1-2m_i$ and $h(x)=f((1-x)/2)$. The domain becomes
\begin{equation}
    x_i\geq0,\qquad x_1+x_2+x_3\leq1,
    \label{eq:triplet-simplex}
\end{equation}
and
\begin{align}
    40202h(x)
        &=26149-480x-29413x^2+16800x^3+10080x^4\notag\\
        &\quad-6048x^5-8736x^6+2016x^7+1920x^8.
    \label{eq:h-polynomial}
\end{align}
The degree-seven Bernstein coefficients of
$40202(\gamma/2-h'(x))$ are
\[
    \left(34538,\frac{300592}{7},\frac{342618}{7},
        \frac{359780}{7},\frac{350062}{7},
        \frac{334968}{7},\frac{362882}{7},55828\right),
\]
and those of $40202(\gamma/2+h'(x))$ are
\[
    \left(33578,\frac{176220}{7},\frac{134194}{7},
        \frac{117032}{7},\frac{126750}{7},
        \frac{141844}{7},\frac{113930}{7},12288\right).
\]
All coefficients are positive. Therefore,
\begin{equation}
    -\frac\gamma2<h'(x)<\frac\gamma2,
    \qquad 0\leq x\leq1.
    \label{eq:h-derivative-bounds}
\end{equation}

For each case, write
$D_j=\gamma C_j-h(x_1)-h(x_2)-h(x_3)$.
For $C_1$, we have
\[
    C_1=\frac32-\frac{x_1+x_2-x_3}{2}.
\]
By \eqref{eq:h-derivative-bounds},
\[
    \frac{\partial D_1}{\partial x_i}
        =-\frac\gamma2-h'(x_i)<0\quad(i=1,2),
    \qquad
    \frac{\partial D_1}{\partial x_3}
        =\frac\gamma2-h'(x_3)>0.
\]
The minimum over \eqref{eq:triplet-simplex} therefore occurs
on $x_3=0$, $x_1+x_2=1$. On this boundary,
\[
    40202D_1(x,1-x,0)=(2x-1)^2R(x),
\]
where
\[
    R(x)=3530-7170x+570x^2+14160x^3
        -9480x^4+2880x^5-960x^6.
\]
Its degree-six Bernstein coefficients are
\[
    (3530,2335,1178,767,1178,2335,3530).
\]
They are positive, so $R(x)>0$ on $[0,1]$ and hence $D_1\geq0$.
The cases $C_2,C_3$ follow by permuting coordinates.

For $C_4$,
\[
    C_4=\frac32+\frac{x_1+x_2+x_3}{2},
    \qquad
    \frac{\partial D_4}{\partial x_i}
        =\frac\gamma2-h'(x_i)>0.
\]
Its minimum is at $(0,0,0)$, where
\[
    D_4(0,0,0)=\frac32\gamma-3h(0)
        =\frac{23727}{40202}>0.
\]
Finally, for $C_5$,
\[
    C_5=\frac32+\frac{x_1+x_2-x_3}{2}.
\]
The gap increases with $x_1,x_2$ and decreases with $x_3$,
again by \eqref{eq:h-derivative-bounds}. Its minimum is
therefore at $(0,0,1)$, where
\[
    D_5(0,0,1)=\gamma-2h(0)-h(1)
        =\frac{1765}{20101}>0.
\]
This proves \eqref{eq:joint-triplet} in all five cases.
\end{proof}

\makeatletter
\newcommand{\setword}[2]{%
  \phantomsection
  #1\def\@currentlabel{\unexpanded{#1}}\label{#2}%
}
\makeatother

\newcommand{\boolNo}{\text{No}} % boolean No
\newcommand{\boolYes}{\text{Yes}} % boolean Yes

\NewDocumentCommand{\localG}{g}{%
    \IfNoValueTF{#1}
    {G}%
    {G_{#1}}%
}
\NewDocumentCommand{\sparseG}{g}{%
    \IfNoValueTF{#1}
    {\widetilde{G}}%
    {\widetilde{G}_{#1}}%
} % sparsified graph in dynamic agreement (DA)
\NewDocumentCommand{\dyCC}{g}{%
    \IfNoValueTF{#1}
    {\textsc{CC}}%
    {\textsc{CC}_{#1}}%
} % dynamic data structure for connected components
\NewDocumentCommand{\journal}{g}{%
    \IfNoValueTF{#1}
    {\mathcal{J}}%
    {\mathcal{J}_{#1}}%
}
\newcommand{\epsDA}{\varepsilon_{\rm DA}} % accuracy parameter for dynamic agreement
\NewDocumentCommand{\interestSet}{g}{%
    \IfNoValueTF{#1}
    {\mathcal{I}}%
    {\mathcal{I}_{#1}}%
} % interesting set in DA
\NewDocumentCommand{\anchorSet}{g}{%
    \IfNoValueTF{#1}
    {\Phi}%
    {\Phi_{#1}}%
} % anchor set in DA


\newcommand{\openN}[2]{N_{#1}(#2)} % open neighborhood of vertex #2 in graph #1
\newcommand{\closedN}[2]{N_{#1}[#2]} % open neighborhood of vertex #2 in graph #1
\newcommand{\snapshotN}[1]{\widetilde{N}(#1)} % open neighborhood of vertex #2 in graph #1


\newcommand{\explicitDA}{\mathrm{ExplicitDynamicAgreement}} % explicit dynamic agreement algorithm
\newcommand{\ADA}{\mathrm{AmplifiedDynamicAgreement}} % explicit dynamic agreement algorithm
\newcommand{\DA}{\mathrm{DynamicAgreement}} % explicit dynamic agreement algorithm
\NewDocumentCommand{\Notify}{g}{%
    \IfNoValueTF{#1}
    {\mathrm{Notify}}%
    {\mathrm{Notify}(#1)}%
}
\NewDocumentCommand{\Anchor}{g}{%
    \IfNoValueTF{#1}
    {\mathrm{Anchor}}%
    {\mathrm{Anchor}(#1)}%
}
\NewDocumentCommand{\Clean}{g}{%
    \IfNoValueTF{#1}
    {\mathrm{Clean}}%
    {\mathrm{Clean}(#1)}%
}
\NewDocumentCommand{\Connect}{g}{%
    \IfNoValueTF{#1}
    {\mathrm{Connect}}%
    {\mathrm{Connect}(#1)}%
}
\NewDocumentCommand{\LabelSync}{g}{%
    \IfNoValueTF{#1}
    {\mathrm{LabelSync}}%
    {\mathrm{LabelSync}(#1)}%
}
\NewDocumentCommand{\ProbAgree}{g}{%
    \IfNoValueTF{#1}
    {\mathrm{ProbablisticAgreement}}%
    {\mathrm{ProbablisticAgreement}(#1)}%
}
\NewDocumentCommand{\Heavy}{g}{%
    \IfNoValueTF{#1}
    {\mathrm{Heavy}}%
    {\mathrm{Heavy}(#1)}%
}

\newcommand{\delt}[1]{\mathrm{Del}_{#1}}
\newcommand{\addt}[1]{\mathrm{Add}_{#1}}
\newcommand{\delIncEdge}[1]{\delta(#1)} % number of edges incident on vertex #1 that are deleted (\in \delt)
\NewDocumentCommand{\ccclass}{g}{%
    \IfNoValueTF{#1}
    {C}%
    {C_{#1}}%
} % connected component class
\NewDocumentCommand{\ccclasses}{g}{%
    \IfNoValueTF{#1}
    {\mathcal C}%
    {\mathcal{C}_{#1}}%
} % connected component class
\newcommand{\comp}[1]{\mathrm{comp}(#1)} % connected component of vertex #1
\NewDocumentCommand{\piece}{ g }{%
    \IfNoValueTF{#1}
    {piece}%
    {piece#1}%
}
\newcommand{\retainedPiece}[1]{R_{#1}} % retained pieces of class #1
\newcommand{\piecesToMerge}{\mathcal{Z}} %
\newcommand{\pieceToMerge}{P} %
\newcommand{\complabel}{\ell} %

\newcommand{\successDA}{\mathcal{E}_{\rm DA}} % success event of dynamic agreement
\newcommand{\jnTime}[1]{\card{\journal(#1)}} % length of the journal at time #1
\newcommand{\nChangedVers}[1]{c_{#1}} % number of vertices whose labels changed at time #1

\newcommand{\edSolver}{\mathcal D}
\newcommand{\edFactor}{\eta} % approximation factor of the explicit correlation-clustering algorithm
\newcommand{\edTime}[1]{T_{\mathcal D}(#1)}
\newcommand{\edcost}[1]{\widehat{\cost}_{#1}}
\newcommand{\edmean}[2]{\widehat{\cost}^{(#1)}_{#2}}
\newcommand{\edscore}[2]{\widehat f_{#1}(#2)}
\newcommand{\edhistory}[2]{\widehat G_{#1}(#2)}
\newcommand{\edminimum}[1]{\widehat M_{#1}}
\newcommand{\edproxy}[2]{Z_{#1}(#2)}
\newcommand{\edlog}[2]{\mathcal L_{#1}(#2)}
\newcommand{\edbucket}[2]{B_{#1}(#2)}
\newcommand{\edatom}[1]{B_{123}(#1)}

% name for variable used in the implementation of LabelSync
\newcommand{\setClasses}{\mathit{Set}_{\ccclasses}} % variable used in the
\newcommand{\groupdict}{\mathit{group\_dict}} % variable used in the
\newcommand{\compdict}{\mathit{comp\_dict}} % variable used in the
\newcommand{\changedList}{\mathit{changed\_list}} % variable used in the


\section{Explicit Dynamic Algorithm for Insertion-Only Streams}\label{sec:explicit_dynamic_algorithm_for_insertion_only_stream}

In this section, we consider the explicit dynamic model for insertion-only streams.
In this model, at each update time $ t $, a new vertex $ v_{t} $ arrives.
The algorithm needs to maintain the solution and reports $ v_t $ and its cluster label, and for each vertex whose cluster identifier has changed, reports the vertex and its new cluster identifier.

Our results is an explicit dynamic algorithm for consensus clustering in the insertion-only stream model, which achieves a strictly better than $ 2 $ approximation ratio with probablity at least $ 1-\zeta $, and polylogarithmic amortized update time in expectation.
\ExplicitDynamic*

We use the vertex-arrival model of \cref{subsec:online-model}, with the input-label access of \cref{subsec:dynamic-model}.
In particular, the input partitions and arrival order are oblivious, the horizon $ n $ is known, and any selected input label of an arriving vertex can be read in constant time.

We obtain \cref{thm:explicit.dynamic} from the following two lemmas.

\begin{lemma}\label{lem:explicit.consensus.clustering}
Suppose that $ \edSolver $ is a randomized explicit dynamic correlation-clustering algorithm in the database model~(\cref{def:query.model.label}) that, on every oblivious stream of at most $ n $ vertex insertions, has expected amortized update time $ \edTime{n}=\polylog(n+2) $ including queries and reports, and with probability at least $ 1/2 $ maintains an $ \edFactor $-approximation simultaneously at all prefixes.
Then, for every $ 0<\eps<2/(3\edFactor+5) $ and $ \zeta\in(0,1) $, there is an explicit dynamic consensus-clustering algorithm for insertion-only streams that achieves approximation ratio $ 2-2/(3\edFactor+5)+\eps $ at all prefixes simultaneously with probability at least $ 1-\zeta $.
Its expected amortized update time is
\begin{align}
    O\!\left(
        \eps^{-3}\log^2\frac{2n}{\zeta}
        \log\frac{2n}{\eps\zeta}\cdot\polylog(n+2)
    \right).
    \label{eq:explicit.consensus.time}
\end{align}
\end{lemma}

\begin{remark}
Given the upper bound $ \edTime{n} $, we can convert the algorithm of \cref{lem:explicit.consensus.clustering} into one whose amortized update time satisfies \eqref{eq:explicit.consensus.time} on every execution, preserving the approximation guarantee with probability at least $ 1-\zeta $.
\end{remark}

\begin{definition}[Database model~\cite{assadi-wang}]\label{def:query.model.label}
At time $ t $, the algorithm accesses the graph of positive edges on $ V_t $ through three queries.
A degree query returns the number of neighbors of a given vertex, an edge query tests whether two given vertices are adjacent, and a neighborhood-sample query returns an independent uniformly random neighbor of a given vertex of positive degree.
Each query costs $ O(\log(|V_t|+2)) $ time.
\end{definition}

\paragraph{From clique buckets to the database model.}\label{par:explicit.bucket.queries}
By the bucket representation of \cref{subsec:offline-buckets,subsec:dynamic-algorithm}, two distinct vertices of a three-input majority graph are adjacent exactly when their input labels agree on some pair of coordinates.
For a vertex $ v $, the intersection of any two of its three pair buckets is the set of vertices agreeing with $ v $ on all three coordinates.
Order each pair bucket by the remaining coordinate and then by vertex index in an order-statistic tree, so this intersection occupies one rank interval.
The open neighborhood is the disjoint union of the common intersection with $ v $ removed and the complement of the intersection in each of the three buckets.
Summing these four sizes gives the degree, and a uniform rank in their concatenation gives a uniform neighbor.
Thus insertions and all three database queries take $ O(\log(|V_t|+2)) $ time, using three bucket memberships per vertex.

\begin{restatable}{lemma}{ExplicitDA}\label{lem:explicit.correlation.clustering}
For some absolute constant $ \edFactor\ge1 $, there is an explicit randomized dynamic correlation clustering algorithm in the database model~(\cref{def:query.model.label}) achieves an $ \edFactor $-approximation simultaneously at all prefixes of any fixed insertion-only stream with probability at least $ 1-n^{-2} $.
%The algorithm uses $ N $ as its size parameter and takes the majority of three independent runs of each agreement or heaviness test.
For every $ t\le n $, the expected total work through insertion $ t $, including all reported label assignments, is $ O(t\polylog(n+2)) $.
\end{restatable}

\begin{proof}[Proof of~\cref{thm:explicit.dynamic}]
Apply \cref{lem:explicit.consensus.clustering} with the algorithm of \cref{lem:explicit.correlation.clustering}.
For fixed $ \eps $ and inverse-polynomial $ \zeta $, \eqref{eq:explicit.consensus.time} is polylogarithmic in $ n $.
\end{proof}

We defer the proof of \cref{lem:explicit.correlation.clustering} to \cref{sec:explicit_dynamic_correlation_clustering_algorithm_for_insertion_only_stream}.
The remainder of this section is devoted to prove \cref{lem:explicit.consensus.clustering}.
The cases $ n=1 $ and $ m\le2 $ are handled by following an input, as in \cref{sec:online}, so assume $ n\ge2 $ and $ m\ge3 $.

\subsection{Algorithm}\label{subsec:explicit.consensus.algorithm}

We follow the history-score selection rule in \cref{sec:oncon-algorithm}, with the following two modifications.
For every triple candidate, maintain its solution by employing $ \edSolver $ instead of $ \Agree $.
In \cref{sec:oncon-algorithm}, the algorithm decides whether to retain the selected candidate or switch to another candidate using history scores computed from exact costs.
As maintaining the exact cost of each candidate is too expensive in dynamic model, we maintain an estimated cost of each candidate using a set of evaluation samples. 
We detail this estimation below.

We reuse the parameter choices from \cref{sec:oncon-algorithm,subsec:oncon-analysis}:
\begin{equation*}
  c=\eta+1,\qquad
  w=\frac{3c}{3c+2},\qquad
  R_0=2-\frac{2}{3\eta+5},\qquad
  R=R_0+\frac{\eps}{2},
\end{equation*}
with $ \eta $ the approximation factor of $ \edSolver $. 
The candidate set is defined as in \cref{sec:oncon-algorithm}, with the parameters changed by a constant factor as follows.

\begin{align}
    \theta=\gamma=\frac{\eps}{40}, \qquad
    K=\left\lceil\frac{8}{\eps}\log\frac{4n}{\zeta}\right\rceil \label{eq:explicit.dynamic.consensus.K}
\end{align}

Sample $ K $ candidates independently as in \cref{sec:oncon-algorithm}: with probability $ w $, follow a uniform input, and otherwise run $ \edSolver $ on the majority graph of a uniform triple of distinct input indices.
Let $ \A=[K] $, and let $ \A_{\rm in} $ and $ \A_{\rm tri} $ be the sets of indices of input and triple candidates, respectively. As in \cref{sec:oncon-algorithm}, write $ A_a(t) $ for candidate $ a $'s partition and $ \lambda_t^a(v) $ for its local identifier of vertex $ v $.

\paragraph{Evaluation samples.}
Set
\begin{align}
    L=\left\lceil\log_2\frac{4n}{\zeta}\right\rceil, \qquad
    g=\left\lceil\frac{2}{\gamma^2}\log\frac{8KnL}{\zeta}\right\rceil \qquad
    s=Lg.
    \label{eq:explicit.consensus.parameters}
\end{align}

Independently of the candidates, we sample $ L $ groups $ (J_{j,1},\ldots,J_{j,g}) $, with all $ J_{j,b} $ independent and uniform in $ [m] $ as evaluation sets use throughout the course of the algorithm.
For a partition $ C $ of $ V_t $, define
\begin{align}
    \edmean{j}{t}(C)&:=\frac1g\sum_{b=1}^g d(C,P_{J_{j,b}}^t),
    &\edcost{t}(C)&:=\min_{j\in[L]}\edmean{j}{t}(C),\label{eq:explicit.consensus.estimates}\\
    \edscore{a}{t}&:=\edcost{t}(A_a(t)),
    &\edhistory{a}{t}&:=\max_{u\le t}\edscore{a}{u},
    &\edminimum{t}&:=\min_{a\in\A}\edhistory{a}{t}.\nonumber
\end{align}
Initially, $ \edhistory{a}{0}=0 $.

The updating rule of the algorithm are given in \cref{alg:explicit.consensus}.

\tienlong{Unreadable beyond this point. The details are correct though.}

\begin{algorithm}[htbp]
    \DontPrintSemicolon
    \KwIn{The arriving vertex $ v_t $, the sampled indices, and the maintained candidate and selection states.}
    \KwOut{The identifier of $ v_t $ and every changed output identifier.}
    Read and store the labels of $ v_t $ at the sampled input indices\;
    \ForEach{$ a\in\A $}{
        Advance candidate $ a $ to time $ t $; for a triple candidate, first update its majority-graph representation\;
        Let $ \edlog{a}{t} $ contain the assignment to $ v_t $ and the reported changes to old vertex labels\;
        Update all group costs by processing $ \edlog{a}{t} $\;
        Compute $ \edscore{a}{t} $ and set $ \edhistory{a}{t}\gets\max\{\edhistory{a}{t-1},\edscore{a}{t}\} $\;
    }
    Set $ \edminimum{t}\gets\min_{a\in\A}\edhistory{a}{t} $ and apply the selection rule in \cref{sec:oncon-algorithm}, using estimated history scores, to choose $ a_t $\;
    \eIf{$ t=1 $ or $ a_t\ne a_{t-1} $}{
        Assign and report $ (a_t,\lambda_t^{a_t}(v)) $ for every $ v\in V_t $\;
    }{
        Apply and report the assignments in $ \edlog{a_t}{t} $, prefixing each local label by $ a_t $\;
    }
    \caption{Explicit dynamic consensus-clustering algorithm.}
    \label{alg:explicit.consensus}
\end{algorithm}

\subsection{Approximation Guarantee}\label{subsec:explicit.consensus.approximation}

For the analysis, define
\begin{align}
    \edproxy{a}{t}:=
    \begin{cases}
        \cost_t(P_{i_a}^t),&a\in\A_{\rm in},\\
        \OPT_t+c|B_{Q_a}^t(F_t^*)|,&a\in\A_{\rm tri}.
    \end{cases}
    \label{eq:explicit.consensus.proxy}
\end{align}
Call $ a $ good for $ t $ if $ \edproxy{a}{t}\le R\OPT_t $ and, when $ a $ is a triple candidate, its solver achieves the $ \edFactor $-approximation at every prefix.

\begin{lemma}[Existence of a good candidate]\label{lem:explicit.good.candidate}
With probability at least $ 1-\zeta/4 $, every prefix $ t\le n $ has a good candidate.
\end{lemma}
\begin{proof}
The proxy depends only on the sampled candidate type and input indices.
The input expectation in \eqref{eq:oncon-input-history}, the majority-error bound in \eqref{eq:oncon-majority-error}, and the cancellation in the proof of \cref{lem:oncon-expectation} give $ \E[\edproxy{a}{t}]\le R_0\OPT_t $.
The fixed-prefix Markov argument in the proof of \cref{lem:oncon-good-history} therefore gives
\begin{align}
    \Prb[\edproxy{a}{t}\le R\OPT_t]\ge\frac{\eps}{4}.
    \label{eq:explicit.consensus.proxy.success}
\end{align}
If $ \OPT_t=0 $, the proxy is zero for every draw, so this bound still holds.

Conditional on any fixed triple, its solver succeeds throughout the stream with probability at least $ 1/2 $.
Consequently, a candidate is good for $ t $ with probability at least $ \eps/8 $, independently across candidates.
By the choice of $ K $ in \eqref{eq:explicit.dynamic.consensus.K}, the probability that any prefix has no good candidate is at most
\begin{align}
    n\exp(-\eps K/8)\le\zeta/4.
    \label{eq:explicit.consensus.candidate.event}
\end{align}
\end{proof}

\begin{lemma}[Accuracy of the estimated costs]\label{lem:explicit.estimated.costs}
With probability at least $ 1-\zeta/2 $, simultaneously for every $ t\le n $ and $ a\in\A $,
\begin{align}
    \edcost{t}(F_t^*)&\le2\OPT_t,\label{eq:explicit.consensus.ceiling}\\
    \left|
        \edcost{t}(A_a(t))-\edcost{t}(F_t^*)
        -\bigl(\cost_t(A_a(t))-\OPT_t\bigr)
    \right|
    &\le\gamma d(A_a(t),F_t^*).
    \label{eq:explicit.consensus.comparison}
\end{align}
\end{lemma}
\begin{proof}
To prove \eqref{eq:explicit.consensus.ceiling}, apply Markov's inequality to each group mean for the fixed partition $ F_t^* $.
Each mean exceeds $ 2\OPT_t $ with probability at most $ 1/2 $, so their minimum does so with probability at most $ 2^{-L} $.
When $ \OPT_t=0 $, all these means are zero.
A union bound over $ t $ gives failure probability at most $ \zeta/4 $.

For \eqref{eq:explicit.consensus.comparison}, condition on all candidate samples and solver choices, leaving the evaluation samples independent.
For each $ a,t,j $, the difference of the two group means averages independent variables $ d(A_a(t),P_{J_{j,b}}^t)-d(F_t^*,P_{J_{j,b}}^t) $.
By the triangle-inequality argument, each variable lies in $ [-d(A_a(t),F_t^*),d(A_a(t),F_t^*)] $.
Hoeffding's inequality bounds the probability of a deviation larger than $ \gamma d(A_a(t),F_t^*) $ in either direction by $ 2\exp(-g\gamma^2/2) $.
If the distance is zero, the comparison is exact.
A union bound over $ KnL $ comparisons gives failure probability at most $ \zeta/4 $.
Finally, if the differences of two lists of numbers all lie in the same interval, the difference of their minima lies in that interval as well.
Taking the minimum over the groups proves \eqref{eq:explicit.consensus.comparison}.

\end{proof}

\begin{lemma}[Estimated history of a good candidate]\label{lem:explicit.estimated.history}
If $ b $ is good for $ t $ and \eqref{eq:explicit.consensus.comparison} holds at $ t $, then
\begin{align}
    \edminimum{t}\le\edhistory{b}{t}
    \le\edcost{t}(F_t^*)+(R-1+3\gamma)\OPT_t.
    \label{eq:explicit.consensus.good.history}
\end{align}
\end{lemma}
\begin{proof}
The first inequality follows from the definition of $ \edminimum{t} $.
If $ b $ is an input candidate, each group cost is nondecreasing under insertions, so $ \edhistory{b}{t}=\edcost{t}(P_{i_b}^t) $ by the input-cost monotonicity argument in the proof of \cref{lem:oncon-expectation}.
Since $ \cost_t(P_{i_b}^t)\le R\OPT_t $, the triangle inequality gives $ d(P_{i_b}^t,F_t^*)\le(R+1)\OPT_t\le3\OPT_t $.
Substituting these bounds in \eqref{eq:explicit.consensus.comparison} proves the second inequality for an input candidate.

For a successful triple candidate, the derivation of \eqref{eq:oncon-history-distance} uses only the solver's simultaneous approximation guarantee, so it gives $ d(A_b(u),F_t^*|_{V_u})\le c|B_{Q_b}^t(F_t^*)| $ for every $ u\le t $.
Each group mean, and hence their minimum, satisfies both the cost comparison and restriction inequalities used in \cref{lem:oncon-history}.
Applying that proof with estimated costs yields
\begin{align}
    \edhistory{b}{t}
    \le\edcost{t}(F_t^*)+c|B_{Q_b}^t(F_t^*)|
    \le\edcost{t}(F_t^*)+(R-1)\OPT_t.
    \nonumber
\end{align}
\end{proof}

\begin{lemma}[Approximation at every prefix]\label{lem:explicit.selected.cost}
With probability at least $ 1-\zeta $, the output satisfies $ \cost_t(C_t)\le(R_0+\eps)\OPT_t $ simultaneously for every $ t\le n $.
\end{lemma}
\begin{proof}
By \cref{lem:explicit.good.candidate,lem:explicit.estimated.costs} and a union bound, with probability at least $ 1-3\zeta/4 $ every prefix has a good candidate and satisfies \eqref{eq:explicit.consensus.ceiling} and \eqref{eq:explicit.consensus.comparison}.
On this event, fix $ t $ and apply \cref{lem:explicit.estimated.history} to a good candidate at $ t $.
The proof of \cref{lem:oncon-selection} uses only the monotonicity of history scores and the fact that a current score is at most its history score.
It therefore applies without change and gives $ \edcost{t}(C_t)\le(1+\theta)\edminimum{t} $.
Combine this inequality with \eqref{eq:explicit.consensus.good.history}, \eqref{eq:explicit.consensus.comparison} for the selected candidate, and $ d(C_t,F_t^*)\le\cost_t(C_t)+\OPT_t $ to obtain
\begin{align}
    (1-\gamma)\cost_t(C_t)
    &\le(1+\gamma)\OPT_t+\theta\edcost{t}(F_t^*)
        +(1+\theta)(R-1+3\gamma)\OPT_t\nonumber\\
    &\le\bigl(R+4\gamma+\theta(R+1+3\gamma)\bigr)\OPT_t.
    \label{eq:explicit.consensus.final.cost}
\end{align}
The last inequality uses \eqref{eq:explicit.consensus.ceiling}.
Since $ \theta=\gamma=\eps/40<1/160 $ and $ R<2 $,
\begin{align}
    \frac{R+4\gamma+\gamma(R+1+3\gamma)}{1-\gamma}
    &=R+\frac{\gamma(2R+5+3\gamma)}{1-\gamma}
    <R+10\gamma
    =R_0+\frac{3\eps}{4}
    <R_0+\eps.
    \nonumber
\end{align}
This holds for every prefix on the stated event, including prefixes with $ \OPT_t=0 $.
\end{proof}

\subsection{Implementation and Update Time}\label{subsec:explicit.consensus.implementation}


\begin{lemma}[Maintaining the estimated costs]\label{lem:explicit.cost.maintenance}
For each candidate, its $ L $ evaluation-group means can be maintained in $ O(ns) $ words and $ O(s\log(n+2)) $ time per reported label assignment, where $ s=Lg $.
Computing the estimated cost from the updated means takes $ O(L) $ time.
\end{lemma}
\begin{proof}
Apply the counters and update formulas of \cref{subsec:dynamic-algorithm} separately to each evaluation group, with group size $ g $.
Restore the input-only term following \eqref{eq:dynamic-score} to obtain the full group mean before taking its minimum or updating a history score.
That term cannot be omitted here, since it varies across groups and times.
For each candidate, keep its previous labels while processing the arriving vertex's assignment and then the old vertices' reported changes.
A change to a previously unused label is an insertion into an empty block.
As in \cref{subsec:dynamic-algorithm}, each assignment or move takes $ O(s\log(n+2)) $ time, including updates to all group sums.
Each of the $ s $ evaluation positions has at most $ n $ nonempty block-intersection records, giving the stated space bound.
Computing the minimum over the $ L $ groups takes $ O(L) $ time per candidate per arrival.
\end{proof}

\begin{lemma}[Cost of rebuilding the explicit output]\label{lem:explicit.rebuilding}
On every execution, the number $ S $ of outer switches satisfies
\begin{align}
    S\le K\left(2+\left\lceil\log_{1+\theta}(gn^2)\right\rceil\right)
    =O\!\left(K\eps^{-1}\log(gn+2)\right).
    \label{eq:explicit.consensus.switches}
\end{align}
Over the first $ t $ insertions, rebuilding the output at switches takes $ O(t(S+1)) $ time.
\end{lemma}
\begin{proof}
Every estimated score, history score, and minimum history score belongs to $ g^{-1}\mathbb Z\cap[0,\binom n2] $.
Applying the phase-count argument in the proof of \cref{lem:oncon-recourse}, with minimum positive score $ 1/g $ and the present value of $ \theta $, gives \eqref{eq:explicit.consensus.switches}.
Each rebuild writes and reports at most $ t $ identifiers during the first $ t $ insertions.
Between switches, output writes and reports are charged to assignments reported by the selected candidate.
\end{proof}

\begin{lemma}[Total update time]\label{lem:explicit.total.work}
For every $ t\le n $, the expected total time for the first $ t $ insertions, including initialization and all output reports, is
\begin{align}
    O\!\left(
        tK(\edTime{n}+1)s\log(n+2)
        +tK\eps^{-1}\log(gn+2)
    \right).
    \label{eq:explicit.consensus.total.work}
\end{align}
\end{lemma}
\begin{proof}
For every fixed choice of candidate samples, the expected work of all solver copies over the first $ t $ insertions, including their queries and reported assignments, is $ O(tK\edTime{n}) $.
The bucket representation implements each database query in $ O(\log(n+2)) $ time, and \cref{lem:explicit.cost.maintenance} charges $ O(s\log(n+2)) $ time to each reported assignment.
Reading the sampled labels, updating input candidates, and computing scores and selections fit in $ O(tKs\log(n+2)) $ time.
By \cref{lem:explicit.rebuilding}, output rebuilds contribute $ O(tK\eps^{-1}\log(gn+2)) $ time.
Adding these bounds proves \eqref{eq:explicit.consensus.total.work}.
\end{proof}

\begin{lemma}[Bounding the work of each solver copy]\label{lem:explicit.capped.work}
Given the prefix expectation bound in \cref{lem:explicit.consensus.clustering}, one can modify \cref{alg:explicit.consensus} so that \eqref{eq:explicit.consensus.time} bounds the amortized update time at every prefix on every execution.
The modified algorithm retains the simultaneous approximation guarantee of \cref{lem:explicit.selected.cost}.
\end{lemma}
\begin{proof}
Double $ K $ in \eqref{eq:explicit.dynamic.consensus.K} and use this larger value when setting $ g $.
Set $ h=1+\lceil\log_2 n\rceil $.
At insertion $ t $, stop a triple solver permanently if its next operation would bring its cumulative work above $ 8ht\edTime{n} $.
Count only the solver's own work against its budget, and buffer its reports until the insertion finishes.
If the solver stops, discard its pending reports and use the first input partition for that candidate from this insertion onward.
Cache that input's labels as vertices arrive, and rebuild the candidate's labels and score counters once when it switches to that input.
Use fresh identifiers for its input blocks, report the resulting assignments, and retain the candidate's previous history score when taking the running maximum.

Fix a sampled triple and let $ W(u) $ be the cumulative work of its uninterrupted solver through insertion $ u $.
For any $ t\le n $, the first endpoint $ u\ge t $ in $ \{\min(2^j,n):0\le j<h\} $ satisfies $ u\le2t $.
If the budget is exceeded at $ t $, monotonicity gives $ W(u)>4hu\edTime{n} $.
Since $ \E[W(u)]\le u\edTime{n} $, Markov's inequality and a union bound over the at most $ h $ endpoints bound the probability of stopping by $ 1/4 $.
Thus, conditional on any triple, its solver both succeeds at all prefixes and finishes the entire stream with probability at least $ 1/2-1/4=1/4 $.
Using this probability in the proof of \cref{lem:explicit.good.candidate} gives success probability at least $ \eps/16 $ per candidate at each fixed prefix, for which the doubled $ K $ suffices.
The modified trajectories are independent of the evaluation samples, so \cref{lem:explicit.estimated.costs,lem:explicit.estimated.history,lem:explicit.selected.cost} apply as before.

Through insertion $ t $, each solver spends at most $ 8ht\edTime{n} $ operations and rebuilds at most once, at a cost of $ O(ts\log(n+2)) $ for that rebuild.
Consequently the accounting of \cref{lem:explicit.total.work} applies on every execution with $ 8h\edTime{n} $ in place of $ \edTime{n} $.
The additional factor $ h=O(\log(n+2)) $ is absorbed by the polylogarithmic factor in \eqref{eq:explicit.consensus.time}.
\end{proof}

\begin{proof}[Completion of the proof of~\cref{lem:explicit.consensus.clustering}]
The approximation guarantee follows from \cref{lem:explicit.selected.cost}.
By \eqref{eq:explicit.dynamic.consensus.K} and \eqref{eq:explicit.consensus.parameters},
\begin{align}
    K&=O\!\left(\eps^{-1}\log\frac{2n}{\zeta}\right),
    &L&=O\!\left(\log\frac{2n}{\zeta}\right),
    &g&=O\!\left(\eps^{-2}\log\frac{2n}{\eps\zeta}\right).
    \nonumber
\end{align}
Substituting $ s=Lg $ and $ \edTime{n}=\polylog(n+2) $ into \eqref{eq:explicit.consensus.total.work} proves \eqref{eq:explicit.consensus.time}.
\end{proof}

\section{Explicit Dynamic Correlation Clustering Algorithm for Insertion-Only Stream}\label{sec:explicit_dynamic_correlation_clustering_algorithm_for_insertion_only_stream} % (fold)

In this section, we prove the following lemma.
\ExplicitDA*

We present the algorithm $ \explicitDA $, which is an explicit algorithm for correlation clustering.
We demonstrate that $ \explicitDA $ obtains the properties stated in \cref{lem:explicit.correlation.clustering}.

Throughout this section, we denote by $ G $ the underlying graph of the problem, and by $ G_{t} $ the induced subgraph of $ G $ on the first $ t $ vertices.
\subsection{Algorithm \texorpdfstring{$ \explicitDA $}{ExplicitDA}}\label{sub:algorithm.explictDA} % (fold)
The algorithm $ \explicitDA $ (\cref{alg.explicit.dynamic.agreement}) builds on an amplified version of $ \DA $ (\cite{CLMP2024Dynamic}), denoted by $ \ADA $, with additional data structures and procedures for reporting label changes.

$ \DA $ is a randomized dynamic correlation-clustering algorithm.
For each fixed time $ t $, its maintained clustering is an $ \edFactor $-approximation with probability at least $ 1-5/n $, for an absolute constant $ \edFactor $ \cite[Theorem~6]{CLMP2024Dynamic}.
Its expected amortized update time is $ \polylog(n) $ \cite[Appendix~C, proof of Theorem~41]{CLMP2024Dynamic}.
The algorithm does not report the vertices whose cluster labels change after an update; instead, its output procedure enumerates the entire clustering when requested \cite[Appendix~D]{CLMP2024Dynamic}.

We amplify $ \DA $ to obtain $ \ADA $, which maintains an $ \edFactor $-approximation simultaneously at all prefixes with probability at least $ 1-n^{-2} $, while retaining $ \polylog(n) $ expected amortized update time.
The amplification and its analysis are given in \cref{sub:amplify.da}.

Whenever a vertex $ v_t $ arrives, $ \explicitDA $ executes $ \ADA $ to update the sparse graph and anchor set.
To report the resulting label changes, it maintains the following additional data structures:
\begin{itemize}
    \item A journal $ \journal $ records every edge insertion, edge deletion, and change of anchor membership during the execution of $ \ADA $.
    For each edge modified during the update, the journal also records whether it was present immediately before its first modification.

    We denote by $ \journal{t} $ the changes recorded in the journal during the execution of $ \ADA $ on the arrival of $ v_t $.

    \item The dynamic connectivity structure of \cite{holm2001poly}, denoted by $ \dyCC $, maintains the connected components of $ \sparseG $ under edge insertions and deletions.
    It supports each edge update in $ O(\log^2 n) $ amortized time and a query $ \comp{x} $ in $ O(\log n) $ time.
    At a fixed moment, $ \comp{x}=\comp{y} $ if and only if $ x $ and $ y $ belong to the same component of the current graph.
    These component identifiers need not remain stable between updates.

    \item After $ \ADA $ finishes an update, we use a $ \LabelSync $ procedure, which takes as input the journal $ \journal $, the $ \dyCC $ data structure, to maintain a data structure that assigns labels to arrived vertices.
        At each time step $ t $, two vertices $ x $ and $ y $ have the same label if and only if they belong to the same connected component of the current sparse graph $ \sparseG{t} $.
        The procedure $ \LabelSync $ updates the component identifiers of the vertices in $ \sparseG $ to reflect the connected components of the current sparse graph, and reports every vertex whose component identifier changed during update $ t $.
    On the good event defined below, two vertices have the same label exactly when they belong to the same component of $ \sparseG{t} $.
\end{itemize}

\begin{algorithm}[htbp]
    \SetKwInOut{KwIn}{Input}
    \SetKwInOut{KwOut}{Output}
    \KwIn{A known horizon $ n $ and an oblivious stream of at most $ n $ vertex insertions, accessed through the database model of \cref{def:query.model.label}.}
    \KwOut{After each insertion, the new vertex and its label, followed by every previously arrived vertex whose label changed and its new label.}
    Initialize the maintained sparse graph $ \sparseG \gets \emptyset $

    Initialize the dynamic connectivity structure $ \dyCC $

    Initialize $ \LabelSync $

    \On{arrival of $ v_t $}{
        Initialize an empty journal $ \journal{t} $ and insert $ v_t $ into $ \dyCC $ as an isolated vertex

        Execute $ \ADA(v_t) $, recording every edge operation and anchor-membership change in $ \journal{t} $ and applying every edge operation to $ \dyCC $

        Execute $ \LabelSync{t} $ using $ \journal{t} $ and $ \dyCC $
    }
    \caption{$ \explicitDA $: Explicit Dynamic Agreement Algorithm for Insertion-Only Streams}
    \label{alg.explicit.dynamic.agreement}
\end{algorithm}

To prove \cref{lem:explicit.correlation.clustering}, we state the following results about $ \ADA $ and $ \LabelSync $.

\begin{restatable}{lemma}{AmplifiedDynamicAgreementLemma}
    \label{lem:ada.guarantees}
    Fix an oblivious insertion-only stream of at most $ n $ vertices.
    For sufficiently large $ n $, with probability at least $ 1-n^{-2} $, the connected components of the sparse graph maintained by $ \ADA $ give an $ \edFactor $-approximation simultaneously at every prefix.
    For every $ t\le n $, the expected total work of $ \ADA $ through insertion $ t $ is $ O(t\polylog(n)) $.
\end{restatable}

The following lemma summarizes the correctness and implementation cost of $ \LabelSync $, established in \cref{lem:labelsync.classes.are.connected.components,lem:labelsync.cost}.

\begin{restatable}{lemma}{LabelSyncLemma}
    \label{lem:labelsync.guarantees}
    For sufficiently large $ n $ and any fixed oblivious insertion-only stream processed by $ \ADA $, with probability at least $ 1-n^{-2} $, $ \LabelSync $ correctly reports the labels of vertices whose labels changed during each update.
    This holds jointly with the approximation guarantee in \cref{lem:ada.guarantees}.
    For every execution, the work of $ \LabelSync $ during update $ t $, including all reports, is
    \begin{align}
        O((\jnTime{t}+1)\log n+\nChangedVers{t}),
    \end{align}
    where $ \jnTime{t} $ is the journal length and $ \nChangedVers{t} $ counts changed labels of previously arrived vertices.
\end{restatable}

\begin{proof}[Proof of~\cref{lem:explicit.correlation.clustering}]
    By \cref{lem:ada.guarantees}, with probability at least $ 1-n^{-2} $, $ \ADA $ maintains an $ \edFactor $-approximation.
    On each of these executions, by~\cref{lem:labelsync.guarantees}, the labels maintained by $ \LabelSync $ describe exactly the clustering maintained by $ \ADA $ at every prefix.

    Fix a prefix of $ t $ insertions, and let $ W_t $ be the work performed by $ \ADA $ during this prefix.
    \cref{lem:ada.guarantees} gives
    \begin{align}
        \mathbb E[W_t]=O(t\polylog(n+2)).
    \end{align}
    Recording an edge operation or an anchor-membership change takes constant additional work.
    Every journal entry records an operation performed by $ \ADA $, so the total journal length
    \begin{align}
        H_t:=\sum_{s=1}^{t}\jnTime{s}
    \end{align}
    satisfies $ H_t=O(W_t+t) $ on every execution.
    In particular, $ \sum_{s=1}^{t}\card{\delt{s}}\le H_t $, and $ \mathbb E[H_t]=O(t\polylog(n+2)) $.

    By \cref{lem:number.of.updated.vertices}, the total number of changed labels of previously arrived vertices satisfies
    \begin{align}
        \sum_{s=1}^{t}\nChangedVers{s}
        &=O\left(t\log n+\log n\sum_{s=1}^{t}\card{\delt{s}}\right) \\
        &=O((H_t+t)\log n).
    \end{align}
    The newly inserted vertices contribute another $ t $ reported label assignments.
    Summing the work bound in \cref{lem:labelsync.guarantees} over the prefix, the total work of $ \LabelSync $, including its reports, is therefore $ O((H_t+t)\log n) $.
    There are at most $ H_t $ journalled edge operations, each requiring $ O(\log^2 n) $ amortized time in $ \dyCC $ by \cite{holm2001poly}; together with vertex initialization, maintaining $ \dyCC $ costs $ O((H_t+t)\log^2 n) $ time over the prefix.
    Thus the total work of $ \explicitDA $ through insertion $ t $ is
    \begin{align}
        O(W_t+(H_t+t)\log^2 n).
    \end{align}
    Taking expectations yields $ O(t\polylog(n+2)) $, including all reported label assignments.
    The graph queries are implemented in the database model of \cref{def:query.model.label}.
    For the remaining bounded horizons, we can maintain an optimal clustering by exhaustive computation with constant work per insertion, including all reports.
\end{proof}

\subsection{Amplify \texorpdfstring{$ \DA $}{DynamicAgreement} Algorithm}
\label{sub:amplify.da}
We describe $ \ADA $.

We first review $ \DA $.
At every time $ t $, $ \DA $ maintains a sparse subgraph $ \sparseG{t} $ of the input graph $ G_t $ and an anchor set $ \anchorSet{t} $ satisfying
\begin{align}
    \text{every edge of } \sparseG{t} \text{ has an endpoint in } \anchorSet{t}. \tag{I} \label{eq:invariant.anchor}
\end{align}
The maintained clustering is the partition into connected components of $ \sparseG{t} $.

Throughout the whole course of updates, $ \DA $ iteratively utilizes two queries: $ \ProbAgree$ (\cite[Algorithm 8]{CLMP2024Dynamic}) and $ \Heavy $ (\cite[Algorithm 9]{CLMP2024Dynamic}).
They show that each query answer correctly with probability at least $ 1-n/2 $ (\cite[Lemma 49, Lemma 50, Lemma 52, Lemma 53]{CLMP2024Dynamic}).

To obtain $ \ADA $, we replace each decision of the randomized tests ProbabilisticAgreement and Heavy by the majority of three independent executions.
One execution of Heavy includes its internal calls to ProbabilisticAgreement, which are left unchanged.
As shown in \cref{lem:da.amplified.tests}, amplification reduces the error probability of each test decision to at most $ 3n^{-6} $, with only constant-factor overhead per decision.
The resulting whole-run guarantee is stated below and proved in \cref{sub:amplify.da}.

\paragraph{Amplifying the randomized tests.}
Dynamic Agreement uses the randomized tests ProbabilisticAgreement and Heavy.
A response is called correct if it satisfies the acceptance and rejection guarantees established for that test in \cite[Appendix~E]{CLMP2024Dynamic}, with the corrections described above.
Either response is permitted in the intermediate regimes where these guarantees prescribe no answer.

\begin{lemma}[Amplified agreement and heaviness tests]
    \label{lem:da.amplified.tests}
    Replace each test decision by the majority of three independent executions of the corresponding test.
    Conditional on the history before the decision, the amplified test gives a correct response with probability at least $ 1-3n^{-6} $.
    Amplification incurs only a constant-factor overhead per test decision.
\end{lemma}

\begin{proof}
    By \cite[Lemmas~49 and~50]{CLMP2024Dynamic}, one execution of ProbabilisticAgreement gives an incorrect response with probability at most $ n^{-3} $.
    The analogous guarantee for Heavy follows from \cite[Lemmas~52 and~53]{CLMP2024Dynamic}, with the corrected rejection guarantee.

    Conditioning on the preceding history fixes the graph and the test input.
    Since the three executions use independent fresh randomness, their errors are independent under this conditioning.
    The majority response can be incorrect only if at least two executions are incorrect.
    A union bound over the three pairs bounds this probability by $ 3n^{-6} $.

    Each amplified decision performs three executions of the original test.
    One execution of Heavy includes its internal calls to ProbabilisticAgreement, which are left unchanged.
\end{proof}

\begin{lemma}[Number of test invocations,
    {\cite[proof of Theorem~41, Appendix~C]{CLMP2024Dynamic}}]
    \label{lem:da.test.count}
    Fix an oblivious insertion-only stream of length at most $ n $.
    Let $ Q_t $ denote the number of agreement and heaviness test invocations during update $ t $.
    Then $ \mathbb E[Q_t]=O(\polylog(n)) $.
    Consequently, for every prefix of length $ t $,
    \begin{align}
        \mathbb E\left[\sum_{s=1}^{t}Q_s\right]
        =O(t\polylog(n)).
    \end{align}
\end{lemma}

\paragraph{The good event for Dynamic Agreement.}
Fix an oblivious insertion-only stream of length $ m\le n $.
Following the proof of \cite[Theorem~6]{CLMP2024Dynamic}, we separate test correctness from the remaining structural events.

Let $ \mathcal E_{\mathrm{test}}(t) $ denote the event that every amplified test decision through update $ t $ gives a correct response.
For each fixed prefix $ t\le m $, we use the following two structural events from that analysis:
\begin{enumerate}
    \item $ \mathcal A_t $ is the event that every reference cluster considered in \cite[Theorem~23]{CLMP2024Dynamic} is contained in one component of the maintained sparse graph.
    The bound in that theorem, together with the small-cluster correction and a union bound over at most $ n $ reference clusters, gives
    \begin{align}
        \Pr[
            \mathcal E_{\mathrm{test}}(t)
            \cap \overline{\mathcal A_t}
        ]
        \le n^{-7}.
    \end{align}

    \item $ \mathcal R_t $ is the aggregate regularity event used in the density analysis of \cite[Definition~28 and Observation~29]{CLMP2024Dynamic}.
    The estimates in \cite[Lemmas~25 and~27]{CLMP2024Dynamic}, aggregated over the relevant vertices and indices as in that analysis, give
    \begin{align}
        \Pr[\overline{\mathcal R_t}]
        \le n^{-100}.
    \end{align}
\end{enumerate}

On $ \mathcal E_{\mathrm{test}}(t)\cap\mathcal A_t\cap\mathcal R_t $, the argument in \cite[proof of Theorem~6]{CLMP2024Dynamic} implies that the sparse connected components give an $ \edFactor $-approximation, for an absolute constant $ \edFactor $.
Moreover, on $ \mathcal E_{\mathrm{test}}(t)\cap\mathcal R_t $, the density analysis in \cite[proof of Theorem~39]{CLMP2024Dynamic} gives the anchor-density property stated in \cref{lem:ccclass.is.dense}.

Define
\begin{align}
    \successDA
    :=
    \mathcal E_{\mathrm{test}}(m)
    \cap
    \bigcap_{t=1}^{m}
    (\mathcal A_t\cap\mathcal R_t).
\end{align}

\begin{lemma}[Anchor density, adapted from {\cite[proof of Theorem~39]{CLMP2024Dynamic}}]
    \label{lem:ccclass.is.dense}
    On $ \successDA $, for every time $ t\le m $, every component $ \ccclass $ of $ \sparseG{t} $ with $ \card{\ccclass}>1 $, and every $ a\in \ccclass\cap\anchorSet{t} $,
    \begin{align}
        |\openN{\sparseG{t}}{a}|\ge\tfrac34\card{\ccclass}-1.
    \end{align}
\end{lemma}

\begin{restatable}{lemma}{successDAProb}
    \label{lem:da.whole.run.success}
    For sufficiently large $ n $, the amplified algorithm satisfies
    \begin{align}
        \Pr[\successDA]\ge1-n^{-2}.
    \end{align}
    On this event, its sparse connected components give an $ \edFactor $-approximation and satisfy the anchor-density property simultaneously at every prefix.
\end{restatable}

\begin{proof}
    The insertion-only algorithm makes at most $ B_n=O(n^2\log^5 n) $ test decisions over the entire stream.
    This deterministic bound follows by counting test invocations in the algorithm and using \cite[Lemma~40]{CLMP2024Dynamic}.
    It holds independently of the test responses.

    By \cref{lem:da.amplified.tests} and a union bound over invocation indices,
    \begin{align}
        \Pr[\overline{\mathcal E_{\mathrm{test}}(m)}]
        \le 3B_n n^{-6}.
    \end{align}
    Combining this with the imported structural bounds gives
    \begin{align}
        \Pr[\overline{\successDA}]
        &\le
        \Pr[\overline{\mathcal E_{\mathrm{test}}(m)}]
        +
        \sum_{t=1}^{m}
        \Pr[
            \mathcal E_{\mathrm{test}}(t)
            \cap\overline{\mathcal A_t}
        ]
        +
        \sum_{t=1}^{m}
        \Pr[\overline{\mathcal R_t}] \\
        &\le 3B_n n^{-6}+n^{-6}+n^{-99} \\
        &=O(n^{-4}\log^5 n)+n^{-6}+n^{-99}
        \le n^{-2}
    \end{align}
    for sufficiently large $ n $.
    The approximation and density conclusions follow from the implications stated above.
\end{proof}

% subsection Auxiliary Lemmas From~ (end)


\subsection{Description of \texorpdfstring{$ \LabelSync $}{LabelSync} Algorithm}\label{sub:description_of_LabelSync} % (fold)

% \paragraph{\texorpdfstring{$ \LabelSync $}{LabelSync} Algorithm}\label{par:_texorpdfstring_labelsync_labelsync_algorithm} % (fold)
At the end of update $ t $, $ \LabelSync $ maintains a partition $ \ccclasses{t} $ of the current vertex set $ V_t $.
We call each member of $ \ccclasses{t} $ a \emph{class}.
Each class $ \ccclass $ has a distinct label $ \complabel(\ccclass) $, and every vertex $ w\in\ccclass $ stores this label as $ \complabel_t(w)=\complabel(\ccclass) $.
The partition $ \ccclasses{t} $ is the solution reported by $ \explicitDA $.
At the end of each update $ t $, the algorithm reports only the vertices that changed clusters, that is, $ (v_t,\complabel_t(v_t)) $ and every $ (w,\complabel_t(w)) $ with $ w\in V_{t-1} $ and $ \complabel_t(w)\ne\complabel_{t-1}(w) $.
We prove in \cref{sub:analysis_for_labelsync} that, on the success event $ \successDA $, the classes in $ \ccclasses{t} $ are exactly the connected components of $ \sparseG{t} $.
% paragraph LabelSync Algorithm (end)

\tienlong{Brevity: $ \ADA(t) $}
\paragraph{Overview of \texorpdfstring{$ \LabelSync $}{LabelSync}}\label{par:overview.of.LabelSync} % (fold)
The input to $ \LabelSync $ includes the sparse graph $ \sparseG{t} $ (after the update from time $ t-1 $ to time $ t $), the journal $ \journal $ of update $ t $, the arriving vertex $ v_t $, the data structure of $ \LabelSync $ at time $ t-1 $.

The goal is to update the data structure to reflect the connected components of $ \sparseG{t} $ and to report the vertices whose labels changed.

On a high level, the algorithm achieves this by performing a \emph{split} and a \emph{merge} operation.
Suppose that at time step $ t-1 $, $ \LabelSync $ successfully maintained the $ \ccclasses{t-1} $ as the connected components of $ \sparseG{t-1} $.
At time step $ t $, vertices in some components in $ \ccclasses{t-1} $ may not belong to the same components in $ \sparseG{t} $, and some vertices in different components in $ \ccclasses{t-1} $ may belong to the same component in $ \sparseG{t} $.
To fix this, we first perform the split operation, which partitions each old class $ \ccclass \in \ccclasses{t-1} $ into \piece{s}, where each \piece{} is a subset of vertices of $ \ccclass $ that belong to the same component in $ \sparseG{t} $.
After that, the merge operation groups \piece{s} that belong to the same component in $ \sparseG{t} $ together.
Finally, we update the labels of vertices in the new classes and report the vertices whose labels changed.

A direct approach to perform the split and merge operations would traverse the entire graph $ \sparseG{t} $ to find the connected components, which in the worst case takes $ O(n\polylog(n)) $ time.
To avoid this, we follow vertices and edges affected by $ \ADA(t) $ to perform the split and merge operations.
Note that, the length of the journal $ \journal{t} $ is at most the update time of $ \ADA $, therefore, we can afford to scan the whole journal for each update $ t $.
By reading off $ \journal{t} $, the algorithm can identify edges and vertices that were updated after the execution of $ \ADA $ for update $ t $.
It suffices to focus on the edges that were deleted to perform the split operation.
In particular, consider a component $ \ccclass \in \ccclasses{t-1} $.
If there are many edges inside $ \ccclass $ deleted by $ \ADA(t) $, we can safely scan the entire component $ \ccclass $ to split it into \piece{s}.
On the other hand, if only a few edges inside $ \ccclass $ were deleted, we demonstrate that in this case, all the vertices in $ \ccclass \cap \anchorSet{t-1} $ remain in the same component in $ \sparseG{t} $.
With the invariant~\eqref{eq:invariant.anchor} being maintained, one endpoint of each deleted edge is in $ \anchorSet{t-1} $.
This enables us to perform the split operation for $ \ccclass $ solely on the deleted edges and their endpoints not in $ \anchorSet{t-1} $.
After the split step, the merge step now operates on the \piece{s} created in the split step and the other components in $ \ccclasses{t-1} $ that have edges added by $ \ADA(t) $.
While merging \piece{s} in the same component of $ \sparseG{t} $ together, we take the label of a largest \piece{} in the component and assign it to all the vertices in the component.
Under this rule, intuitively, every time a vertex changes its label, the size of its component at least doubles.
Thus, on avarage, each vertex can change its label at most $ O(\log n) $ times, which gives a bound on the number of vertices reported by $ \LabelSync $.

% paragraph Overview of LabelSync (end)

\paragraph{Detailed Description of \texorpdfstring{$ \LabelSync $}{LabelSync}}\label{par:detailed_description_of_labelsync_} % (fold)
We provide the algorithm in more details with its main steps, data structure and the implementation.
The summary of the algorithm is given in \cref{alg.label.sync}.

Prior to the execution of $ \LabelSync{t} $, each class $ \ccclass \in \ccclasses{t-1} $ is represented by a record containing its label, size, and a doubly linked list of its vertices.
A set data structure is maintained to store the records of classes in $ \ccclasses{t-1} $.
We refer to this set as $ \setClasses$.
Each vertex $ w $ of $ \ccclass $ keeps a record of its label, a pointer to the record of $ \ccclass $, a pointer to $ w $'s position in $ \ccclass $'s doubly linked list, $ w $'s membership in $ \anchorSet{t-1} $, and its degree $ \deg_{t-1}(w):=\card{\openN{\sparseG{t-1}}{w}} $ in the sparse graph.

The algorithm consists of the following four steps.

\noindent\setword{\textbf{Step 1 (Identify Updated Edges)}}{StepOneLabelSync}:
In this step, we identify the edges that are updated after executing $ \ADA $ to handle update $ t $.
In particular, we compute
\begin{align}
    \delt{t}:=E(\sparseG{t-1})\setminus E(\sparseG{t}), \nonumber
    \qquad
    \addt{t}:=E(\sparseG{t})\setminus E(\sparseG{t-1}), \nonumber
\end{align}
the sets of edges deleted and added between $ \sparseG{t-1} $ and $ \sparseG{t} $ by $ \ADA $, respectively.
Note that, during the execution of $ \ADA $, an edge between two vertices $ u $ and $ v $ may be deleted and then added back, or added and then deleted, multiple times.
Here, $ \delt{t} $ and $ \addt{t} $ only record the edges that are present in $ \sparseG{t-1} $ (before update $ t $) but not in $ \sparseG{t} $, and vice versa, after $ \ADA $ finishes.

Additionally, for each old class $ \ccclass\in\ccclasses{t-1} $, let $ D_{\ccclass}:=\{(x,y)\in\delt{t}:x,y\in\ccclass\} $ be its deleted edges, and let $ d_{\ccclass}:=\card{D_{\ccclass}} $.
For each endpoint $ w $ of an edge in $ D_{\ccclass} $, let $ \delIncEdge{w}:=\card{\{e\in D_{\ccclass}:w\in e\}} $.

\textit{Implementation and data structures.}
We store $ \delt{t} $ and $ \addt{t} $ as lists, with each edge occurring at most once.
For each edge considered during the execution of $ \ADA $ for update $ t $, the journal $ \journal $ records its status in $ \sparseG{t-1} $ and $ \sparseG{t} $.
Hence, one pass over the journal establishes $ \delt{t} $ and $ \addt{t} $.
For each old class $ \ccclass $ with $ D_{\ccclass}\ne\emptyset $, we have a dictionary mapping its label $ \complabel(\ccclass) $ to the pair $ (D_{\ccclass},d_{\ccclass}) $, where $ D_{\ccclass} $ is a list of deleted edges and $ d_{\ccclass}=\card{D_{\ccclass}} $.
A second dictionary is employed to map each endpoint $ w $ of a deleted edge to its count $ \delIncEdge{w} $.
We fill both dictionaries in one pass over $ \delt{t} $.

\noindent\setword{\textbf{Step 2 (Split)}}{StepTwoLabelSync}:
We partition each old class into temporary sets called \emph{\piece{s}}, such that each \piece{} lies entirely within one component of $ \sparseG{t} $.

To do that, observe that conditioned on $ \successDA $, the \piece{s} of an old class $ \ccclass $ will be exactly
\begin{align}
    \mathcal P_{\ccclass}
    =\{\ccclass\cap K:K\text{ is a component of }\sparseG{t},\ \ccclass\cap K\ne\emptyset\}. \nonumber
\end{align}
Among the \piece{s} in $ \mathcal P_{\ccclass} $, we choose one to inherit the old label $ \complabel(\ccclass) $.
We refer to it as the \emph{retained \piece{}} of $ \ccclass $ and denote it by $ \retainedPiece{\ccclass} $.
The other \piece{s} initially have no assigned label.

If $ d_{\ccclass}=0 $, the class is represented by the single retained \piece{} $ \retainedPiece{\ccclass}=\ccclass $.
For each class with $ d_{\ccclass}>0 $, we distinguish two cases.
\begin{itemize}
    \item If $ d_{\ccclass}\ge\card{\ccclass}/2-1 $, we traverse $ \ccclass $ and group its vertices into \piece{s} such that, $ x,y \in \ccclass $ are grouped into the same \piece{} if and only if $ \comp{x}=\comp{y} $.

    A largest \piece{} is chosen as $ \retainedPiece{\ccclass} $.
    % Since $ \card{\ccclass}\le 2d_{\ccclass}+2\le 4d_{\ccclass} $, the number of vertices examined is proportional to $ d_{\ccclass} $.
    \item If $ d_{\ccclass}<\card{\ccclass}/2-1 $, choose an endpoint $ r\in\ccclass\cap\anchorSet{t-1} $ of an edge in $ D_{\ccclass} $.
    Such an endpoint exists because every old edge has an endpoint in $ \anchorSet{t-1} $, by definition of the anchor set.
    Define
    \begin{align}
        W_{\ccclass}:=\{w|\ \exists e\in D_{\ccclass}: w\in e, \ \text{and }
        w \in \ccclass\setminus\anchorSet{t-1},\ \delIncEdge{w}=\deg_{t-1}(w)\}, \nonumber
    \end{align}
    the set of vertices in $ \ccclass\setminus \anchorSet{t-1} $ whose incident edges at time $ t-1 $ are all deleted during update $ t $. \tienlong{Need to explain about $ \delIncEdge{w}=\deg_{t-1}(w) $}
    We will show that (\cref{lem:labelsync.robust.components}), on $ \successDA $, every vertex of $ \ccclass\setminus W_{\ccclass} $ remains connected to $ r $ after the deletions.
    For each vertex $ w\in W_{\ccclass} $, we group it with the same \piece{} as $ r $ if and only if $ \comp{w}=\comp{r} $.
    For the remaining vertices in $ W_{\ccclass} $, we group them into \piece{s}
    \begin{align}
        P_{\ccclass,\kappa}:=\{w\in W_{\ccclass}:\comp{w}=\kappa\}, \nonumber
    \end{align}
    for each $ \kappa\ne\comp{r} $ attained in $ W_{\ccclass} $.
    The retained \piece{} is
    \[
        \retainedPiece{\ccclass}:=\ccclass\setminus\bigcup_{\kappa\ne\comp{r}}P_{\ccclass,\kappa},
    \]
    where the union ranges only over the identifiers attained in $ W_{\ccclass} $.
    % We store $ R_{\ccclass} $ implicitly, with representative $ r $ and size $ \card{\ccclass}-\sum_{\kappa\ne\comp{r}}\card{P_{\ccclass,\kappa}} $.
    % The set $ W_{\ccclass} $ is found from the endpoints of $ D_{\ccclass} $, and its vertices are grouped using a hash map indexed by their component identifiers.
    % Thus this case does not require listing $ \ccclass $.
\end{itemize}

The \piece{s} of all old classes, together with the singleton $ \{v_t\} $, form a partition $ \mathcal P $ of $ V_t $. \tienlong{Not clear}

\textit{Implementation and data structures.}
For each class $ \ccclass $ participating in this step, we remove $ \ccclass $ from $ \setClasses $.

For each such $ \ccclass $, to group vertices of $ \ccclass $ with the same component identifier into one \piece{}, we utilize a dictionary mapping each identifier to the record containing a doubly linked list of vertices with that identifier.
We refer to this dictionary as $ \groupdict(\ccclass) $ for each such $ \ccclass $.

This dictionary is constructed as follows.

If $ d_{\ccclass}\ge\card{\ccclass}/2-1 $, the algorithm builds $ \groupdict(\ccclass) $ by traversing the vertices of $ \ccclass $.
In particular, for each vertex $ w\in \ccclass $, we add $ w $ to the doubly linked list of the record corresponding to the key $ \comp{w} $ in $ \groupdict(\ccclass) $.
We remove $ w $ from $ \ccclass $'s list, update $ w $'s pointers to the corresponding record in $ \groupdict(\ccclass) $ and its new position in the list of this record.
After the traversal, the record of a largest \piece{} in $ \groupdict(\ccclass) $ takes the label of $ \ccclass $, and the other records have no assigned label.

If $ 0<d_{\ccclass}<\card{\ccclass}/2-1 $, the algorithm first finds $ r $ and constructs the set $ W_{\ccclass} $ by examining each endpoints of $ D_{\ccclass} $.
We build $ \groupdict(\ccclass) $ just like in the previous case, but by traversing only the vertices in $ W_{\ccclass} $ instead of the whole class $ \ccclass $.
Also, during the traversal, we do not create a key in $ \groupdict(\ccclass) $ and the corresponding record for vertices with the same component identifier as $ r $, and also not remove them from the list of $ \ccclass $'s record.
After the traversal, the list of the record of the original class $ \ccclass $ now contains only the vertices of $ \retainedPiece{\ccclass} $.
Hence, we can use the record of $ \ccclass $ to represent $ \retainedPiece{\ccclass} $.
This record is added to $ \groupdict(\ccclass) $ with key $ \comp{r} $.

Note that, in both cases, whenever a vertex is moved from one list into another, we also update the size the corresponding records containing such lists.

% The singleton $ \{v_t\} $ initially has no assigned label.
% For each represented \piece{} $ P $, we store its size, a representative $ x_P\in P $, and its inherited label $ \lambda(P) $, where $ \lambda(P)=\bot $ if it has none.
% The inherited labels are temporary attributes of \piece{s}; vertex labels remain unchanged until Step 4.
% Parts other than implicitly represented retained \piece{s} are stored as explicit vertex lists.
% A hash map records the \piece{} containing each old vertex outside a retained part.
% Every other old vertex belongs to the retained \piece{} of its old class, so the \piece{} containing any given vertex can be found without listing its class.

\noindent\setword{\textbf{Step 3 (Merge)}}{StepThreeLabelSync}:
Let $ \piecesToMerge $ be a set containing all \piece{s} constructed for classes with $ d_{\ccclass}>0 $, the singleton $ \{v_t\} $, and the \piece{s} containing each endpoint of an edge in $ \addt{t} $.
This step groups \piece{s} in $ \piecesToMerge $ that belong to the same component of $ \sparseG{t} $ and assigns a label to each group.
Classes (\piece{s}) that are not in $ \piecesToMerge $ remain unchanged, and their labels are not modified.

For each $ \pieceToMerge \in \piecesToMerge $, define its \emph{component key} as $ \comp{x} $, for any $ x\in\pieceToMerge $.
The component key is well-defined because all vertices of a \piece{} share the same component identifier in $ \sparseG{t} $.

After all the \piece{s} in $ \piecesToMerge $ have been assigned a component key, we group the \piece{s} by their component keys.
For each such key $ \kappa $, choose a largest \piece{} $ \pieceToMerge^* $ with key $ \kappa $, breaking ties arbitrarily.
If $ \pieceToMerge^* $ has no assigned label, we assign it a previously unused label.
Vertices from the lists of all other \piece{s} with key $ \kappa $ are inserted into the list of $ \pieceToMerge^{*} $, which represents the new class.
The label of these vertices is updated to the label of $ \pieceToMerge^{*} $.

\textit{Implementation and data structures.}
We store $ \piecesToMerge $ in a set data structure containing records, initially containing a record for the singleton $ \{v_t\} $.
Other \piece{s} are added to $ \piecesToMerge $ as follows.

Recall that after~\ref{StepTwoLabelSync}, for each participating class $ \ccclass $, $ \groupdict(\ccclass) $ is a dictionary containing the records of all \piece{s} split from $ \ccclass $.
We add all such records to $ \piecesToMerge $ by traversing $ \groupdict(\ccclass) $.

For each endpoint $ x $ of an added edge (edges in $ \addt{t} $), based on its record, we can find the record of the \piece{} containing $ x $ and add it to $ \piecesToMerge $.
We also remove that record from $ \setClasses $, if it is currently in $ \setClasses $.

To group the \piece{s} in $ \piecesToMerge $ by their component keys, we use a dictionary mapping each key to a list of records of \piece{s} with that key.
We iterate over $ \piecesToMerge $, compute the component key of each \piece{}, and add its record to the list corresponding to that key in the dictionary.
We refer to this dictionary as $ \compdict $.

For each list of records with the same component key, we scan the list to find a largest \piece{} $ \pieceToMerge^* $ and assign it a new label if it does not have.
If $ \pieceToMerge^* $ is assigned a new label, so are its vertices.
After that, the lists of vertices of all the other \piece{s} in the same list with $ \pieceToMerge^* $ are inserted into the list of vertices of $ \pieceToMerge^* $, and the labels of these vertices are updated to the label of $ \pieceToMerge^* $.

Note that, records of all such vertices are also updated so that their pointers point to the record of the new class and the positions in the list.
The record of $ \pieceToMerge^{*} $ increases its size to reflect the number of vertices its list now contains.
$ \setClasses $ adds to itself the record of $ \pieceToMerge^* $.

To keep track of the vertices whose labels changed, we maintain a list $ \changedList $, which is initially empty.
While assigning labels to \piece{s} in the same component, we add a vertex with it label to $ \changedList $ whenever the \piece{} containing it is assigned a new label or when it is merged into another \piece{}.

\noindent\setword{\textbf{Step 4 (Report and Sync)}}{StepFourLabelSync}:
This step reports the newly arrived vertex $ v_{t} $ and its assigned label, and lists each vertex whose label changed together with its new label.
At the end, we also update the anchor memberships in $ \anchorSet{t} $ and degrees of vertices in $ \sparseG{t} $ for vertices with these updated information.

\textit{Implementation and data structures.}
Reporting the vertices with changed labels is proceeded by one pass over $ \changedList $, which keeps track of all vertices whose labels changed in~\ref{StepThreeLabelSync}.

To update the degrees in $ \sparseG{t} $ and membership in $ \anchorSet{t} $, we traverse the journal $ \journal $ and update the such values for only the vertices appearing in $ \journal $.
If $ v_t $ does not appear in $ \journal $, its degree is set to $ 0 $ and it is not a member of $ \anchorSet{t} $.

\tienlong{Remember to revise the input and output of the algorithm}
\begin{algorithm}[htbp]
  \small
  \DontPrintSemicolon
  \SetKwInOut{KwIn}{Input}
  \SetKwInOut{KwOut}{Output}
  \KwIn{The journal $ \journal $ of update $ t $, the arriving vertex $ v_t $, the stored state at time $ t-1 $, and connectivity queries in $ \sparseG{t} $.}
  \KwOut{$ (v_t,\complabel_t(v_t)) $ and every $ (w,\complabel_t(w)) $ with $ w\in V_{t-1} $ and $ \complabel_t(w)\ne\complabel_{t-1}(w) $.}

    $ \piecesToMerge\gets\{\{v_t\}\} $, $ \complabel(\{ v_{t} \})\gets \perp $

    \tcc{\ref{StepOneLabelSync}}
    \tcp{Make one scan over $ \journal $ to compute the following}
    $ \delt{t} \gets E(\sparseG{t-1})\setminus E(\sparseG{t}) $

    $ \addt{t} \gets E(\sparseG{t}) \setminus E(\sparseG{t-1}) $

    \tcp{Make one scan over $ \delt{t} $ to compute the following}
    $ D_{\ccclass} = \{(x,y)\in\delt{t}:x,y\in\ccclass\} $, $ d_{\ccclass} = \card{D_{\ccclass}} $ for each $ \ccclass \in \ccclasses{t-1} $ such that $ D_{\ccclass} \neq \emptyset $

    $ \delIncEdge{w}:=\card{\{e\in D_{\ccclass}:w\in e\}} $ for each $ w $ such that $ \delIncEdge{w} > 0 $

    \tcc{\ref{StepTwoLabelSync}}
    \ForEach{$ \ccclass $ with $ d_{\ccclass}>0 $}{\nllabel{line:labelsync.split.begin}
        \uIf{$ d_{\ccclass}\ge \card{\ccclass}/2-1 $}{\nllabel{line:labelsync.split.large.deletion.begin}
            \tcp{Make one scan over $ \ccclass $ to decompose it into \piece{s}.}
            $ \mathcal P_{\ccclass} \gets \{ \ccclass_{1}, \ccclass_{2}, \dots \} $ such that $ \ccclass = \cup_{\ccclass_{i}\in \mathcal P_{\ccclass}} \ccclass_{i} $, and $ x,y \in \ccclass_{i} $ if and only if $ \comp{x} = \comp{y} $


            $ \retainedPiece{\ccclass} \gets \argmax_{\ccclass_{i} \in \mathcal P_{\ccclass}}\card{\ccclass_{i}} $ is the retained \piece{} of $ \ccclass $ \tcp*[r]{Break ties arbitrarily.}
        \nllabel{line:labelsync.split.large.deletion.end}
        }
        \Else{\nllabel{line:labelsync.split.few.deletion.begin}
            \tcp{Make one scan over $ D_{\ccclass} $ to compute the following}
            One vertex $ r\in \ccclass\cap\anchorSet{t-1} $ incident to an edge of $ D_{\ccclass} $

            $ W_{\ccclass}\gets\{w\in\bigcup_{e\in D_{\ccclass}}e: w\notin\anchorSet{t-1},\ \delIncEdge{w}=\deg_{t-1}(w)\} $\;

            \tcp{Make one scan over $ W_{\ccclass} $ to decompose $ \ccclass $ into \piece{s}.}
            $ \mathcal P_{\ccclass} \gets \{ \ccclass_{1}, \ccclass_{2}, \dots \} $ such that $ \ccclass = \cup_{\ccclass_{i}\in \mathcal P_{\ccclass}} \ccclass_{i} $, and $ x,y \in \ccclass_{i} $ if and only if $ \comp{x} = \comp{y} $

            $ \retainedPiece{\ccclass} \gets \ccclass_{i} := \{ x: \comp{x} = \comp{r} \} $ is the retained \piece{} of $ \ccclass $
        }
        $ \complabel(\retainedPiece{\ccclass}) \gets \complabel(\ccclass) $

        $ \complabel(\ccclass_{i}) \gets \perp $ for all $ \ccclass_{i} \in \mathcal P_{ \ccclass } \setminus \{ \retainedPiece{\ccclass} \} $\;
        $ \piecesToMerge \gets \piecesToMerge \cup \{ \ccclass_{i}|\  \ccclass_{i}\in \mathcal P_{ \ccclass } \} $
    }\nllabel{line:labelsync.split.end}

    \tcc{\ref{StepThreeLabelSync}}
    \tcp{Make one scan over $ \addt{t} $ to add \piece{s} containing endpoints of edges in $ \addt{t} $ to $ \piecesToMerge $.}
    $ \piecesToMerge\gets\piecesToMerge\cup\{P: P \text{ is a \piece{} containing a vertex } x\in e\in \addt{t} \text{ and } P \text{ is not already in }\piecesToMerge\} $ \nllabel{line:labelsync.merge.begin} \nllabel{line:labelsync.piecesToMerge.include.addt}

    \BlankLine
    \tcp{Compute the component key of each \piece{} in $ \piecesToMerge $.}
    \ForEach{$ \pieceToMerge \in\piecesToMerge $}{ 
        $ \mathrm{key}(\pieceToMerge )\gets\comp{x} $, for any $ x\in\pieceToMerge $ \tcp*[r]{note that $ \comp{x} = \comp{y} $ for any $ x,y\in\pieceToMerge $} \nllabel{line:labelsync.assign.key}
    }

    $ \changedList \gets \{  \} $ \tcp*[r]{List to keep track of vertices whose labels changed.}
    \BlankLine
    \tcp{Group \piece{s} in $ \piecesToMerge $ by their component keys and assign labels.}
    \ForEach{distinct component key $ \kappa $ represented in $ \piecesToMerge $}{
        $ \piecesToMerge_\kappa\gets\{\pieceToMerge \in\piecesToMerge:\mathrm{key}(\pieceToMerge )=\kappa\} $\;\nllabel{line:labelsync.group.key}

        $ \pieceToMerge^* \gets \argmax_{\pieceToMerge \in\piecesToMerge_\kappa} \card{\pieceToMerge} $ \tcp*[r]{Break ties arbitrarily.}

        \uIf{$ \complabel(\pieceToMerge^*) = \perp $}{
            $ \complabel(\pieceToMerge^*) \gets \text{a previously unused label} $ \tcp*[r]{Assign a new label to $ \pieceToMerge^* $}

            Add $ (v, \complabel(\pieceToMerge^*)) $ to $ \changedList $, and update the label of $ v $ to $ \complabel(\pieceToMerge^*) $, for all $ v\in\pieceToMerge^* $
        }

        Add $ (v, \complabel(\pieceToMerge^*)) $ to $ \changedList $ for all $ v\in\pieceToMerge $, for all $ \pieceToMerge \in\piecesToMerge_\kappa\setminus\{\pieceToMerge^*\} $ \nllabel{line:labelsync.merge.add.small.pieces.changedlist}

        Merge $ \pieceToMerge $ to $ \pieceToMerge^* $ for all $ \pieceToMerge \in\piecesToMerge_\kappa\setminus\{\pieceToMerge^*\} $ and update the labels of vertices in $ \pieceToMerge $ to $ \complabel(\pieceToMerge^*) $


    }\nllabel{line:labelsync.merge.end}

    \tcc{\ref{StepFourLabelSync}}
    Make one scan over $ \journal{t} $ to update the degrees of vertices in $ \sparseG{t} $ and their membership in $ \anchorSet{t} $

    \Return{$ \changedList $}
  \caption{$ \LabelSync $ at the end of update $ t $.}
  \label{alg.label.sync}
\end{algorithm}
% paragraph Detailed Description of LabelSync (end)


% subsection Description of LabelSync (end)

\subsection{Analysis for \texorpdfstring{$ \LabelSync $}{LabelSync}}
\label{sub:analysis_for_labelsync}

\paragraph{Correctness.}
\label{par:correctness}
We use the invariant that every edge of $ \sparseG{t} $ has an endpoint in $ \anchorSet{t} $.
Let $ \successDA $ denote the success event from the analysis of $ \explicitDA $; on this event, the density bound in \cref{lem:ccclass.is.dense} holds at every time step.

\begin{lemma} \label{lem:labelsync.robust.components}
Let $ \sparseG $ be a simple undirected graph and let $ \anchorSet $ be a vertex set containing at least one endpoint of every edge of $ \sparseG $.
Let $ \ccclass $ be a connected component with $ \card{\ccclass}>1 $, and suppose that
\[
  |\openN{\sparseG}{a}|\ge\tfrac34\card{\ccclass}-1
  \qquad\text{for every }a\in \ccclass\cap\anchorSet.
\]
For any $ D\subseteq E(\sparseG[\ccclass]) $ with $ |D|<\card{\ccclass}/2-1 $, the graph $ \sparseG[\ccclass]\setminus D $ has one component containing $ \ccclass\cap\anchorSet $.
Every other component is a singleton outside $ \anchorSet $, and there are at most $ |D| $ such singletons.
\end{lemma}
\begin{proof}
    As $ \card{\ccclass} > 1 $, $ \ccclass $ has an least one edge.
    By definition, $ \anchorSet $ contains at least one vertex of all edges.
    Hence, $ \ccclass \cap \anchorSet \neq \emptyset $.
    We proceed by cases.

    $ \bullet $ If $ \card{\ccclass \cap \anchorSet} = 1 $, suppose $ \ccclass \cap \anchorSet = \{ a \} $.
    By definition of $ \anchorSet $, this implies that all edges of $ \ccclass $ are incident on $ a $, namely, $ \ccclass $ is a star graph with $ a $ as the center.
    Hence, deleting $ D $ results in a graph with one component containing $ a $ and all other components are singletons, and there are at most $ \card{D} $ of them.

    $ \bullet $ If $ \card{\ccclass \cap \anchorSet} > 1 $, consider two vertices $ a \neq b \in \ccclass \cap \anchorSet $. We show that the number of edge-disjoint paths of length at most $ 2 $ between $ a $ and $ b $ in $ \sparseG[\ccclass] $ is more than $ \card{D} $.

    Note that, for any two vertices $ w, w' \in \openN{\sparseG}{a} \cap \openN{\sparseG}{b} $, the paths $ (a,w,b) $ and $ (a,w',b) $ are of length $ 2 $ and edge-disjoint.
    If $ (a,b) $ is not an edge in $ \sparseG $, the number of such that is $ \card{\openN{\sparseG}{a} \cap \openN{\sparseG}{b}} \geq 2(\frac{3}{4}\card{\ccclass}-1) - (\card{\ccclass} - 2) \geq \frac{1}{2}\card{C} > \card{D} $.
    If $ (a,b) $ is an edge in $ \sparseG $, then $ \card{\openN{\sparseG}{a} \cap \openN{\sparseG}{b}} + \geq 2(\frac{3}{4}\card{\ccclass}-1) - \card{\ccclass} \geq \frac{1}{2}\card{C} - 2 $.
    The number of edge-disjoint paths of length at most $ 2 $ between $ a $ and $ b $, which also includes $ (a,b) $ in this case, is $ \card{\openN{\sparseG}{a} \cap \openN{\sparseG}{b}} + 1 \geq \frac{1}{2}\card{\ccclass} - 1 > \card{D} $.
    Hence, after deleting the edges in $ D $, there is still a path between $ a $ and $ b $.
    As this holds for any two vertices in $ \ccclass \cap \anchorSet $, all vertices in $ \ccclass \cap \anchorSet $ are in the same connected component of $ \sparseG[\ccclass]\setminus D $.

    After deleting $ D $, for any vertex in $ \ccclass \setminus \anchorSet $ that is not in the same component as $ \ccclass \cap \anchorSet $, all of its incident edges must be in $ D $, as the other endpoint of any such edge is in $ \ccclass \cap \anchorSet $.
    This also implies that all other components are singletons, and there are at most $ \card{D} $ of them.

\end{proof}
\begin{lemma}
\label{lem:labelsync.classes.are.connected.components}
Conditioned on $ \successDA $, after every update $ t $, the classes maintained by $ \LabelSync $ are exactly the connected components of $ \sparseG{t} $.
\end{lemma}
\begin{proof}
    We prove by induction on time step $ t $.

    At the base case $ t=1 $, the lemma holds as both $ \sparseG{1} $ and the classes maintained by $ \LabelSync{1} $ are both $ \{v_{1}\} $.

    For the inductive step, assume that the lemma holds for time step $ t-1 $, that is, every class maintained by $ \LabelSync{t-1} $ is a connected component of $ \sparseG{t-1} $.

    We show that it also holds for time step $ t $.

    $ \bullet $ First, we demonstrate that, after the split step (\cref{line:labelsync.split.begin}-\cref{line:labelsync.split.end}), every vertex in same \piece{} is in the same connected component of $ \sparseG{t} $. 

    Indeed, for each class $ \ccclass $ not participating in the split step ($ d_{\ccclass} = 0 $), they have no deleted edges.
    By induction hypothesis, $ \ccclass $ is a connected component of $ \sparseG{t-1} $. With none of its edges deleted, $ \ccclass $ remains a connected component of $ \sparseG{t} $.
    Hence, all its vertices are in the same connected component of $ \sparseG{t} $.

    It remains to consider classes $ \ccclass $ participating in the split step ($ d_{\ccclass} \geq 1 $).
    If $ d_{\ccclass} \geq \frac{1}{2}\card{\ccclass}-1 $, the algorithm iterates over all the vertices in $ \ccclass $, grouping vertices in the same connected component of $ \sparseG{t} $ together (\cref{line:labelsync.split.large.deletion.begin}-\cref{line:labelsync.split.large.deletion.end}).
    Note that, determining whether two vertices $ x,y $ are in the same connected component of $ \sparseG{t} $ can be done by checking if $ \comp{x} = \comp{y} $. 
    If $ d_{\ccclass} < \frac{1}{2}\card{\ccclass} - 1 $, conditioned on $ \successDA $, by~\cref{lem:ccclass.is.dense}, we have $ \openN{\sparseG{t-1}}{a} \geq \frac{3}{4}\card{\ccclass}-1 $ for any $ a \in \ccclass \cap \anchorSet{t-1} $.
    Now, by \cref{lem:labelsync.robust.components}, all vertices in $ \ccclass \cap \anchorSet{t-1} $ are in the same connected component of $ \sparseG{t} $, and all other vertices in $ \ccclass $ that are not in the same component as $ \ccclass \cap \anchorSet{t-1} $ are singletons.

    $ \bullet $ To finish the proof, it suffices to show that, the merge step (\cref{line:labelsync.merge.begin}-\cref{line:labelsync.merge.end}) merges all \piece{s} that are in the same connected component of $ \sparseG{t} $ and assigns them the same label, and assigns different labels to \piece{s} in different components. 
    This step begins with a set $ \piecesToMerge $, consisting of all the \piece{s} created in the split step, and the singleton $ \{v_t\} $; then adding to $ \piecesToMerge $ all the \piece{s} containing endpoints of an edge in $ \addt{t} $ (\cref{line:labelsync.merge.begin}-\cref{line:labelsync.piecesToMerge.include.addt}).
    All the \piece{s} participating in the merge step are the \piece{s} in $ \piecesToMerge $.

    If a \piece{} does not participate in the merge step, that is, is not included in the set $ \piecesToMerge $, then all of its vertices are not incident to any edge in $ \addt{t} $.
    By induction hypothesis, this \piece{} is a connected component of $ \sparseG{t-1} $.
    As no edges are added to this \piece{}, it remains a connected component of $ \sparseG{t} $ and keeps its label.

    For each \piece{} $ \pieceToMerge \in \piecesToMerge $, \cref{line:labelsync.assign.key} assigns a key $ \mathrm{key}(\pieceToMerge) = \comp{x} $, for some $ x \in \pieceToMerge $.
    We have shown above that all vertices in $ \pieceToMerge $ are in the same connected component of $ \sparseG{t} $, and hence, they have the same component identifier.
    It follows that, two \piece{s} share the same key if and only if they are in the same connected component of $ \sparseG{t} $.
    By grouping \piece{s} by their keys, the algorithm merges all \piece{s} which forms exactly one connected component of $ \sparseG{t} $ (\cref{line:labelsync.group.key}-\cref{line:labelsync.merge.end}).

    Finally, we argue that different components receive different labels. Indeed, each old label is inherited by exactly one retained \piece{}, and fresh labels are distinct and have never been used before.
    The classes after update $ t $ are therefore exactly the components of $ \sparseG{t} $.
\end{proof}

\begin{lemma}\label{lem:number.of.updated.vertices}
    For every $ t\leq n $, the total number of label changes of arrived vertices after executing $ \LabelSync $ (\cref{alg.label.sync}) for the first $ t $ updates is at most $ t\log_{2}n + 4(\log_{2}n + 3)\sum_{t'\leq t} \card{\delt{t'}} $.
\end{lemma}
\newcommand{\newpiecesize}[1]{\sigma_{#1}}
\newcommand{\smallpiecesize}[1]{\pi_{#1}}
\newcommand{\potential}[1]{\mathrm{Budget}(#1)}
\begin{proof}
    Fix any update time $ t'\in [1,t] $.
    We count the number of label changes after executing $ \LabelSync $ for update $ t' $, that is, the number of vertices with different labels between time $ t'-1 $ and $ t' $.

    At a high level, we upper bound this quantity by the number of vertices participating in the split and merge steps of $ \LabelSync $ for update $ t' $.

    Formally, consider the execution of $ \LabelSync $ for update $ t' $.
    Let $ \newpiecesize{t'} $ be the total size of the new \piece{s} created in the split step (\cref{line:labelsync.split.begin}-\cref{line:labelsync.split.end}). 
    In the merge step (\cref{line:labelsync.merge.begin}-\cref{line:labelsync.merge.end}), denote by $ \smallpiecesize{t'} $ be the total size of all the \piece{s} in $ \piecesToMerge $ that are not $ \pieceToMerge^* $.

    We first claim that, the number of vertices with changed labels is at most $ \newpiecesize{t'} + \smallpiecesize{t'} $. Indeed, the vertices with changed labels must belong to:
    \begin{itemize}
        \item Either one of the \piece{s} created in the split step, which also includes \piece{s} $ \piecesToMerge^{*} $ that requires a fresh label in the merge steps. The total size of these \piece{s} is at most $ \newpiecesize{t'} $.
        \item Or one of the \piece{s} in $ \piecesToMerge $ that are not $ \pieceToMerge^* $, which has total size at most $ \smallpiecesize{t'} $.
    \end{itemize}

    It suffices to demonstrate that $ \sum_{t'\leq t} \newpiecesize{t'} + \sum_{t'\leq t}\smallpiecesize{t'} \leq t \log_{2}n + 4(\log_{2} n + 3) \sum_{t'\leq t}\card{\delt{t'}} $.

    \paragraph{Bounding $ \newpiecesize{t'} $.}\label{par:bounding.newpiecesize} % (fold)
    In the split step, if $ d_{\ccclass} \geq \card{\ccclass}/2-1 $, the algorithm traverses all vertices in $ \ccclass $ and groups them into \piece{s} according to their identifiers. 
    The total size of the new \piece{s} created in this case is therefore at most $ \card{\ccclass} \leq 2d_{\ccclass}+2 \leq 4d_{\ccclass} $.
    In the remaining case, if $ d_{\ccclass} < \card{\ccclass}/2-1 $, the new \piece{s} are created grouping vertices in $ W_{\ccclass} $ according to their identifiers.
    The total size of the new \piece{s} created in this case is at most $ \card{W_{\ccclass}} $.
    By definition of $ W_{\ccclass} $ and the invariant~\eqref{eq:invariant.anchor}, we have $ \card{W_{\ccclass}} \leq d_{\ccclass} $.
    In either cases, the total size of the new \piece{s} created for class $ \ccclass $ is at most $ 4d_{\ccclass} $.
    Thus, the total size of the new \piece{s} created in the split step is
    \begin{align}
        \newpiecesize{t'} \leq \sum_{\ccclass \text{ participated in the split step}} 4d_{\ccclass} = 4\card{\delt{t'}}. \label{eq:bound.newpiecesize}
    \end{align}
    % paragraph Bounding newpiecesize (end)

    \paragraph{Bounding $ \smallpiecesize{t'}.$}\label{par:bounding_smallpiecesize_t_} % (fold)
    We show an upper bound for $ \sum_{t'\leq t} \smallpiecesize{t'} $ instead of bounding $ \smallpiecesize{t'} $ individually for each $ t' $.

    Consider the following scenario: throughout the whole course of the algorithm, for all time step $ t' $, there are no classes participating in the split step.
    Observe the merge step, a vertex $ v $ changes its label only if the \piece{} containing $ v $ joins a \piece{} with a larger size.
    This operation at least doubles the size of the \piece{} containing $ v $.
    Since the split opertation is never performed, it follows that any vertex can change its label at most $ \log_{2}n $ times.

    To account for the split step, we employ a potential function argument, which enables us to track the number of label changes allowed for each vertex
    Intuitively, every newly arrived vertex $ v $ is assigned a potential $ \potential{\{v\}} $ of $ \log_{2}n $, which corresponds to the maximum number of label changes it can undergo if the class containing $ v $ never participates in the split step.
    A split step performed on a class reduces the potential of each vertex in that class, while a merge step increases the potential of each vertex in the classes that get merged into a larger class.

    Formally, we define the potential function as follows. For a partition $ \ccclasses $ of a set $ V $, the potential of a class $ \ccclass\in\ccclasses $ is defined as
    \begin{align}
        \potential{\ccclass} = \card{\ccclass}\log_{2} \frac{n}{\card{\ccclass}}, \nonumber
    \end{align}
    and the potential of the whole partition is
    \begin{align}
        \potential{\ccclasses} = \sum_{\ccclass\in\ccclasses} \potential{\ccclass} = \sum_{\ccclass\in\ccclasses} \card{\ccclass}\log_{2} \frac{n}{\card{\ccclass}}. \nonumber
    \end{align}
    For each newly arrived vertex $ v $, assign a potential $ \potential{\{v\}} = \log_{2}n $.

    We analyze the change in potential of the partition maintained by $ \LabelSync $ at following three points in time: $ t'-1 $, after the split step of update $ t' $, and after the merge step of update $ t' $.
    At time $ t'-1 $, the partition is $ \potential{\ccclasses{t'-1}} $, which is the partition of classes maintained by $ \LabelSync $ at time $ t'-1 $, serving as the input to $ \LabelSync{t'} $. The partition after the merge step of update $ t' $ is the output of $ \LabelSync{t'} $. The partition after the split step of update $ t' $, denoted by $ \widetilde{\ccclasses{t'}} $, consists of $ v_{t'} $, the new \piece{s} created in the split step, and the \piece{s} that are not split.

    The split step transforming $ \ccclasses{t'-1} $ into $ \widetilde{\ccclasses{t'}} $ increases the potential. In particular, for each class $ \ccclass \in \ccclass{t'-1} $ split into $ \pieceToMerge_{1}, \pieceToMerge_{2}, \dots  $, with $ \pieceToMerge_{1} $ being the retained \piece{}, the potential increases by
    \begin{align}
        \sum_{i} \card{\pieceToMerge_{i}}\log_{2} \frac{n}{\card{\pieceToMerge_{i}}} - \card{\ccclass}\log_{2} \frac{n}{\card{\ccclass}} &= \sum_{i}\card{\pieceToMerge_{i}} \log_{2}\dfrac{\card{\ccclass}}{\card{\pieceToMerge_{i}}} \nonumber \\
                                                                                                                                                 &\leq \card{\pieceToMerge_{1}}\log_{2}\left(1+\dfrac{\card{\ccclass}-\card{\pieceToMerge_{1}}}{\card{\pieceToMerge_{1}}}\right) + (\card{\ccclass} - \card{\pieceToMerge_{1}})\log_{2}n  \nonumber \\
                                                                                                                                                 &\leq (\card{\ccclass} - \card{\pieceToMerge_{1}})(\log_{2}e + \log_{2}n). \nonumber
    \end{align}
    Note that, newly arrived vertex $ v_{t'} $ contributes $ \log_{2}n $ to the potential, and $ \sum_{\ccclass\in \ccclass{t'-1}}\card{\ccclass} - \card{\pieceToMerge_{1}} \leq \newpiecesize{t'} $.
    Hence,
    \begin{align}
        \potential{\widetilde{\ccclasses{t'}}} - \potential{\ccclasses{t'-1}} \leq \newpiecesize{t'}(\log_{2}e + \log_{2}n) + \log_{2}n. \label{eq.split.increases.potential}
    \end{align}

    The merge step transforming $ \widetilde{\ccclasses{t'}} $ into $ \ccclasses{t'} $ decreases the potential, and we argue that the decrease is at least $ \smallpiecesize{t'} $. 
    Indeed, consider a merge of $ \pieceToMerge_{1}, \pieceToMerge_{2}, \dots, $ with $ \pieceToMerge^* $ into $ \ccclass' $. Here, $ \card{\pieceToMerge^*} \geq \pieceToMerge_{i} $, and hence $ \card{\ccclass}\geq 2 \card{\pieceToMerge_{i}} $.
    Thus, for each such merge, the potential decreases by at least
    \begin{align}
        \sum_{i} \card{\pieceToMerge_{i}}\log_{2} \frac{n}{\card{\pieceToMerge_{i}}} - \card{\ccclass'}\log_{2} \frac{n}{\card{\ccclass'}} &\geq \sum_{i}\card{\pieceToMerge_{i}} \log_{2}\dfrac{\card{\ccclass'}}{\card{\pieceToMerge_{i}}} \geq \sum_{i} \card{\pieceToMerge_{i}}. \nonumber
    \end{align}
    By definition, $ \smallpiecesize{t'} $ covers the total size of all such $ \pieceToMerge_{i} $.
    It follows that $ \potential{\widetilde{\ccclasses{t'}}} - \potential{\ccclasses{t'}} \geq \smallpiecesize{t'} $.

    Taking the sum over all $ t'\leq t $, and note that $ \potential{\ccclasses{0}} = 0 \leq \potential{\ccclasses{t}} $, we have
    { \small
    \begin{align}
        \sum_{t'\leq t} \smallpiecesize{t'} &\leq (\potential{\widetilde{\ccclasses{1}}} - \potential{\ccclasses{1}}) + (\potential{\widetilde{\ccclasses{2}}} - \potential{\ccclasses{2}}) + \dots + (\potential{\widetilde{\ccclasses{t}}} - \potential{\ccclasses{t}}) \nonumber \\
                                            &\leq (\potential{\widetilde{\ccclasses{1}}} - \potential{\ccclasses{0}}) + (\potential{\widetilde{\ccclasses{2}}} - \potential{\ccclasses{1}}) + \dots + (\potential{\widetilde{\ccclasses{t}}} - \potential{\ccclasses{t-1}}) \nonumber \\
                                            &\leq \sum_{t'\leq t} \newpiecesize{t'}(\log_{2}e + \log_{2}n) + t\log n, \label{eq:bound.smallpiecesize}
    \end{align} 
    }
    where the last inequality follows from~\eqref{eq.split.increases.potential}.

    To finish the proof, combining~\eqref{eq:bound.newpiecesize} and~\eqref{eq:bound.smallpiecesize}, we obtain
    \begin{align}
        \sum_{t'\leq t} \newpiecesize{t'} + \sum_{t'\leq t}\smallpiecesize{t'} &\leq \sum_{t'\leq t} 4\card{\delt{t'}} + \sum_{t'\leq t} \newpiecesize{t'}(\log_{2}e + \log_{2}n) + t\log n \nonumber \\
                                                                                  &\leq 4(\log_{2}n + 3)\sum_{t'\leq t}\card{\delt{t'}} + t\log n. \nonumber
    \end{align}
    % paragraph Bounding $ \smallpiecesize{t'} $ (end)
\end{proof}

\begin{lemma}\label{lem:labelsync.cost}
    Suppose that after the update time $ t $, the length of the journal is $ \jnTime{t} $ and the number of vertices whose cluster label changes is $ \nChangedVers{t} $. 
    \tienlong{May be later remove this number of vertices whose cluster label changes, it seems confusing when bounding the running tim of the algorithm, or may be not, for example, suffix tree finding k matches}
    Then the running time of $ \LabelSync $ is $ O(\jnTime{t}\log n + \nChangedVers{t}) $.
\end{lemma}

\begin{proof}
    Write $ D= \card{\delt{t}} $ and $ A= \card{\addt{t}} $.
    Since the journal records all the edges that are removed or added, we have $ D+A\leq \jnTime{t} $.

    Scanning the journal.
    Identifying the deleted and added edges takes $ O(\jnTime{t}) $ dictionary operations.
    Scanning $\delt{t}$ then computes the counters and groups deleted edges by old class in $O(D)$ operations.
    Neither operation scans all vertices or all classes.

    Processing splits. There are at most $D$ processed classes, since each has $d_C\ge1$, and $\sum_C d_C\le D.$
    For a class in the many-deletions case, $|C|\le4d_C$, so enumerating and grouping its vertices takes $O(d_C\log n)$ time.
    In the few-deletions case, scanning its deleted edges identifies $W_C$ and $r$ in $O(d_C)$ operations, and querying and grouping the stranded vertices takes $O(d_C\log n)$ time.
    The retained piece remains in the old class list without enumerating it in the few-deletions case.
    Moving vertices into new piece lists takes time proportional to their number.

    Consequently, the entire split step takes $O(D\log n)$ time.
    Moreover, if $\sigma_t$ is the total number of vertices in new pieces, then $\sigma_t\le4D.$

    Bounding and constructing $\piecesToMerge$.
    There are at most $\sigma_t\le4D$ new pieces, since each is nonempty, and at most $D$ retained pieces, one per processed class.
    Added-edge endpoints introduce at most $2A$ further pieces, and $v_t$ contributes one singleton.
    Hence $|\piecesToMerge|\le4D+D+2A+1=O(\jnTime{t}+1)$.
    Using the piece lookup and deduplication described above, constructing this collection takes $O(\jnTime{t}+1)$ dictionary operations.

    Grouping and choosing labels.
    Querying one representative per piece and grouping by component identifier takes $O(|\piecesToMerge|\log n)$ time.
    Finding a largest piece in each group and deciding its final label takes $O(|\piecesToMerge|)$ further operations.

    Enumerating relabelled vertices.
    Every piece is enumerated directly from its linked list only when its vertices receive a different label.
    This includes a winning new piece that receives a fresh label; the newly inserted vertex contributes one additional operation.
    Thus all enumeration takes $O(\nChangedVers{t}+1)$ time.

    Reporting and bookkeeping.
    Applying the collected moves and reporting labels takes $O(\nChangedVers{t}+1)$ operations.
    Updating flags and degrees for journal vertices takes $O(\jnTime{t}+1)$ operations.

    Combining these bounds gives $O((\jnTime{t}+1)\log n+\nChangedVers{t})$.
    Separately, there are at most $\jnTime{t}$ journalled edge operations, each requiring $O(\log^2 n)$ amortized connectivity-update time.
\end{proof}

% section Explicit Dynamic Correlation Clustering Algorithm for Insertion-Only Stream (end)




\end{document}
